% -----------------------------------------------------------
% -----------------------------------------------------------
\part{The Life of Christ and Synopsis of the Gospels}
% -----------------------------------------------------------
% -----------------------------------------------------------
	\label{LCHRIST}
	
	This list of events comes from a synopsis of the gospels in the book \textit{Charting the New Testament}, by John Welch and John Hall, 2002.
	
	\begin{multicols}{3}
	
	\begin{enumerate}
	
		\footnotesize
	
			% 1
		\item The Word Becomes Flesh
			\begin{compactitem}
				\item[\textbf{Matthew}] ---
				\item[\textbf{Mark}] ---
				\item[\textbf{Luke}] ---
				\item[\textbf{John}] \hyperref[JOHN-1-1]{1:1--14}
			\end{compactitem}
			
			% 2
		\item Luke's Sources and Purpose
			\begin{compactitem}
				\item[\textbf{Matthew}] ---
				\item[\textbf{Mark}] ---
				\item[\textbf{Luke}] \hyperref[LUKE-1-1]{1:1--4}
				\item[\textbf{John}] ---
			\end{compactitem}
			
			% 3
		\item Birth of John the Baptist Foretold
			\begin{compactitem}
				\item[\textbf{Matthew}] ---
				\item[\textbf{Mark}] ---
				\item[\textbf{Luke}] \hyperref[LUKE-1-5]{1:5--25}
				\item[\textbf{John}] ---
			\end{compactitem}
			
			% 4
		\item Enunciation to Mary
			\begin{compactitem}
				\item[\textbf{Matthew}] ---
				\item[\textbf{Mark}] ---
				\item[\textbf{Luke}] \hyperref[LUKE-1-26]{1:26--38}
				\item[\textbf{John}] ---
			\end{compactitem}
			
			% 5
		\item Mary Visits Elizabeth
			\begin{compactitem}
				\item[\textbf{Matthew}] ---
				\item[\textbf{Mark}] ---
				\item[\textbf{Luke}] \hyperref[LUKE-1-39]{1:39--45}
				\item[\textbf{John}] ---
			\end{compactitem}
			
			% 6
		\item Mary's Song of Praise
			\begin{compactitem}
				\item[\textbf{Matthew}] ---
				\item[\textbf{Mark}] ---
				\item[\textbf{Luke}] \hyperref[LUKE-1-46]{1:46--56}
				\item[\textbf{John}] ---
			\end{compactitem}
			
			% 7
		\item Birth of John the Baptist
			\begin{compactitem}
				\item[\textbf{Matthew}] ---
				\item[\textbf{Mark}] ---
				\item[\textbf{Luke}] \hyperref[LUKE-1-57]{1:57--63}
				\item[\textbf{John}] ---
			\end{compactitem}
			
			% 8
		\item Prophecy of Zacharias
			\begin{compactitem}
				\item[\textbf{Matthew}] ---
				\item[\textbf{Mark}] ---
				\item[\textbf{Luke}] \hyperref[LUKE-1-64]{1:64--80}
				\item[\textbf{John}] ---
			\end{compactitem}
			
			% 9
		\item Genealogy of Jesus
			\begin{compactitem}
				\item[\textbf{Matthew}] \hyperref[MATTHEW-1-1]{1:1--17}
				\item[\textbf{Mark}] ---
				\item[\textbf{Luke}] \hyperref[LUKE-3-23]{3:23--38}
				\item[\textbf{John}] ---
			\end{compactitem}
			
			% 10
		\item Birth of Jesus
			\begin{compactitem}
				\item[\textbf{Matthew}] \hyperref[MATTHEW-1-18]{1:18--25}
				\item[\textbf{Mark}] ---
				\item[\textbf{Luke}] \hyperref[LUKE-2-1]{2:1--7}
				\item[\textbf{John}] ---
			\end{compactitem}
			
	\columnbreak
			
			% 11
		\item Angels Visit the Shepherds
			\begin{compactitem}
				\item[\textbf{Matthew}] ---
				\item[\textbf{Mark}] ---
				\item[\textbf{Luke}] \hyperref[LUKE-2-8]{2:8--14}
				\item[\textbf{John}] ---
			\end{compactitem}
			
			% 12
		\item Shepherds Visit Jesus
			\begin{compactitem}
				\item[\textbf{Matthew}] ---
				\item[\textbf{Mark}] ---
				\item[\textbf{Luke}] \hyperref[LUKE-2-15]{2:15--20}
				\item[\textbf{John}] ---
			\end{compactitem}
			
			% 13
		\item Presentation of Jesus in the Temple
			\begin{compactitem}
				\item[\textbf{Matthew}] ---
				\item[\textbf{Mark}] ---
				\item[\textbf{Luke}] \hyperref[LUKE-2-21]{2:21--39}
				\item[\textbf{John}] ---
			\end{compactitem}
			
			% 14
		\item Wise Men Visit Jesus
			\begin{compactitem}
				\item[\textbf{Matthew}] \hyperref[MATTHEW-2-1]{2:1--12}
				\item[\textbf{Mark}] ---
				\item[\textbf{Luke}] ---
				\item[\textbf{John}] ---
			\end{compactitem}
			
			% 15
		\item Flight to Egypt
			\begin{compactitem}
				\item[\textbf{Matthew}] \hyperref[MATTHEW-2-13]{2:13--15}
				\item[\textbf{Mark}] ---
				\item[\textbf{Luke}] ---
				\item[\textbf{John}] ---
			\end{compactitem}
			
			% 16
		\item Slaying of the Infants by Herod
			\begin{compactitem}
				\item[\textbf{Matthew}] \hyperref[MATTHEW-2-16]{2:16--18}
				\item[\textbf{Mark}] ---
				\item[\textbf{Luke}] ---
				\item[\textbf{John}] ---
			\end{compactitem}
			
			% 17
		\item Return to Nazareth
			\begin{compactitem}
				\item[\textbf{Matthew}] \hyperref[MATTHEW-2-19]{2:19--23}
				\item[\textbf{Mark}] ---
				\item[\textbf{Luke}] \hyperref[LUKE-2-39]{2:39}
				\item[\textbf{John}] ---
			\end{compactitem}
			
			% 18
		\item Jesus' Childhood
			\begin{compactitem}
				\item[\textbf{Matthew}] ---
				\item[\textbf{Mark}] ---
				\item[\textbf{Luke}] \hyperref[LUKE-2-40]{2:40}, \hyperref[LUKE-2-52]{52}
				\item[\textbf{John}] ---
			\end{compactitem}
	
			% 19
		\item Jesus as a Boy in the Temple
			\begin{compactitem}
				\item[\textbf{Matthew}] ---
				\item[\textbf{Mark}] ---
				\item[\textbf{Luke}] \hyperref[LUKE-2-41]{2:41--52}
				\item[\textbf{John}] ---
			\end{compactitem}
			
			% 20
		\item \textit{Preaching of John the Baptist}
			\begin{compactitem}
				\item[\textbf{Matthew}] \hyperref[MATTHEW-3-1]{3:1--12}
				\item[\textbf{Mark}] \hyperref[MARK-1-1]{1:1--8}
				\item[\textbf{Luke}] \hyperref[LUKE-3-2]{3:2--18}
				\item[\textbf{John}] \hyperref[JOHN-1-6]{1:6--8}, \hyperref[JOHN-1-15]{1:15--28}
			\end{compactitem}
			
	\columnbreak

			% 21
		\item John Teaches Generosity and Moderation
			\begin{compactitem}
				\item[\textbf{Matthew}] ---
				\item[\textbf{Mark}] ---
				\item[\textbf{Luke}] \hyperref[LUKE-3-10]{3:10--14}
				\item[\textbf{John}] ---
			\end{compactitem}
			
			% 22
		\item John's Messianic Preaching
			\begin{compactitem}
				\item[\textbf{Matthew}] \hyperref[MATTHEW-3-11]{3:11--12}
				\item[\textbf{Mark}] ---
				\item[\textbf{Luke}] \hyperref[LUKE-3-15]{3:15--18}
				\item[\textbf{John}] \hyperref[JOHN-1-19]{1:19--28}
			\end{compactitem}
			
			% 23
		\item \textit{John's Imprisonment}
			\begin{compactitem}
				\item[\textbf{Matthew}] \hyperref[MATTHEW-14-3]{14:3--5}
				\item[\textbf{Mark}] \hyperref[MARK-6-17]{6:17--19}
				\item[\textbf{Luke}] \hyperref[LUKE-3-19]{3:19--20}
				\item[\textbf{John}] ---
			\end{compactitem}
			
			% 24
		\item John Recognizes Jesus as the Lamb of God
			\begin{compactitem}
				\item[\textbf{Matthew}] ---
				\item[\textbf{Mark}] ---
				\item[\textbf{Luke}] ---
				\item[\textbf{John}] \hyperref[JOHN-1-29]{1:29--34}
			\end{compactitem}
			
			% 25
		\item \textit{Baptism of Jesus}
			\begin{compactitem}
				\item[\textbf{Matthew}] \hyperref[MATTHEW-3-13]{3:13--17}
				\item[\textbf{Mark}] \hyperref[MARK-1-9]{1:9--11}
				\item[\textbf{Luke}] \hyperref[LUKE-3-21]{3:21--22}
				\item[\textbf{John}] ---
			\end{compactitem}
			
			% 26
		\item \textit{Temptation of Christ}
			\begin{compactitem}
				\item[\textbf{Matthew}] \hyperref[MATTHEW-4-1]{4:1--11}
				\item[\textbf{Mark}] \hyperref[MARK-1-12]{1:12--13}
				\item[\textbf{Luke}] \hyperref[LUKE-4-1]{4:1--13}
				\item[\textbf{John}] ---
			\end{compactitem}
			
			% 27
		\item \textit{First Preaching in Galilee}
			\begin{compactitem}
				\item[\textbf{Matthew}] \hyperref[MATTHEW-4-12]{4:12--17}
				\item[\textbf{Mark}] \hyperref[MARK-1-14]{1:14--15}
				\item[\textbf{Luke}] \hyperref[LUKE-4-14]{4:14--15}
				\item[\textbf{John}] ---
			\end{compactitem}
	
			% 28
		\item \textit{Jesus Calls His First Disciples}
			\begin{compactitem}
				\item[\textbf{Matthew}] \hyperref[MATTHEW-4-18]{4:18--22}
				\item[\textbf{Mark}] \hyperref[MARK-1-16]{1:16--20}
				\item[\textbf{Luke}] \hyperref[LUKE-5-1]{5:1--11}
				\item[\textbf{John}] \hyperref[JOHN-1-35]{1:35--51}
			\end{compactitem}
			
			% 29
		\item Wedding of Cana
			\begin{compactitem}
				\item[\textbf{Matthew}] ---
				\item[\textbf{Mark}] ---
				\item[\textbf{Luke}] ---
				\item[\textbf{John}] \hyperref[JOHN-2-1]{2:1--11}
			\end{compactitem}
			
			% 30
		\item Teaching with Authority and Casting Out a Devil in Capernaum
			\begin{compactitem}
				\item[\textbf{Matthew}] ---
				\item[\textbf{Mark}] \hyperref[MARK-1-21]{1:21--28}
				\item[\textbf{Luke}] \hyperref[LUKE-4-31]{4:31--37}
				\item[\textbf{John}] ---
			\end{compactitem}
			
	\columnbreak
			
			% 31
		\item \textit{Healing Peter's Mother-in-Law}
			\begin{compactitem}
				\item[\textbf{Matthew}] \hyperref[MATTHEW-8-14]{8:14--15}
				\item[\textbf{Mark}] \hyperref[MARK-1-29]{1:29--31}
				\item[\textbf{Luke}] \hyperref[LUKE-4-38]{4:38--39}
				\item[\textbf{John}] ---
			\end{compactitem}
			
			% 32
		\item \textit{Evening of Healing}
			\begin{compactitem}
				\item[\textbf{Matthew}] \hyperref[MATTHEW-8-16]{8:16--17}
				\item[\textbf{Mark}] \hyperref[MARK-1-32]{1:32--34}
				\item[\textbf{Luke}] \hyperref[LUKE-4-40]{4:40--41}
				\item[\textbf{John}] ---
			\end{compactitem}
			
			% 33
		\item Departure from Capernaum
			\begin{compactitem}
				\item[\textbf{Matthew}] ---
				\item[\textbf{Mark}] \hyperref[MARK-1-35]{1:35--38}
				\item[\textbf{Luke}] \hyperref[LUKE-4-42]{4:42--43}
				\item[\textbf{John}] ---
			\end{compactitem}
			
			% 34
		\item \textit{Preaching and Healing in Galilee}
			\begin{compactitem}
				\item[\textbf{Matthew}] \hyperref[MATTHEW-4-23]{4:23--25}
				\item[\textbf{Mark}] \hyperref[MARK-1-39]{1:39--44}
				\item[\textbf{Luke}] \hyperref[LUKE-4-42]{4:42--44}; \hyperref[LUKE-6-17]{6:17--19}
				\item[\textbf{John}] ---
			\end{compactitem}
			
			% 35
		\item Miraculous Catch of Fish
			\begin{compactitem}
				\item[\textbf{Matthew}] ---
				\item[\textbf{Mark}] ---
				\item[\textbf{Luke}] \hyperref[LUKE-5-1]{5:1--11}
				\item[\textbf{John}] ---
			\end{compactitem}
			
			% 36
		\item Jesus Instructs Nicodemus to Be Born of the Spirit
			\begin{compactitem}
				\item[\textbf{Matthew}] ---
				\item[\textbf{Mark}] ---
				\item[\textbf{Luke}] ---
				\item[\textbf{John}] \hyperref[JOHN-3-1]{3:1--21}
			\end{compactitem}
			
			% 37
		\item Sermon on the Mount Introduction
			\begin{compactitem}
				\item[\textbf{Matthew}] \hyperref[MATTHEW-5-1]{5:1--2}
				\item[\textbf{Mark}] ---
				\item[\textbf{Luke}] ---
				\item[\textbf{John}] ---
			\end{compactitem}
			
			% 38
		\item The Beatitudes
			\begin{compactitem}
				\item[\textbf{Matthew}] \hyperref[MATTHEW-5-3]{5:3--12}
				\item[\textbf{Mark}] ---
				\item[\textbf{Luke}] \hyperref[LUKE-6-20]{6:20--23}
				\item[\textbf{John}] ---
			\end{compactitem}
			
			% 39
		\item \textit{Salt of the Earth, Light of the World}
			\begin{compactitem}
				\item[\textbf{Matthew}] \hyperref[MATTHEW-5-13]{5:13--16}; \hyperref[MATTHEW-6-22]{6:22--23}
				\item[\textbf{Mark}] \hyperref[MARK-9-50]{9:50}
				\item[\textbf{Luke}] \hyperref[LUKE-14-34]{14:34--35}
				\item[\textbf{John}] ---
			\end{compactitem}
			
			% 40
		\item Jesus Came to Fulfil the Law
			\begin{compactitem}
				\item[\textbf{Matthew}] \hyperref[MATTHEW-5-17]{5:17--20}
				\item[\textbf{Mark}] ---
				\item[\textbf{Luke}] ---
				\item[\textbf{John}] ---
			\end{compactitem}
		
	\columnbreak
			
			% 41
		\item Be Not Angry, Agree with Adversary
			\begin{compactitem}
				\item[\textbf{Matthew}] \hyperref[MATTHEW-5-21]{5:21--26}
				\item[\textbf{Mark}] ---
				\item[\textbf{Luke}] \hyperref[LUKE-12-57]{12:57--59}
				\item[\textbf{John}] ---
			\end{compactitem}
			
			% 42
		\item The Higher Law on Adultery
			\begin{compactitem}
				\item[\textbf{Matthew}] \hyperref[MATTHEW-5-27]{5:27--30}
				\item[\textbf{Mark}] ---
				\item[\textbf{Luke}] ---
				\item[\textbf{John}] ---
			\end{compactitem}
			
			% 43
		\item \textit{The Higher Law on Divorce}
			\begin{compactitem}
				\item[\textbf{Matthew}] \hyperref[MATTHEW-5-31]{5:31--32}; \hyperref[MATTHEW-19-9]{19:9}
				\item[\textbf{Mark}] \hyperref[MARK-10-11]{10:11--12}
				\item[\textbf{Luke}] \hyperref[LUKE-16-18]{16:18}
				\item[\textbf{John}] ---
			\end{compactitem}
			
			% 44
		\item Swear Not at All
			\begin{compactitem}
				\item[\textbf{Matthew}] \hyperref[MATTHEW-5-33]{5:33--37}
				\item[\textbf{Mark}] ---
				\item[\textbf{Luke}] ---
				\item[\textbf{John}] ---
			\end{compactitem}
			
			% 45
		\item The Higher Law on Retaliation
			\begin{compactitem}
				\item[\textbf{Matthew}] \hyperref[MATTHEW-5-38]{5:38--42}
				\item[\textbf{Mark}] ---
				\item[\textbf{Luke}] \hyperref[LUKE-6-29]{6:29--30}
				\item[\textbf{John}] ---
			\end{compactitem}
			
			% 46
		\item Love Your Enemies, Bless Them that Curse You
			\begin{compactitem}
				\item[\textbf{Matthew}] \hyperref[MATTHEW-5-43]{5:43--48}; \hyperref[MATTHEW-7-12]{7:12}
				\item[\textbf{Mark}] ---
				\item[\textbf{Luke}] \hyperref[LUKE-6-27]{6:27--36}
				\item[\textbf{John}] ---
			\end{compactitem}
			
			% 47
		\item Give Alms in Secret
			\begin{compactitem}
				\item[\textbf{Matthew}] \hyperref[MATTHEW-6-1]{6:1--4}
				\item[\textbf{Mark}] ---
				\item[\textbf{Luke}] ---
				\item[\textbf{John}] ---
			\end{compactitem}
			
			% 48
		\item On Prayer
			\begin{compactitem}
				\item[\textbf{Matthew}] \hyperref[MATTHEW-6-5]{6:5--15}; \hyperref[MATTHEW-7-7]{7:7--11}
				\item[\textbf{Mark}] ---
				\item[\textbf{Luke}] \hyperref[LUKE-11-1]{11:1--13}
				\item[\textbf{John}] ---
			\end{compactitem}
			
			% 49
		\item The Lord's Prayer
			\begin{compactitem}
				\item[\textbf{Matthew}] \hyperref[MATTHEW-6-9]{6:9--15}
				\item[\textbf{Mark}] ---
				\item[\textbf{Luke}] \hyperref[LUKE-11-2]{11:2--4}
				\item[\textbf{John}] ---
			\end{compactitem}
			
			% 50
		\item On Fasting
			\begin{compactitem}
				\item[\textbf{Matthew}] \hyperref[MATTHEW-6-16]{6:16--18}
				\item[\textbf{Mark}] ---
				\item[\textbf{Luke}] ---
				\item[\textbf{John}] ---
			\end{compactitem}
			
	\columnbreak
			
			% 51
		\item Lay Up Treasures in Heaven
			\begin{compactitem}
				\item[\textbf{Matthew}] \hyperref[MATTHEW-6-19]{6:19--21}
				\item[\textbf{Mark}] ---
				\item[\textbf{Luke}] \hyperref[LUKE-12-33]{12:33--34}
				\item[\textbf{John}] ---
			\end{compactitem}
			
			% 52
		\item The Eye is the Light of the Body
			\begin{compactitem}
				\item[\textbf{Matthew}] \hyperref[MATTHEW-6-22]{6:22--23}
				\item[\textbf{Mark}] ---
				\item[\textbf{Luke}] \hyperref[LUKE-11-34]{11:34--36}
				\item[\textbf{John}] ---
			\end{compactitem}
			
			% 53
		\item No Man Can Serve Both God and Mammon
			\begin{compactitem}
				\item[\textbf{Matthew}] \hyperref[MATTHEW-6-24]{6:24}
				\item[\textbf{Mark}] ---
				\item[\textbf{Luke}] \hyperref[LUKE-16-13]{16:13}
				\item[\textbf{John}] ---
			\end{compactitem}
			
			% 54
		\item Take No Thought for the Morrow
			\begin{compactitem}
				\item[\textbf{Matthew}] \hyperref[MATTHEW-6-25]{6:25--34}; \hyperref[MATTHEW-6-19]{19--21}
				\item[\textbf{Mark}] ---
				\item[\textbf{Luke}] \hyperref[LUKE-12-22]{12:22--34}
				\item[\textbf{John}] ---
			\end{compactitem}
			
			% 55
		\item Ye Shall Be Judged with What Judgment Ye Judge
			\begin{compactitem}
				\item[\textbf{Matthew}] \hyperref[MATTHEW-7-1]{7:1--5}
				\item[\textbf{Mark}] ---
				\item[\textbf{Luke}] \hyperref[LUKE-6-37]{6:37--42}
				\item[\textbf{John}] ---
			\end{compactitem}
			
			% 56
		\item Cast Not Your Pears before Swine
			\begin{compactitem}
				\item[\textbf{Matthew}] \hyperref[MATTHEW-7-6]{7:6}
				\item[\textbf{Mark}] ---
				\item[\textbf{Luke}] ---
				\item[\textbf{John}] ---
			\end{compactitem}
			
			% 57
		\item Ask, Seek, Knock
			\begin{compactitem}
				\item[\textbf{Matthew}] \hyperref[MATTHEW-7-7]{7:7--12}
				\item[\textbf{Mark}] ---
				\item[\textbf{Luke}] \hyperref[LUKE-11-9]{11:9--13}
				\item[\textbf{John}] ---
			\end{compactitem}

			% 58
		\item The Golden Rule
			\begin{compactitem}
				\item[\textbf{Matthew}] \hyperref[MATTHEW-7-12]{7:12}
				\item[\textbf{Mark}] ---
				\item[\textbf{Luke}] \hyperref[LUKE-6-31]{6:31}
				\item[\textbf{John}] ---
			\end{compactitem}
			
			% 59
		\item Strait is the Gate, Narrow the Way
			\begin{compactitem}
				\item[\textbf{Matthew}] \hyperref[MATTHEW-7-13]{7:13--14}; \hyperref[MATTHEW-7-21]{21--23}
				\item[\textbf{Mark}] ---
				\item[\textbf{Luke}] \hyperref[LUKE-13-22]{13:22--30}
				\item[\textbf{John}] ---
			\end{compactitem}
			
			% 60
		\item A Tree Bears Only One Kind of Fruit
			\begin{compactitem}
				\item[\textbf{Matthew}] \hyperref[MATTHEW-7-15]{7:15--20}; \hyperref[MATTHEW-12-34]{12:34--35}
				\item[\textbf{Mark}] ---
				\item[\textbf{Luke}] \hyperref[LUKE-6-43]{6:43--45}
				\item[\textbf{John}] ---
			\end{compactitem}
			
	\columnbreak
			
			% 61
		\item I Never Knew You
			\begin{compactitem}
				\item[\textbf{Matthew}] \hyperref[MATTHEW-7-21]{7:21--23}
				\item[\textbf{Mark}] ---
				\item[\textbf{Luke}] \hyperref[LUKE-13-25]{13:25--27}
				\item[\textbf{John}] ---
			\end{compactitem}
			
			% 62
		\item Wise Man Builds House on Rock
			\begin{compactitem}
				\item[\textbf{Matthew}] \hyperref[MATTHEW-7-24]{7:24--28}
				\item[\textbf{Mark}] ---
				\item[\textbf{Luke}] \hyperref[LUKE-6-46]{6:46--49}
				\item[\textbf{John}] ---
			\end{compactitem}	
			
			% 63
		\item Jesus Taught as One with Authority
			\begin{compactitem}
				\item[\textbf{Matthew}] \hyperref[MATTHEW-7-28]{7:28--29}
				\item[\textbf{Mark}] ---
				\item[\textbf{Luke}] ---
				\item[\textbf{John}] ---
			\end{compactitem}
			
			% 64
		\item Samaritan Woman at the Well
			\begin{compactitem}
				\item[\textbf{Matthew}] ---
				\item[\textbf{Mark}] ---
				\item[\textbf{Luke}] ---
				\item[\textbf{John}] \hyperref[JOHN-4-1]{4:1--42}
			\end{compactitem}
			
			% 65
		\item \textit{Jesus Heals Faithful Leper}
			\begin{compactitem}
				\item[\textbf{Matthew}] \hyperref[MATTHEW-8-1]{8:1--4}
				\item[\textbf{Mark}] \hyperref[MARK-1-40]{1:40--45}
				\item[\textbf{Luke}] \hyperref[LUKE-5-12]{5:12--16}
				\item[\textbf{John}] ---
			\end{compactitem}
			
			% 66
		\item Jesus Heals Centurion's Servant
			\begin{compactitem}
				\item[\textbf{Matthew}] \hyperref[MATTHEW-8-5]{8:5--13}
				\item[\textbf{Mark}] ---
				\item[\textbf{Luke}] \hyperref[LUKE-7-1]{7:1--10}
				\item[\textbf{John}] \hyperref[JOHN-4-43]{4:43--54}
			\end{compactitem}
			
			% 67
		\item Jesus Raises Widow's Son at Nain
			\begin{compactitem}
				\item[\textbf{Matthew}] ---
				\item[\textbf{Mark}] ---
				\item[\textbf{Luke}] \hyperref[LUKE-7-11]{7:11--17}
				\item[\textbf{John}] ---
			\end{compactitem}
			
			% 68
		\item \textit{Jesus Heals the Sick, Casts Out Devils}
			\begin{compactitem}
				\item[\textbf{Matthew}] \hyperref[MATTHEW-8-16]{8:16--17}
				\item[\textbf{Mark}] \hyperref[MARK-1-32]{1:32--34}
				\item[\textbf{Luke}] \hyperref[LUKE-4-40]{4:40--41}
				\item[\textbf{John}] ---
			\end{compactitem}
			
			% 69
		\item Would-be Followers
			\begin{compactitem}
				\item[\textbf{Matthew}] \hyperref[MATTHEW-8-18]{8:18--22}
				\item[\textbf{Mark}] ---
				\item[\textbf{Luke}] \hyperref[LUKE-9-57]{9:57--62}
				\item[\textbf{John}] ---
			\end{compactitem}
			
			% 70
		\item \textit{Stilling the Tempest}
			\begin{compactitem}
				\item[\textbf{Matthew}] \hyperref[MATTHEW-8-23]{8:23--27}
				\item[\textbf{Mark}] \hyperref[MARK-4-36]{4:36--41}
				\item[\textbf{Luke}] \hyperref[LUKE-8-22]{8:22--25}
				\item[\textbf{John}] ---
			\end{compactitem}
			
	\columnbreak

			% 71
		\item \textit{Gadarene Demoniacs}
			\begin{compactitem}
				\item[\textbf{Matthew}] \hyperref[MATTHEW-8-28]{8:28--34}
				\item[\textbf{Mark}] \hyperref[MARK-5-1]{5:1--20}
				\item[\textbf{Luke}] \hyperref[LUKE-8-26]{8:26--39}
				\item[\textbf{John}] ---
			\end{compactitem}
			
			% 72
		\item \textit{Easier to Heal Sick or Forgive Sins?}
			\begin{compactitem}
				\item[\textbf{Matthew}] \hyperref[MATTHEW-9-1]{9:1--8}
				\item[\textbf{Mark}] \hyperref[MARK-2-1]{2:1--12}
				\item[\textbf{Luke}] \hyperref[LUKE-5-17]{5:17--26}
				\item[\textbf{John}] ---
			\end{compactitem}
			
			% 73
		\item \textit{Matthew and Tax Collecting}
			\begin{compactitem}
				\item[\textbf{Matthew}] \hyperref[MATTHEW-9-9]{9:9--13}
				\item[\textbf{Mark}] \hyperref[MARK-2-13]{2:13--17}
				\item[\textbf{Luke}] \hyperref[LUKE-5-27]{5:27--32}
				\item[\textbf{John}] ---
			\end{compactitem}
			
			% 74
		\item \textit{Question about Fasting}
			\begin{compactitem}
				\item[\textbf{Matthew}] \hyperref[MATTHEW-9-14]{9:14--17}
				\item[\textbf{Mark}] \hyperref[MARK-2-18]{2:18--22}
				\item[\textbf{Luke}] \hyperref[LUKE-5-33]{5:33--39}
				\item[\textbf{John}] ---
			\end{compactitem}
			
			% 75
		\item \textit{Jairus' Daughter}
			\begin{compactitem}
				\item[\textbf{Matthew}] \hyperref[MATTHEW-9-18]{9:18--19}, \hyperref[MATTHEW-9-23]{9:23--26}
				\item[\textbf{Mark}] \hyperref[MARK-5-21]{5:21--24}, \hyperref[MARK-5-35]{5:35--43}
				\item[\textbf{Luke}] \hyperref[LUKE-8-41]{8:41--42}, \hyperref[LUKE-8-49]{8:49--56}
				\item[\textbf{John}] ---
			\end{compactitem}
			
			% 76
		\item \textit{Woman Who Touched Jesus}
			\begin{compactitem}
				\item[\textbf{Matthew}] \hyperref[MATTHEW-9-20]{9:20--22}
				\item[\textbf{Mark}] \hyperref[MARK-5-25]{5:25--34}
				\item[\textbf{Luke}] \hyperref[LUKE-8-43]{8:43--48}
				\item[\textbf{John}] ---
			\end{compactitem}
			
			% 77
		\item Two Blind Men Healed
			\begin{compactitem}
				\item[\textbf{Matthew}] \hyperref[MATTHEW-9-27]{9:27--31}
				\item[\textbf{Mark}] ---
				\item[\textbf{Luke}] ---
				\item[\textbf{John}] ---
			\end{compactitem}
			
			% 78
		\item Dumb Man Healed
			\begin{compactitem}
				\item[\textbf{Matthew}] \hyperref[MATTHEW-9-32]{9:32--34}
				\item[\textbf{Mark}] ---
				\item[\textbf{Luke}] ---
				\item[\textbf{John}] ---
			\end{compactitem}
			
			% 79
		\item Compassion of Jesus
			\begin{compactitem}
				\item[\textbf{Matthew}] \hyperref[MATTHEW-9-35]{9:35--38}
				\item[\textbf{Mark}] ---
				\item[\textbf{Luke}] ---
				\item[\textbf{John}] ---
			\end{compactitem}
			
			% 80
		\item \textit{Calling the Twelve}
			\begin{compactitem}
				\item[\textbf{Matthew}] \hyperref[MATTHEW-10-1]{10:1--4}
				\item[\textbf{Mark}] \hyperref[MARK-3-13]{3:13--19}
				\item[\textbf{Luke}] \hyperref[LUKE-6-12]{6:12--16}
				\item[\textbf{John}] ---
			\end{compactitem}
			
	\columnbreak
	
			% 81
		\item \textit{Mission of the Twelve}
			\begin{compactitem}
				\item[\textbf{Matthew}] \hyperref[MATTHEW-10-1]{10:1--15}
				\item[\textbf{Mark}] \hyperref[MARK-6-7]{6:7--13}
				\item[\textbf{Luke}] \hyperref[LUKE-9-1]{9:1--6}
				\item[\textbf{John}] ---
			\end{compactitem}
			
			% 82
		\item \textit{Coming Persecutions}
			\begin{compactitem}
				\item[\textbf{Matthew}] \hyperref[MATTHEW-10-16]{10:16--25}
				\item[\textbf{Mark}] \hyperref[MARK-13-9]{13:9--13}
				\item[\textbf{Luke}] \hyperref[LUKE-21-12]{21:12--17}
				\item[\textbf{John}] ---
			\end{compactitem}
			
			% 83
		\item Confessing Christ before Men
			\begin{compactitem}
				\item[\textbf{Matthew}] \hyperref[MATTHEW-10-32]{10:32--33}; \hyperref[MATTHEW-12-32]{12:32}
				\item[\textbf{Mark}] 
				\item[\textbf{Luke}] \hyperref[LUKE-12-8]{12:8--12}
				\item[\textbf{John}] ---
			\end{compactitem}
			
			% 84
		\item Whom to Fear
			\begin{compactitem}
				\item[\textbf{Matthew}] \hyperref[MATTHEW-10-26]{10:26--33}
				\item[\textbf{Mark}] ---
				\item[\textbf{Luke}] \hyperref[LUKE-12-4]{12:4--7}
				\item[\textbf{John}] ---
			\end{compactitem}
			
			% 85
		\item Not Peace But a Sword
			\begin{compactitem}
				\item[\textbf{Matthew}] \hyperref[MATTHEW-10-34]{10:34--39}
				\item[\textbf{Mark}] ---
				\item[\textbf{Luke}] \hyperref[LUKE-12-49]{12:49--53}; \hyperref[LUKE-14-26]{14:26--27}
				\item[\textbf{John}] ---
			\end{compactitem}
			
			% 86
		\item Conditions of Discipleship
			\begin{compactitem}
				\item[\textbf{Matthew}] \hyperref[MATTHEW-10-37]{10:37--38}
				\item[\textbf{Mark}] ---
				\item[\textbf{Luke}] \hyperref[LUKE-14-26]{14:26--33}
				\item[\textbf{John}] ---
			\end{compactitem}
			
			% 87
		\item Rewards of Discipleship
			\begin{compactitem}
				\item[\textbf{Matthew}] \hyperref[MATTHEW-10-40]{10:40--11:1}
				\item[\textbf{Mark}] \hyperref[MARK-9-41]{9:41}
				\item[\textbf{Luke}] ---
				\item[\textbf{John}] ---
			\end{compactitem}
			
			% 88
		\item Messengers from John the Baptist
			\begin{compactitem}
				\item[\textbf{Matthew}] \hyperref[MATTHEW-11-2]{11:2--19}
				\item[\textbf{Mark}] ---
				\item[\textbf{Luke}] \hyperref[LUKE-7-18]{7:18--23}
				\item[\textbf{John}] ---
			\end{compactitem}
			
			% 89
		\item Christ's Testimony of John
			\begin{compactitem}
				\item[\textbf{Matthew}] ---
				\item[\textbf{Mark}] ---
				\item[\textbf{Luke}] \hyperref[LUKE-7-24]{7:24--35}
				\item[\textbf{John}] \hyperref[JOHN-3-22]{3:22--30}
			\end{compactitem}
			
			% 90
		\item Law and the Kingdom of God
			\begin{compactitem}
				\item[\textbf{Matthew}] \hyperref[MATTHEW-11-12]{11:12--13}
				\item[\textbf{Mark}] ---
				\item[\textbf{Luke}] \hyperref[LUKE-16-14]{16:14--18}
				\item[\textbf{John}] ---
			\end{compactitem}
			
	\columnbreak
	
			% 91
		\item On Testimony
			\begin{compactitem}
				\item[\textbf{Matthew}] ---
				\item[\textbf{Mark}] ---
				\item[\textbf{Luke}] ---
				\item[\textbf{John}] \hyperref[JOHN-3-31]{3:31--35}
			\end{compactitem}
			
			% 92
		\item Woes to Unrepentant Cities
			\begin{compactitem}
				\item[\textbf{Matthew}] \hyperref[MATTHEW-11-20]{11:20--24}
				\item[\textbf{Mark}] ---
				\item[\textbf{Luke}] \hyperref[LUKE-10-13]{10:13--15}
				\item[\textbf{John}] ---
			\end{compactitem}
			
			% 93
		\item Christ's Thanksgiving Prayer
			\begin{compactitem}
				\item[\textbf{Matthew}] \hyperref[MATTHEW-11-25]{11:25--27}
				\item[\textbf{Mark}] ---
				\item[\textbf{Luke}] ---
				\item[\textbf{John}] ---
			\end{compactitem}
			
			% 94
		\item Healing at the Pool of Bethesda
			\begin{compactitem}
				\item[\textbf{Matthew}] ---
				\item[\textbf{Mark}] ---
				\item[\textbf{Luke}] ---
				\item[\textbf{John}] \hyperref[JOHN-5-2]{5:2--18}
			\end{compactitem}
			
			% 95
		\item Authority of the Son
			\begin{compactitem}
				\item[\textbf{Matthew}] ---
				\item[\textbf{Mark}] ---
				\item[\textbf{Luke}] ---
				\item[\textbf{John}] \hyperref[JOHN-5-19]{5:19--29}
			\end{compactitem}
			
			% 96
		\item Come to Me and Rest
			\begin{compactitem}
				\item[\textbf{Matthew}] \hyperref[MATTHEW-11-25]{11:25--30}
				\item[\textbf{Mark}] ---
				\item[\textbf{Luke}] ---
				\item[\textbf{John}] ---
			\end{compactitem}
			
			% 97
		\item Witnesses to Jesus
			\begin{compactitem}
				\item[\textbf{Matthew}] ---
				\item[\textbf{Mark}] ---
				\item[\textbf{Luke}] ---
				\item[\textbf{John}] \hyperref[JOHN-5-30]{5:30--47}
			\end{compactitem}
			
			% 98
		\item \textit{Corn on the Sabbath}
			\begin{compactitem}
				\item[\textbf{Matthew}] \hyperref[MATTHEW-12-1]{12:1--8}
				\item[\textbf{Mark}] \hyperref[MARK-2-23]{2:23--28}
				\item[\textbf{Luke}] \hyperref[LUKE-6-1]{6:1--5}
				\item[\textbf{John}] ---
			\end{compactitem}
			
			% 99
		\item \textit{Healing a Withered Hand}
			\begin{compactitem}
				\item[\textbf{Matthew}] \hyperref[MATTHEW-12-9]{12:9--14}
				\item[\textbf{Mark}] \hyperref[MARK-3-1]{3:1--6}
				\item[\textbf{Luke}] \hyperref[LUKE-6-6]{6:6--11}
				\item[\textbf{John}] ---
			\end{compactitem}
			
			% 100
		\item The Chosen Healer-Servant
			\begin{compactitem}
				\item[\textbf{Matthew}] \hyperref[MATTHEW-12-15]{12:15--21}
				\item[\textbf{Mark}] ---
				\item[\textbf{Luke}] ---
				\item[\textbf{John}] ---
			\end{compactitem}
			
	\columnbreak
			
			% 101
		\item The Seaside Multitude
			\begin{compactitem}
				\item[\textbf{Matthew}] ---
				\item[\textbf{Mark}] \hyperref[MARK-3-7]{3:7--12}
				\item[\textbf{Luke}] ---
				\item[\textbf{John}] ---
			\end{compactitem}
			
			% 102
		\item The Servant from Capernaum
			\begin{compactitem}
				\item[\textbf{Matthew}] ---
				\item[\textbf{Mark}] ---
				\item[\textbf{Luke}] \hyperref[LUKE-7-1]{7:1--10}
				\item[\textbf{John}] ---
			\end{compactitem}
			
			% 103
		\item A Sinful Woman Forgiven
			\begin{compactitem}
				\item[\textbf{Matthew}] ---
				\item[\textbf{Mark}] ---
				\item[\textbf{Luke}] \hyperref[LUKE-7-36]{7:36--50}
				\item[\textbf{John}] ---
			\end{compactitem}
			
			% 104
		\item Women Ministering to Jesus
			\begin{compactitem}
				\item[\textbf{Matthew}] ---
				\item[\textbf{Mark}] ---
				\item[\textbf{Luke}] \hyperref[LUKE-8-2]{8:2--3}
				\item[\textbf{John}] ---
			\end{compactitem}
			
			% 105
		\item The Accusation of Devilishness
			\begin{compactitem}
				\item[\textbf{Matthew}] \hyperref[MATTHEW-12-22]{12:22--23}
				\item[\textbf{Mark}] ---
				\item[\textbf{Luke}] \hyperref[LUKE-11-14]{11:14--15}
				\item[\textbf{John}] ---
			\end{compactitem}
			
			% 106
		\item \textit{The Beelzebub Controversy}
			\begin{compactitem}
				\item[\textbf{Matthew}] \hyperref[MATTHEW-12-22]{12:22--23}
				\item[\textbf{Mark}] \hyperref[MARK-3-22]{3:22--30}
				\item[\textbf{Luke}] \hyperref[LUKE-11-14]{11:14--23}
				\item[\textbf{John}] ---
			\end{compactitem}
			
			% 107
		\item A Tree and Its Fruits
			\begin{compactitem}
				\item[\textbf{Matthew}] \hyperref[MATTHEW-12-33]{12:33--37}
				\item[\textbf{Mark}] ---
				\item[\textbf{Luke}] ---
				\item[\textbf{John}] ---
			\end{compactitem}
			
			% 108
		\item \textit{The Demand for a Sign}
			\begin{compactitem}
				\item[\textbf{Matthew}] \hyperref[MATTHEW-12-38]{12:38--42}
				\item[\textbf{Mark}] \hyperref[MARK-8-11]{8:11--12}
				\item[\textbf{Luke}] \hyperref[LUKE-11-29]{11:29--32}
				\item[\textbf{John}] ---
			\end{compactitem}
			
			% 109
		\item Return of the Unclean Spirit
			\begin{compactitem}
				\item[\textbf{Matthew}] \hyperref[MATTHEW-12-43]{12:43--45}
				\item[\textbf{Mark}] ---
				\item[\textbf{Luke}] \hyperref[LUKE-11-24]{11:24--26}
				\item[\textbf{John}] ---
			\end{compactitem}
			
			% 110
		\item \textit{The Brothers of Jesus}
			\begin{compactitem}
				\item[\textbf{Matthew}] \hyperref[MATTHEW-12-46]{12:46--50}
				\item[\textbf{Mark}] \hyperref[MARK-3-31]{3:31--35}
				\item[\textbf{Luke}] \hyperref[LUKE-8-19]{8:19--21}
				\item[\textbf{John}] \hyperref[JOHN-2-12]{2:12}
			\end{compactitem}
			
		% 111
		\item \textit{Parable of the Sower}
			\begin{compactitem}
				\item[\textbf{Matthew}] \hyperref[MATTHEW-13-1]{13:1--9}
				\item[\textbf{Mark}] \hyperref[MARK-4-3]{4:3--9}
				\item[\textbf{Luke}] \hyperref[LUKE-8-4]{8:4--8}
				\item[\textbf{John}] ---
			\end{compactitem}
			
		% 112
		\item \textit{Purpose of Parables}
			\begin{compactitem}
				\item[\textbf{Matthew}] \hyperref[MATTHEW-13-10]{13:10--17}
				\item[\textbf{Mark}] \hyperref[MARK-4-10]{4:10--12}
				\item[\textbf{Luke}] \hyperref[LUKE-8-9]{8:9--10}
				\item[\textbf{John}] ---
			\end{compactitem}
			
		% 113
		\item Blessedness of Discipleship
			\begin{compactitem}
				\item[\textbf{Matthew}] \hyperref[MATTHEW-13-16]{13:16--17}
				\item[\textbf{Mark}] ---
				\item[\textbf{Luke}] \hyperref[LUKE-10-23]{10:23--24}
				\item[\textbf{John}] ---
			\end{compactitem}
			
		% 114
		\item \textit{The Sower Interpreted}
			\begin{compactitem}
				\item[\textbf{Matthew}] \hyperref[MATTHEW-13-18]{13:18--23}
				\item[\textbf{Mark}] \hyperref[MARK-4-13]{4:13--20}
				\item[\textbf{Luke}] \hyperref[LUKE-8-11]{8:11--15}
				\item[\textbf{John}] ---
			\end{compactitem}
			
		% 115
		\item Light under a Bushel
			\begin{compactitem}
				\item[\textbf{Matthew}] ---
				\item[\textbf{Mark}] \hyperref[MARK-4-21]{4:21--25}
				\item[\textbf{Luke}] \hyperref[LUKE-8-16]{8:16--18}; \hyperref[LUKE-11-33]{11:33}
				\item[\textbf{John}] ---
			\end{compactitem}
			
		% 116
		\item Weeds among the Wheat
			\begin{compactitem}
				\item[\textbf{Matthew}] \hyperref[MATTHEW-13-24]{13:24--30}
				\item[\textbf{Mark}] ---
				\item[\textbf{Luke}] ---
				\item[\textbf{John}] ---
			\end{compactitem}
			
		% 117
		\item The Growing Seed
			\begin{compactitem}
				\item[\textbf{Matthew}] ---
				\item[\textbf{Mark}] \hyperref[MARK-4-26]{4:26--29}
				\item[\textbf{Luke}] ---
				\item[\textbf{John}] ---
			\end{compactitem}
			
		% 118
		\item The Wheat and the Tares
			\begin{compactitem}
				\item[\textbf{Matthew}] \hyperref[MATTHEW-13-24]{13:24--30}
				\item[\textbf{Mark}] ---
				\item[\textbf{Luke}] ---
				\item[\textbf{John}] ---
			\end{compactitem}
			
		% 119
		\item \textit{The Mustard Seed}
			\begin{compactitem}
				\item[\textbf{Matthew}] \hyperref[MATTHEW-13-31]{13:31--32}
				\item[\textbf{Mark}] \hyperref[MARK-4-30]{4:30--32}
				\item[\textbf{Luke}] \hyperref[LUKE-13-18]{13:18--19}
				\item[\textbf{John}] ---
			\end{compactitem}
			
		% 120
		\item The Leaven
			\begin{compactitem}
				\item[\textbf{Matthew}] \hyperref[MATTHEW-13-33]{13:33}
				\item[\textbf{Mark}] ---
				\item[\textbf{Luke}] \hyperref[LUKE-13-20]{13:20--21}
				\item[\textbf{John}] ---
			\end{compactitem}
			
	\columnbreak

		% 121
		\item The Use of Parables
			\begin{compactitem}
				\item[\textbf{Matthew}] \hyperref[MATTHEW-13-34]{13:34--35}
				\item[\textbf{Mark}] \hyperref[MARK-4-33]{4:33--34}
				\item[\textbf{Luke}] ---
				\item[\textbf{John}] ---
			\end{compactitem}
			
		% 122
		\item Interpretation of the Tares
			\begin{compactitem}
				\item[\textbf{Matthew}] \hyperref[MATTHEW-13-36]{13:36--43}
				\item[\textbf{Mark}] ---
				\item[\textbf{Luke}] ---
				\item[\textbf{John}] ---
			\end{compactitem}
			
		% 123
		\item Pearl of Great Price
			\begin{compactitem}
				\item[\textbf{Matthew}] \hyperref[MATTHEW-13-44]{13:44--46}
				\item[\textbf{Mark}] ---
				\item[\textbf{Luke}] ---
				\item[\textbf{John}] ---
			\end{compactitem}
			
		% 124
		\item Parable of the Great Net
			\begin{compactitem}
				\item[\textbf{Matthew}] \hyperref[MATTHEW-13-47]{13:47--50}
				\item[\textbf{Mark}] ---
				\item[\textbf{Luke}] ---
				\item[\textbf{John}] ---
			\end{compactitem}
			
		% 125
		\item End of the Parable Sermon
			\begin{compactitem}
				\item[\textbf{Matthew}] \hyperref[MATTHEW-13-51]{13:51--52}
				\item[\textbf{Mark}] ---
				\item[\textbf{Luke}] ---
				\item[\textbf{John}] ---
			\end{compactitem}
			
		% 126
		\item \textit{Rejection at Nazareth}
			\begin{compactitem}
				\item[\textbf{Matthew}] \hyperref[MATTHEW-13-53]{13:53--58}
				\item[\textbf{Mark}] \hyperref[MARK-6-1]{6:1--6}
				\item[\textbf{Luke}] \hyperref[LUKE-4-16]{4:16--30}
				\item[\textbf{John}] \hyperref[JOHN-7-1]{7:1--9}
			\end{compactitem}
			
		% 127
		\item Herod's Opinion of Christ
			\begin{compactitem}
				\item[\textbf{Matthew}] \hyperref[MATTHEW-14-1]{14:1--2}
				\item[\textbf{Mark}] \hyperref[MARK-6-14]{6:14--16}
				\item[\textbf{Luke}] ---
				\item[\textbf{John}] ---
			\end{compactitem}
			
		% 128
		\item \textit{Death of John the Baptist}
			\begin{compactitem}
				\item[\textbf{Matthew}] \hyperref[MATTHEW-14-1]{14:1--12}
				\item[\textbf{Mark}] \hyperref[MARK-6-14]{6:14--29}
				\item[\textbf{Luke}] \hyperref[LUKE-9-7]{9:7--9}
				\item[\textbf{John}] ---
			\end{compactitem}
			
		% 129
		\item \textit{Feeding the Five Thousand}
			\begin{compactitem}
				\item[\textbf{Matthew}] \hyperref[MATTHEW-14-13]{14:13--21}
				\item[\textbf{Mark}] \hyperref[MARK-6-30]{6:30--44}
				\item[\textbf{Luke}] \hyperref[LUKE-9-10]{9:10--17}
				\item[\textbf{John}] \hyperref[JOHN-6-1]{6:1--14}
			\end{compactitem}
			
		% 130
		\item Walking on the Water
			\begin{compactitem}
				\item[\textbf{Matthew}] \hyperref[MATTHEW-14-22]{14:22--33}
				\item[\textbf{Mark}] \hyperref[MARK-6-45]{6:45--51}
				\item[\textbf{Luke}] ---
				\item[\textbf{John}] \hyperref[JOHN-6-15]{6:15--21}
			\end{compactitem}
	
	\columnbreak
	
		% 131
		\item Healings at Gennesaret
			\begin{compactitem}
				\item[\textbf{Matthew}] \hyperref[MATTHEW-14-34]{14:34--36}
				\item[\textbf{Mark}] \hyperref[MARK-6-53]{6:53--56}
				\item[\textbf{Luke}] ---
				\item[\textbf{John}] ---
			\end{compactitem}
			
		% 132
		\item Traditions of the Jews
			\begin{compactitem}
				\item[\textbf{Matthew}] \hyperref[MATTHEW-15-1]{15:1--20}
				\item[\textbf{Mark}] \hyperref[MARK-7-1]{7:1--23}
				\item[\textbf{Luke}] ---
				\item[\textbf{John}] ---
			\end{compactitem}
			
		% 133
		\item The Gentile Woman
			\begin{compactitem}
				\item[\textbf{Matthew}] \hyperref[MATTHEW-15-21]{15:21--28}
				\item[\textbf{Mark}] \hyperref[MARK-7-24]{7:24--30}
				\item[\textbf{Luke}] ---
				\item[\textbf{John}] ---
			\end{compactitem}
			
		% 134
		\item Healing of a Deaf Mute
			\begin{compactitem}
				\item[\textbf{Matthew}] ---
				\item[\textbf{Mark}] \hyperref[MARK-7-32]{7:32--37}
				\item[\textbf{Luke}] ---
				\item[\textbf{John}] ---
			\end{compactitem}
			
		% 135
		\item Healing of Many Persons
			\begin{compactitem}
				\item[\textbf{Matthew}] \hyperref[MATTHEW-15-29]{15:29--31}
				\item[\textbf{Mark}] ---
				\item[\textbf{Luke}] ---
				\item[\textbf{John}] ---
			\end{compactitem}
			
		% 136
		\item Feeding the Four Thousand
			\begin{compactitem}
				\item[\textbf{Matthew}] \hyperref[MATTHEW-15-32]{15:32--39}
				\item[\textbf{Mark}] \hyperref[MARK-8-1]{8:1--9}
				\item[\textbf{Luke}] ---
				\item[\textbf{John}] ---
			\end{compactitem}
			
		% 137
		\item \textit{Demand for Signs of the Times}
			\begin{compactitem}
				\item[\textbf{Matthew}] \hyperref[MATTHEW-16-1]{16:1--4}
				\item[\textbf{Mark}] \hyperref[MARK-8-11]{8:11--12}
				\item[\textbf{Luke}] \hyperref[LUKE-12-54]{12:54--56}
				\item[\textbf{John}] ---
			\end{compactitem}
			
		% 138
		\item A Discourse on Leaven
			\begin{compactitem}
				\item[\textbf{Matthew}] \hyperref[MATTHEW-16-5]{16:5--12}
				\item[\textbf{Mark}] \hyperref[MARK-8-14]{8:14--21}
				\item[\textbf{Luke}] ---
				\item[\textbf{John}] ---
			\end{compactitem}
		
		% 139
		\item The True Bread of Life
			\begin{compactitem}
				\item[\textbf{Matthew}] ---
				\item[\textbf{Mark}] ---
				\item[\textbf{Luke}] ---
				\item[\textbf{John}] \hyperref[JOHN-6-22]{6:22---58}
			\end{compactitem}
			
		% 140
		\item Words of Eternal Life
			\begin{compactitem}
				\item[\textbf{Matthew}] ---
				\item[\textbf{Mark}] ---
				\item[\textbf{Luke}] ---
				\item[\textbf{John}] \hyperref[JOHN-6-60]{6:60--71}
			\end{compactitem}
	
	\columnbreak
	
		% 141
		\item Feast of the Tabernacles
			\begin{compactitem}
				\item[\textbf{Matthew}] ---
				\item[\textbf{Mark}] ---
				\item[\textbf{Luke}] ---
				\item[\textbf{John}] \hyperref[JOHN-7-10]{7:10--24}
			\end{compactitem}
			
		% 142
		\item Is This the Christ?
			\begin{compactitem}
				\item[\textbf{Matthew}] ---
				\item[\textbf{Mark}] ---
				\item[\textbf{Luke}] ---
				\item[\textbf{John}] \hyperref[JOHN-7-25]{7:25--31}
			\end{compactitem}
			
		% 143
		\item Officers Sent to Arrest Jesus
			\begin{compactitem}
				\item[\textbf{Matthew}] ---
				\item[\textbf{Mark}] ---
				\item[\textbf{Luke}] ---
				\item[\textbf{John}] \hyperref[JOHN-7-32]{7:32--36}
			\end{compactitem}
			
		% 144
		\item Rivers of Living Waters 
			\begin{compactitem}
				\item[\textbf{Matthew}] ---
				\item[\textbf{Mark}] ---
				\item[\textbf{Luke}] ---
				\item[\textbf{John}] \hyperref[JOHN-7-37]{7:37--39}
			\end{compactitem}
			
		% 145
		\item Division among the People
			\begin{compactitem}
				\item[\textbf{Matthew}] ---
				\item[\textbf{Mark}] ---
				\item[\textbf{Luke}] ---
				\item[\textbf{John}] \hyperref[JOHN-7-40]{7:40--44}
			\end{compactitem}
			
		% 146
		\item Unbelief of the Popular Leaders
			\begin{compactitem}
				\item[\textbf{Matthew}] ---
				\item[\textbf{Mark}] ---
				\item[\textbf{Luke}] ---
				\item[\textbf{John}] \hyperref[JOHN-7-45]{7:45--52}
			\end{compactitem}
			
		% 147
		\item The Woman Taken in Adultery
			\begin{compactitem}
				\item[\textbf{Matthew}] ---
				\item[\textbf{Mark}] ---
				\item[\textbf{Luke}] ---
				\item[\textbf{John}] \hyperref[JOHN-8-3]{8:3--11}
			\end{compactitem}
			
		% 148
		\item Light of the World
			\begin{compactitem}
				\item[\textbf{Matthew}] ---
				\item[\textbf{Mark}] ---
				\item[\textbf{Luke}] ---
				\item[\textbf{John}] \hyperref[JOHN-8-12]{8:12--20}; \hyperref[JOHN-9-1]{9:1--12}
			\end{compactitem}
			
		% 149
		\item Healing a Man Born Blind
			\begin{compactitem}
				\item[\textbf{Matthew}] ---
				\item[\textbf{Mark}] ---
				\item[\textbf{Luke}] ---
				\item[\textbf{John}] \hyperref[JOHN-9-1]{9:1--12}
			\end{compactitem}
			
		% 150
		\item The Blind Man at Bethsaida
			\begin{compactitem}
				\item[\textbf{Matthew}] ---
				\item[\textbf{Mark}] \hyperref[MARK-8-22]{8:22--26}
				\item[\textbf{Luke}] ---
				\item[\textbf{John}] ---
			\end{compactitem}
	
	\columnbreak
	
		% 151
		\item Pharisees Investigate
			\begin{compactitem}
				\item[\textbf{Matthew}] ---
				\item[\textbf{Mark}] ---
				\item[\textbf{Luke}] ---
				\item[\textbf{John}] \hyperref[JOHN-9-13]{9:13--34}
			\end{compactitem}
			
		% 152
		\item \textit{Thou Art the Christ}
			\begin{compactitem}
				\item[\textbf{Matthew}] \hyperref[MATTHEW-16-13]{16:13--20}
				\item[\textbf{Mark}] \hyperref[MARK-8-27]{8:27--30}
				\item[\textbf{Luke}] \hyperref[LUKE-9-18]{9:18--21}
				\item[\textbf{John}] ---
			\end{compactitem}
			
		% 153
		\item \textit{First Prediction of the Passion}
			\begin{compactitem}
				\item[\textbf{Matthew}] \hyperref[MATTHEW-16-21]{16:21--28}
				\item[\textbf{Mark}] \hyperref[MARK-8-31]{8:31--9:1}
				\item[\textbf{Luke}] \hyperref[LUKE-9-22]{9:22--27}
				\item[\textbf{John}] ---
			\end{compactitem}
			
		% 154
		\item Where I Am Going You Cannot Come
			\begin{compactitem}
				\item[\textbf{Matthew}] ---
				\item[\textbf{Mark}] ---
				\item[\textbf{Luke}] ---
				\item[\textbf{John}] \hyperref[JOHN-8-21]{8:21--30}
			\end{compactitem}
			
		% 155
		\item Conditions of Discipleship
			\begin{compactitem}
				\item[\textbf{Matthew}] \hyperref[MATTHEW-16-24]{16:24--28}
				\item[\textbf{Mark}] \hyperref[MARK-8-34]{8:34--9:1}
				\item[\textbf{Luke}] ---
				\item[\textbf{John}] \hyperref[JOHN-8-31]{8:31--38}
			\end{compactitem}
			
		% 156
		\item The Truth Shall Make You Free
			\begin{compactitem}
				\item[\textbf{Matthew}] ---
				\item[\textbf{Mark}] ---
				\item[\textbf{Luke}] ---
				\item[\textbf{John}] \hyperref[JOHN-8-32]{8:32}
			\end{compactitem}
			
		% 157
		\item Sons of the Devil
			\begin{compactitem}
				\item[\textbf{Matthew}] ---
				\item[\textbf{Mark}] ---
				\item[\textbf{Luke}] ---
				\item[\textbf{John}] \hyperref[JOHN-8-39]{8:39--47}
			\end{compactitem}
			
		% 158
		\item Before Abraham Was, I Am
			\begin{compactitem}
				\item[\textbf{Matthew}] ---
				\item[\textbf{Mark}] ---
				\item[\textbf{Luke}] ---
				\item[\textbf{John}] \hyperref[JOHN-8-48]{8:48--59}
			\end{compactitem}
			
		% 159
		\item \textit{The Son of God Transfigured}
			\begin{compactitem}
				\item[\textbf{Matthew}] \hyperref[MATTHEW-17-1]{17:1--13}
				\item[\textbf{Mark}] \hyperref[MARK-9-2]{9:2--13}
				\item[\textbf{Luke}] \hyperref[LUKE-9-28]{9:28--36}
				\item[\textbf{John}] ---
			\end{compactitem}
			
		% 160
		\item The Coming of Elijah
			\begin{compactitem}
				\item[\textbf{Matthew}] \hyperref[MATTHEW-17-10]{17:10--13}
				\item[\textbf{Mark}] \hyperref[MARK-9-11]{9:11--13}
				\item[\textbf{Luke}] ---
				\item[\textbf{John}] ---
			\end{compactitem}
	
	\columnbreak
	
		% 161
		\item \textit{Healing of an Epileptic Child}
			\begin{compactitem}
				\item[\textbf{Matthew}] \hyperref[MATTHEW-17-14]{17:14--21}
				\item[\textbf{Mark}] \hyperref[MARK-9-14]{9:14--29}
				\item[\textbf{Luke}] \hyperref[LUKE-9-38]{9:38--42}
				\item[\textbf{John}] ---
			\end{compactitem}
			
		% 162
		\item \textit{Second Prophecy of the Passion}
			\begin{compactitem}
				\item[\textbf{Matthew}] \hyperref[MATTHEW-17-22]{17:22--23}
				\item[\textbf{Mark}] \hyperref[MARK-9-30]{9:30--32}
				\item[\textbf{Luke}] \hyperref[LUKE-9-43]{9:43--45}
				\item[\textbf{John}] ---
			\end{compactitem}
			
		% 163
		\item The Temple Tax
			\begin{compactitem}
				\item[\textbf{Matthew}] \hyperref[MATTHEW-17-24]{17:24--27}
				\item[\textbf{Mark}] ---
				\item[\textbf{Luke}] ---
				\item[\textbf{John}] ---
			\end{compactitem}
			
		% 164
		\item \textit{Dispute about Greatness}
			\begin{compactitem}
				\item[\textbf{Matthew}] \hyperref[MATTHEW-18-1]{18:1--4}
				\item[\textbf{Mark}] \hyperref[MARK-9-33]{9:33--35}
				\item[\textbf{Luke}] \hyperref[LUKE-9-46]{9:46--48}; \hyperref[LUKE-22-24]{22:24--30}
				\item[\textbf{John}] ---
			\end{compactitem}
			
		% 165
		\item The Strange Exorcist
			\begin{compactitem}
				\item[\textbf{Matthew}] ---
				\item[\textbf{Mark}] \hyperref[MARK-9-38]{9:38--41}
				\item[\textbf{Luke}] \hyperref[LUKE-9-49]{9:49--50}
				\item[\textbf{John}] ---
			\end{compactitem}
			
		% 166
		\item \textit{About Taking Offense}
			\begin{compactitem}
				\item[\textbf{Matthew}] \hyperref[MATTHEW-18-6]{18:6--10}
				\item[\textbf{Mark}] \hyperref[MARK-9-42]{9:42--48}
				\item[\textbf{Luke}] \hyperref[LUKE-17-1]{17:1--2}
				\item[\textbf{John}] ---
			\end{compactitem}
			
		% 167
		\item Salt of the Earth, Purity
			\begin{compactitem}
				\item[\textbf{Matthew}] ---
				\item[\textbf{Mark}] \hyperref[MARK-9-49]{9:49--50}
				\item[\textbf{Luke}] ---
				\item[\textbf{John}] ---
			\end{compactitem}
			
		% 168
		\item The Lost Sheep
			\begin{compactitem}
				\item[\textbf{Matthew}] \hyperref[MATTHEW-18-11]{18:11--14}
				\item[\textbf{Mark}] ---
				\item[\textbf{Luke}] \hyperref[LUKE-15-1]{15:1--7}
				\item[\textbf{John}] ---
			\end{compactitem}
			
		% 169
		\item Parable of the Sheepfold
			\begin{compactitem}
				\item[\textbf{Matthew}] ---
				\item[\textbf{Mark}] ---
				\item[\textbf{Luke}] ---
				\item[\textbf{John}] \hyperref[JOHN-10-1]{10:1--6}
			\end{compactitem}
			
		% 170
		\item Spiritual Blindness
			\begin{compactitem}
				\item[\textbf{Matthew}] ---
				\item[\textbf{Mark}] ---
				\item[\textbf{Luke}] ---
				\item[\textbf{John}] \hyperref[JOHN-9-35]{9:35--41}
			\end{compactitem}
		
	\columnbreak
	
		% 171
		\item The Good Shepherd
			\begin{compactitem}
				\item[\textbf{Matthew}] ---
				\item[\textbf{Mark}] ---
				\item[\textbf{Luke}] ---
				\item[\textbf{John}] \hyperref[JOHN-10-7]{10:7--21}
			\end{compactitem}
			
		% 172
		\item Christ Rejected by the Jews
			\begin{compactitem}
				\item[\textbf{Matthew}] ---
				\item[\textbf{Mark}] ---
				\item[\textbf{Luke}] ---
				\item[\textbf{John}] \hyperref[JOHN-10-22]{10:22--42}
			\end{compactitem}
			
		% 173
		\item On Reproving a Brother
			\begin{compactitem}
				\item[\textbf{Matthew}] \hyperref[MATTHEW-18-15]{18:15--20}
				\item[\textbf{Mark}] ---
				\item[\textbf{Luke}] \hyperref[LUKE-17-3]{17:3--4}
				\item[\textbf{John}] ---
			\end{compactitem}
			
		% 174
		\item Parable of the Unmerciful Servant
			\begin{compactitem}
				\item[\textbf{Matthew}] \hyperref[MATTHEW-18-21]{18:21--35}
				\item[\textbf{Mark}] ---
				\item[\textbf{Luke}] ---
				\item[\textbf{John}] ---
			\end{compactitem}
			
		% 175
		\item The Samaritan Villages
			\begin{compactitem}
				\item[\textbf{Matthew}] ---
				\item[\textbf{Mark}] ---
				\item[\textbf{Luke}] \hyperref[LUKE-9-51]{9:51--56}
				\item[\textbf{John}] \hyperref[JOHN-4-1]{4:1--42}
			\end{compactitem}
			
		% 176
		\item Mission of the Seventy
			\begin{compactitem}
				\item[\textbf{Matthew}] ---
				\item[\textbf{Mark}] ---
				\item[\textbf{Luke}] \hyperref[LUKE-10-1]{10:1--20}
				\item[\textbf{John}] ---
			\end{compactitem}
			
		% 177
		\item Rejoicing of Jesus
			\begin{compactitem}
				\item[\textbf{Matthew}] ---
				\item[\textbf{Mark}] ---
				\item[\textbf{Luke}] \hyperref[LUKE-10-21]{10:21--24}
				\item[\textbf{John}] ---
			\end{compactitem}
			
		% 178
		\item The Good Samaritan
			\begin{compactitem}
				\item[\textbf{Matthew}] ---
				\item[\textbf{Mark}] ---
				\item[\textbf{Luke}] \hyperref[LUKE-10-25]{10:25--37}
				\item[\textbf{John}] ---
			\end{compactitem}
			
		% 179
		\item Visiting Mary and Martha
			\begin{compactitem}
				\item[\textbf{Matthew}] ---
				\item[\textbf{Mark}] ---
				\item[\textbf{Luke}] \hyperref[LUKE-10-38]{10:38--42}
				\item[\textbf{John}] ---
			\end{compactitem}
			
		% 180
		\item Blessedness of Christ's Mother
			\begin{compactitem}
				\item[\textbf{Matthew}] ---
				\item[\textbf{Mark}] ---
				\item[\textbf{Luke}] \hyperref[LUKE-11-27]{11:27--28}
				\item[\textbf{John}] ---
			\end{compactitem}
	
	\columnbreak
	
		% 181
		\item Warning against Hypocrisy
			\begin{compactitem}
				\item[\textbf{Matthew}] ---
				\item[\textbf{Mark}] ---
				\item[\textbf{Luke}] \hyperref[LUKE-12-1]{12:1--3}
				\item[\textbf{John}] ---
			\end{compactitem}
			
		% 182
		\item Parable of the Rich Fool
			\begin{compactitem}
				\item[\textbf{Matthew}] ---
				\item[\textbf{Mark}] ---
				\item[\textbf{Luke}] \hyperref[LUKE-12-13]{12:13--21}
				\item[\textbf{John}] ---
			\end{compactitem}
			
		% 183
		\item Repent or Perish
			\begin{compactitem}
				\item[\textbf{Matthew}] ---
				\item[\textbf{Mark}] ---
				\item[\textbf{Luke}] \hyperref[LUKE-13-1]{13:1--5}
				\item[\textbf{John}] ---
			\end{compactitem}
			
		% 184
		\item Parable of the Barren Fig Tree
			\begin{compactitem}
				\item[\textbf{Matthew}] ---
				\item[\textbf{Mark}] ---
				\item[\textbf{Luke}] \hyperref[LUKE-13-6]{13:6--9}
				\item[\textbf{John}] ---
			\end{compactitem}
			
		% 185
		\item Healing a Crippled Woman on the Sabbath
			\begin{compactitem}
				\item[\textbf{Matthew}] ---
				\item[\textbf{Mark}] ---
				\item[\textbf{Luke}] \hyperref[LUKE-13-10]{13:10--17}
				\item[\textbf{John}] ---
			\end{compactitem}
			
		% 186
		\item Healing a Man with Dropsy
			\begin{compactitem}
				\item[\textbf{Matthew}] ---
				\item[\textbf{Mark}] ---
				\item[\textbf{Luke}] \hyperref[LUKE-14-1]{14:1--4}
				\item[\textbf{John}] ---
			\end{compactitem}
			
		% 187
		\item A Lesson to Guests and Hosts
			\begin{compactitem}
				\item[\textbf{Matthew}] ---
				\item[\textbf{Mark}] ---
				\item[\textbf{Luke}] \hyperref[LUKE-14-7]{14:7--14}
				\item[\textbf{John}] ---
			\end{compactitem}
			
		% 188
		\item Parable of the Lost Coin
			\begin{compactitem}
				\item[\textbf{Matthew}] ---
				\item[\textbf{Mark}] ---
				\item[\textbf{Luke}] \hyperref[LUKE-15-8]{15:8--10}
				\item[\textbf{John}] ---
			\end{compactitem}
			
		% 189
		\item Parable of the Prodigal Son
			\begin{compactitem}
				\item[\textbf{Matthew}] ---
				\item[\textbf{Mark}] ---
				\item[\textbf{Luke}] \hyperref[LUKE-15-11]{15:11--32}
				\item[\textbf{John}] ---
			\end{compactitem}
			
		% 190
		\item Parable of the Dishonest Steward
			\begin{compactitem}
				\item[\textbf{Matthew}] ---
				\item[\textbf{Mark}] ---
				\item[\textbf{Luke}] \hyperref[LUKE-16-1]{16:1--12}
				\item[\textbf{John}] ---
			\end{compactitem}
	
	\columnbreak
	
		% 191
		\item Lazarus and the Rich Man
			\begin{compactitem}
				\item[\textbf{Matthew}] ---
				\item[\textbf{Mark}] ---
				\item[\textbf{Luke}] \hyperref[LUKE-16-19]{16:19--31}
				\item[\textbf{John}] ---
			\end{compactitem}
			
		% 192
		\item Unprofitable Servants
			\begin{compactitem}
				\item[\textbf{Matthew}] ---
				\item[\textbf{Mark}] ---
				\item[\textbf{Luke}] \hyperref[LUKE-17-5]{17:5--10}
				\item[\textbf{John}] ---
			\end{compactitem}
			
		% 193
		\item Cleansing of Ten Lepers
			\begin{compactitem}
				\item[\textbf{Matthew}] ---
				\item[\textbf{Mark}] ---
				\item[\textbf{Luke}] \hyperref[LUKE-17-11]{17:11--19}
				\item[\textbf{John}] ---
			\end{compactitem}
			
		% 194
		\item Parable of the Widow and Judge
			\begin{compactitem}
				\item[\textbf{Matthew}] ---
				\item[\textbf{Mark}] ---
				\item[\textbf{Luke}] \hyperref[LUKE-18-1]{18:1--8}
				\item[\textbf{John}] ---
			\end{compactitem}
			
		% 195
		\item Parable of the Pharisee and the Publican
			\begin{compactitem}
				\item[\textbf{Matthew}] ---
				\item[\textbf{Mark}] ---
				\item[\textbf{Luke}] \hyperref[LUKE-18-9]{18:9--14}
				\item[\textbf{John}] ---
			\end{compactitem}
			
		% 196
		\item Marriage and Divorce
			\begin{compactitem}
				\item[\textbf{Matthew}] \hyperref[MATTHEW-19-1]{19:1--10}
				\item[\textbf{Mark}] \hyperref[MARK-10-2]{10:2--12}
				\item[\textbf{Luke}] ---
				\item[\textbf{John}] ---
			\end{compactitem}
			
		% 197
		\item \textit{Suffer the Little Children}
			\begin{compactitem}
				\item[\textbf{Matthew}] \hyperref[MATTHEW-19-13]{19:13--15}
				\item[\textbf{Mark}] \hyperref[MARK-10-13]{10:13--16}
				\item[\textbf{Luke}] \hyperref[LUKE-18-15]{18:15--17}
				\item[\textbf{John}] ---
			\end{compactitem}
			
		% 198
		\item \textit{The Rich Young Ruler}
			\begin{compactitem}
				\item[\textbf{Matthew}] \hyperref[MATTHEW-19-16]{19:16--30}
				\item[\textbf{Mark}] \hyperref[MARK-10-17]{10:17--31}
				\item[\textbf{Luke}] \hyperref[LUKE-18-18]{18:18--30}
				\item[\textbf{John}] ---
			\end{compactitem}
			
		% 199
		\item Parable of Laborers
			\begin{compactitem}
				\item[\textbf{Matthew}] \hyperref[MATTHEW-20-1]{20:1--16}
				\item[\textbf{Mark}] ---
				\item[\textbf{Luke}] ---
				\item[\textbf{John}] ---
			\end{compactitem}
			
		% 200
		\item \textit{The Third Prophecy of the Passion}
			\begin{compactitem}
				\item[\textbf{Matthew}] \hyperref[MATTHEW-20-17]{20:17--19}
				\item[\textbf{Mark}] \hyperref[MARK-10-32]{10:32--34}
				\item[\textbf{Luke}] \hyperref[LUKE-18-31]{18:31--34}
				\item[\textbf{John}] ---
			\end{compactitem}
	
	\columnbreak
	
		% 201
		\item Death of Lazarus
			\begin{compactitem}
				\item[\textbf{Matthew}] ---
				\item[\textbf{Mark}] ---
				\item[\textbf{Luke}] ---
				\item[\textbf{John}] \hyperref[JOHN-11-1]{11:1--17}
			\end{compactitem}
			
		% 202
		\item Resurrection and the Life
			\begin{compactitem}
				\item[\textbf{Matthew}] ---
				\item[\textbf{Mark}] ---
				\item[\textbf{Luke}] ---
				\item[\textbf{John}] \hyperref[JOHN-11-18]{11:18--27}
			\end{compactitem}
			
		% 203
		\item Jesus Weeps 
			\begin{compactitem}
				\item[\textbf{Matthew}] ---
				\item[\textbf{Mark}] ---
				\item[\textbf{Luke}] ---
				\item[\textbf{John}] \hyperref[JOHN-11-28]{11:28--37}
			\end{compactitem}
			
		% 204
		\item Lazarus Raised
			\begin{compactitem}
				\item[\textbf{Matthew}] ---
				\item[\textbf{Mark}] ---
				\item[\textbf{Luke}] ---
				\item[\textbf{John}] \hyperref[JOHN-11-38]{11:38--46}
			\end{compactitem}
			
		% 205
		\item James and John
			\begin{compactitem}
				\item[\textbf{Matthew}] \hyperref[MATTHEW-20-20]{20:20--28}
				\item[\textbf{Mark}] \hyperref[MARK-10-35]{10:35--45}
				\item[\textbf{Luke}] ---
				\item[\textbf{John}] ---
			\end{compactitem}
			
		% 206
		\item \textit{Healing Bartimaeus}
			\begin{compactitem}
				\item[\textbf{Matthew}] \hyperref[MATTHEW-20-29]{20:29--34}
				\item[\textbf{Mark}] \hyperref[MARK-10-46]{10:46--52}
				\item[\textbf{Luke}] \hyperref[LUKE-18-35]{18:35--43}
				\item[\textbf{John}] ---
			\end{compactitem}
			
		% 207
		\item Zacchaeus
			\begin{compactitem}
				\item[\textbf{Matthew}] ---
				\item[\textbf{Mark}] ---
				\item[\textbf{Luke}] \hyperref[LUKE-19-2]{19:2--10}
				\item[\textbf{John}] ---
			\end{compactitem}
			
		% 208
		\item Parable of the Talens
			\begin{compactitem}
				\item[\textbf{Matthew}] \hyperref[MATTHEW-25-14]{25:14--30}
				\item[\textbf{Mark}] ---
				\item[\textbf{Luke}] \hyperref[LUKE-19-11]{19:11--27}
				\item[\textbf{John}] ---
			\end{compactitem}
			
		% 209
		\item Plot against Lazarus
			\begin{compactitem}
				\item[\textbf{Matthew}] ---
				\item[\textbf{Mark}] ---
				\item[\textbf{Luke}] ---
				\item[\textbf{John}] \hyperref[JOHN-12-9]{12:9--11}
			\end{compactitem}
			
		% 210
		\item \textit{The Triumphal Entry}
			\begin{compactitem}
				\item[\textbf{Matthew}] \hyperref[MATTHEW-21-1]{21:1--11}
				\item[\textbf{Mark}] \hyperref[MARK-11-1]{11:1--11}
				\item[\textbf{Luke}] \hyperref[LUKE-19-28]{19:28--38}
				\item[\textbf{John}] \hyperref[JOHN-12-12]{12:12--19}
			\end{compactitem}
	
	\columnbreak
	
		% 211
		\item The Hour Is Come 
			\begin{compactitem}
				\item[\textbf{Matthew}] ---
				\item[\textbf{Mark}] ---
				\item[\textbf{Luke}] ---
				\item[\textbf{John}] \hyperref[JOHN-12-20]{12:20--26}
			\end{compactitem}
			
		% 212
		\item Voice from Heaven
			\begin{compactitem}
				\item[\textbf{Matthew}] ---
				\item[\textbf{Mark}] ---
				\item[\textbf{Luke}] ---
				\item[\textbf{John}] \hyperref[JOHN-12-27]{12:27--30}
			\end{compactitem}
			
		% 213
		\item Christ to Be Lifted Up
			\begin{compactitem}
				\item[\textbf{Matthew}] ---
				\item[\textbf{Mark}] ---
				\item[\textbf{Luke}] ---
				\item[\textbf{John}] \hyperref[JOHN-12-27]{12:27--36}
			\end{compactitem}
			
		% 214
		\item The Unbelief of the Jews
			\begin{compactitem}
				\item[\textbf{Matthew}] ---
				\item[\textbf{Mark}] ---
				\item[\textbf{Luke}] ---
				\item[\textbf{John}] \hyperref[JOHN-12-36]{12:36--43}
			\end{compactitem}
			
		% 215
		\item Judgment by Jesus' Word
			\begin{compactitem}
				\item[\textbf{Matthew}] ---
				\item[\textbf{Mark}] ---
				\item[\textbf{Luke}] ---
				\item[\textbf{John}] \hyperref[JOHN-12-44]{12:44--50}
			\end{compactitem}
			
		% 216
		\item Prediction of Jerusalem's Fall
			\begin{compactitem}
				\item[\textbf{Matthew}] ---
				\item[\textbf{Mark}] ---
				\item[\textbf{Luke}] \hyperref[LUKE-19-39]{19:39--44}
				\item[\textbf{John}] ---
			\end{compactitem}
			
		% 217
		\item Cursing the Fig Tree
			\begin{compactitem}
				\item[\textbf{Matthew}] \hyperref[MATTHEW-21-18]{21:18--22}
				\item[\textbf{Mark}] \hyperref[MARK-11-12]{11:12--14}; \hyperref[MARK-11-20]{20--24}
				\item[\textbf{Luke}] ---
				\item[\textbf{John}] ---
			\end{compactitem}
			
		% 218
		\item \textit{Cleansing the Temple}
			\begin{compactitem}
				\item[\textbf{Matthew}] \hyperref[MATTHEW-21-12]{21:12--17}
				\item[\textbf{Mark}] \hyperref[MARK-11-15]{11:15--19}
				\item[\textbf{Luke}] \hyperref[LUKE-19-45]{19:45--48}
				\item[\textbf{John}] \hyperref[JOHN-2-13]{2:13--22}
			\end{compactitem}
			
		% 219
		\item \textit{Lesson of the Fig Tree}
			\begin{compactitem}
				\item[\textbf{Matthew}] \hyperref[MATTHEW-21-20]{21:20--22}; \hyperref[MATTHEW-24-32]{24:32--35}
				\item[\textbf{Mark}] \hyperref[MARK-11-20]{11:20--26}; \hyperref[MARK-13-28]{13:28--31}
				\item[\textbf{Luke}] \hyperref[LUKE-21-29]{21:29--33}
				\item[\textbf{John}] ---
			\end{compactitem}
			
		% 220
		\item \textit{The Question about Authority}
			\begin{compactitem}
				\item[\textbf{Matthew}] \hyperref[MATTHEW-21-23]{21:23--37}
				\item[\textbf{Mark}] \hyperref[MARK-11-27]{11:27--33}
				\item[\textbf{Luke}] \hyperref[LUKE-20-1]{20:1--8}
				\item[\textbf{John}] ---
			\end{compactitem}
	
	\columnbreak
	
		% 221
		\item Parable of the Two Sons 
			\begin{compactitem}
				\item[\textbf{Matthew}] \hyperref[MATTHEW-21-28]{21:28--32}
				\item[\textbf{Mark}] ---
				\item[\textbf{Luke}] ---
				\item[\textbf{John}] ---
			\end{compactitem}
			
		% 222
		\item \textit{Parable of the Wicked Husbandman}
			\begin{compactitem}
				\item[\textbf{Matthew}] \hyperref[MATTHEW-21-33]{21:33--36}
				\item[\textbf{Mark}] \hyperref[MARK-12-1]{12:1--12}
				\item[\textbf{Luke}] \hyperref[LUKE-20-9]{20:9--19}
				\item[\textbf{John}] ---
			\end{compactitem}
			
		% 223
		\item Parable of the Marriage Feast 
			\begin{compactitem}
				\item[\textbf{Matthew}] \hyperref[MATTHEW-22-1]{22:1--14}
				\item[\textbf{Mark}] ---
				\item[\textbf{Luke}] \hyperref[LUKE-14-16]{14:16--24}
				\item[\textbf{John}] ---
			\end{compactitem}
			
		% 224
		\item \textit{Tribute to Caesar}
			\begin{compactitem}
				\item[\textbf{Matthew}] \hyperref[MATTHEW-22-15]{22:15--22}
				\item[\textbf{Mark}] \hyperref[MARK-12-13]{12:13--17}
				\item[\textbf{Luke}] \hyperref[LUKE-20-20]{20:20--26}
				\item[\textbf{John}] ---
			\end{compactitem}
			
		% 225
		\item \textit{Question on Resurrection} 
			\begin{compactitem}
				\item[\textbf{Matthew}] \hyperref[MATTHEW-22-23]{22:23--33}
				\item[\textbf{Mark}] \hyperref[MARK-12-18]{12:18--27}
				\item[\textbf{Luke}] \hyperref[LUKE-20-27]{20:27--40}
				\item[\textbf{John}] ---
			\end{compactitem}
			
		% 226
		\item \textit{The Great Commandment}
			\begin{compactitem}
				\item[\textbf{Matthew}] \hyperref[MATTHEW-22-34]{22:34--40}
				\item[\textbf{Mark}] \hyperref[MARK-12-28]{12:28--34}
				\item[\textbf{Luke}] \hyperref[LUKE-10-25]{10:25--28}
				\item[\textbf{John}] ---
			\end{compactitem}
			
		% 227
		\item Washing the Disciples' Feet 
			\begin{compactitem}
				\item[\textbf{Matthew}] ---
				\item[\textbf{Mark}] ---
				\item[\textbf{Luke}] ---
				\item[\textbf{John}] \hyperref[JOHN-13-1]{13:1--20}
			\end{compactitem}
			
		% 228
		\item \textit{Jesus Foretells His Betrayal}
			\begin{compactitem}
				\item[\textbf{Matthew}] \hyperref[MATTHEW-26-20]{26:20--25}
				\item[\textbf{Mark}] \hyperref[MARK-14-17]{14:17--21}
				\item[\textbf{Luke}] \hyperref[LUKE-22-21]{22:21--23}
				\item[\textbf{John}] \hyperref[JOHN-13-21]{13:21--30}
			\end{compactitem}
			
		% 229
		\item The New Commandment
			\begin{compactitem}
				\item[\textbf{Matthew}] ---
				\item[\textbf{Mark}] ---
				\item[\textbf{Luke}] ---
				\item[\textbf{John}] \hyperref[JOHN-13-31]{13:31--35}
			\end{compactitem}
			
		% 230
		\item \textit{Aboud David's Son} 
			\begin{compactitem}
				\item[\textbf{Matthew}] \hyperref[MATTHEW-22-41]{22:41--46}
				\item[\textbf{Mark}] \hyperref[MARK-12-35]{12:35--37}
				\item[\textbf{Luke}] \hyperref[LUKE-20-41]{20:41--44}
				\item[\textbf{John}] ---
			\end{compactitem}
	
	\columnbreak
	
		% 231
		\item The Way to the Father
			\begin{compactitem}
				\item[\textbf{Matthew}] ---
				\item[\textbf{Mark}] ---
				\item[\textbf{Luke}] ---
				\item[\textbf{John}] \hyperref[JOHN-14-1]{14:1--14}
			\end{compactitem}
			
		% 232
		\item Promise of the Spirit 
			\begin{compactitem}
				\item[\textbf{Matthew}] ---
				\item[\textbf{Mark}] ---
				\item[\textbf{Luke}] ---
				\item[\textbf{John}] \hyperref[JOHN-14-14]{14:14--31}
			\end{compactitem}
			
		% 233
		\item The True Vine 
			\begin{compactitem}
				\item[\textbf{Matthew}] ---
				\item[\textbf{Mark}] ---
				\item[\textbf{Luke}] ---
				\item[\textbf{John}] \hyperref[JOHN-15-1]{15:1--17}
			\end{compactitem}
			
		% 234
		\item The World's Hatred 
			\begin{compactitem}
				\item[\textbf{Matthew}] ---
				\item[\textbf{Mark}] ---
				\item[\textbf{Luke}] ---
				\item[\textbf{John}] \hyperref[JOHN-15-18]{15:18--16:4}
			\end{compactitem}
			
		% 235
		\item \textit{Woes against the Scribes and Pharisees}
			\begin{compactitem}
				\item[\textbf{Matthew}] \hyperref[MATTHEW-23-1]{23:1--36}
				\item[\textbf{Mark}] \hyperref[MARK-12-38]{12:38--40}
				\item[\textbf{Luke}] \hyperref[LUKE-11-39]{11:39--52}; \hyperref[LUKE-20-45]{20:45--47}
				\item[\textbf{John}] ---
			\end{compactitem}
			
		% 236
		\item The Lament over Jerusalem
			\begin{compactitem}
				\item[\textbf{Matthew}] \hyperref[MATTHEW-23-37]{23:37--39}
				\item[\textbf{Mark}] ---
				\item[\textbf{Luke}] \hyperref[LUKE-13-34]{13:34--35}
				\item[\textbf{John}] ---
			\end{compactitem}
		
		% 237
		\item The Widow's Mite 
			\begin{compactitem}
				\item[\textbf{Matthew}] ---
				\item[\textbf{Mark}] \hyperref[MARK-12-41]{12:41--44}
				\item[\textbf{Luke}] \hyperref[LUKE-21-1]{21:1--4}
				\item[\textbf{John}] ---
			\end{compactitem}
			
		% 238
		\item Work of the Spirit
			\begin{compactitem}
				\item[\textbf{Matthew}] ---
				\item[\textbf{Mark}] ---
				\item[\textbf{Luke}] ---
				\item[\textbf{John}] \hyperref[JOHN-16-7]{16:7--15}
			\end{compactitem}
			
		% 239
		\item \textit{Prophecy of the Temple's Destruction}
			\begin{compactitem}
				\item[\textbf{Matthew}] \hyperref[MATTHEW-24-1]{24:1--2}
				\item[\textbf{Mark}] \hyperref[MARK-13-1]{13:1--2}
				\item[\textbf{Luke}] \hyperref[LUKE-21-5]{21:5--6}; \hyperref[LUKE-21-20]{20--24}
				\item[\textbf{John}] ---
			\end{compactitem}
			
		% 240
		\item Sorrow Will Turn to Joy 
			\begin{compactitem}
				\item[\textbf{Matthew}] ---
				\item[\textbf{Mark}] ---
				\item[\textbf{Luke}] ---
				\item[\textbf{John}] \hyperref[JOHN-16-16]{16:16--24}
			\end{compactitem}
	
	\columnbreak
	
		% 241
		\item I Have Overcome the World
			\begin{compactitem}
				\item[\textbf{Matthew}] ---
				\item[\textbf{Mark}] ---
				\item[\textbf{Luke}] ---
				\item[\textbf{John}] \hyperref[JOHN-16-25]{16:25--33}
			\end{compactitem}
			
		% 242
		\item \textit{The Last Days}
			\begin{compactitem}
				\item[\textbf{Matthew}] \hyperref[MATTHEW-24-3]{24:3--14}
				\item[\textbf{Mark}] \hyperref[MARK-13-3]{13:3--31}
				\item[\textbf{Luke}] \hyperref[LUKE-21-7]{21:7--19}
				\item[\textbf{John}] ---
			\end{compactitem}
			
		% 243
		\item \textit{The Need for Watchfulness}
			\begin{compactitem}
				\item[\textbf{Matthew}] \hyperref[MATTHEW-24-15]{24:15--42}
				\item[\textbf{Mark}] \hyperref[MARK-13-32]{13:32--37}
				\item[\textbf{Luke}] \hyperref[LUKE-21-34]{21:34--38}
				\item[\textbf{John}] ---
			\end{compactitem}
			
		% 244
		\item The Watchful Householder
			\begin{compactitem}
				\item[\textbf{Matthew}] \hyperref[MATTHEW-24-43]{24:43--44}
				\item[\textbf{Mark}] ---
				\item[\textbf{Luke}] \hyperref[LUKE-12-35]{12:35--40}
				\item[\textbf{John}] ---
			\end{compactitem}
			
		% 245
		\item Coming of the Kingdom 
			\begin{compactitem}
				\item[\textbf{Matthew}] \hyperref[MATTHEW-24-23]{24:23--28}; \hyperref[MATTHEW-24-37]{37--41}
				\item[\textbf{Mark}] ---
				\item[\textbf{Luke}] \hyperref[LUKE-17-20]{17:20--37}
				\item[\textbf{John}] ---
			\end{compactitem}
			
		% 246
		\item \textit{The Coming of the Son of Man} 
			\begin{compactitem}
				\item[\textbf{Matthew}] \hyperref[MATTHEW-24-29]{24:29--31}
				\item[\textbf{Mark}] \hyperref[MARK-13-24]{13:24--27}
				\item[\textbf{Luke}] \hyperref[LUKE-21-25]{21:25--28}
				\item[\textbf{John}] ---
			\end{compactitem}
			
		% 247
		\item The Faithful Servant
			\begin{compactitem}
				\item[\textbf{Matthew}] \hyperref[MATTHEW-24-45]{24:45--51}
				\item[\textbf{Mark}] ---
				\item[\textbf{Luke}] \hyperref[LUKE-12-41]{12:41--48}
				\item[\textbf{John}] ---
			\end{compactitem}
			
		% 248
		\item The Ten Virgins 
			\begin{compactitem}
				\item[\textbf{Matthew}] \hyperref[MATTHEW-25-1]{25:1--13}
				\item[\textbf{Mark}] ---
				\item[\textbf{Luke}] ---
				\item[\textbf{John}] ---
			\end{compactitem}
			
		% 249
		\item The Last Judgment
			\begin{compactitem}
				\item[\textbf{Matthew}] \hyperref[MATTHEW-25-31]{25:31--46}
				\item[\textbf{Mark}] ---
				\item[\textbf{Luke}] ---
				\item[\textbf{John}] ---
			\end{compactitem}
			
		% 250
		\item \textit{Plot to Kill Jesus}
			\begin{compactitem}
				\item[\textbf{Matthew}] \hyperref[MATTHEW-26-1]{26:1--5}, \hyperref[MATTHEW-26-14]{14--16}
				\item[\textbf{Mark}] \hyperref[MARK-14-1]{14:1--2}, \hyperref[MARK-14-10]{10--11}
				\item[\textbf{Luke}] \hyperref[LUKE-22-1]{22:1--2}
				\item[\textbf{John}] \hyperref[JOHN-11-46]{11:46--57}
			\end{compactitem}
	
	\columnbreak
	
		% 251
		\item Anointing at Bethany 
			\begin{compactitem}
				\item[\textbf{Matthew}] \hyperref[MATTHEW-26-1]{26:1--13}
				\item[\textbf{Mark}] \hyperref[MARK-14-3]{14:3--9}
				\item[\textbf{Luke}] ---
				\item[\textbf{John}] \hyperref[JOHN-12-1]{12:1--8}
			\end{compactitem}
			
		% 252
		\item \textit{Betrayal by Judas}
			\begin{compactitem}
				\item[\textbf{Matthew}] \hyperref[MATTHEW-26-14]{26:14--16}
				\item[\textbf{Mark}] \hyperref[MARK-14-10]{14:10--11}
				\item[\textbf{Luke}] \hyperref[LUKE-22-3]{22:3--6}
				\item[\textbf{John}] \hyperref[JOHN-18-1]{18:1--3}
			\end{compactitem}
			
		% 253
		\item \textit{Preparations for Passover} 
			\begin{compactitem}
				\item[\textbf{Matthew}] \hyperref[MATTHEW-26-17]{26:17--25}
				\item[\textbf{Mark}] \hyperref[MARK-14-12]{14:12--21}
				\item[\textbf{Luke}] \hyperref[LUKE-22-7]{22:7--13}
				\item[\textbf{John}] \hyperref[JOHN-13-21]{13:21--30}
			\end{compactitem}
			
		% 254
		\item \textit{The Last Supper}
			\begin{compactitem}
				\item[\textbf{Matthew}] \hyperref[MATTHEW-26-17]{26:17--30}
				\item[\textbf{Mark}] \hyperref[MARK-14-12]{14:12--26}
				\item[\textbf{Luke}] \hyperref[LUKE-22-7]{22:7--23}
				\item[\textbf{John}] \hyperref[JOHN-13-21]{13:21--30}
			\end{compactitem}
			
		% 255
		\item The Great Prayer 
			\begin{compactitem}
				\item[\textbf{Matthew}] ---
				\item[\textbf{Mark}] ---
				\item[\textbf{Luke}] ---
				\item[\textbf{John}] \hyperref[JOHN-17-1]{17:1--26}
			\end{compactitem}
			
		% 256
		\item \textit{Peter's Denial Foretold}
			\begin{compactitem}
				\item[\textbf{Matthew}] \hyperref[MATTHEW-26-31]{26:31--35}
				\item[\textbf{Mark}] \hyperref[MARK-14-27]{14:27--31}
				\item[\textbf{Luke}] \hyperref[LUKE-22-31]{22:31--34}
				\item[\textbf{John}] \hyperref[JOHN-13-36]{13:36--38}
			\end{compactitem}
			
		% 257
		\item Purse, Scrip, and Sword 
			\begin{compactitem}
				\item[\textbf{Matthew}] ---
				\item[\textbf{Mark}] ---
				\item[\textbf{Luke}] \hyperref[LUKE-22-35]{22:35--38}
				\item[\textbf{John}] ---
			\end{compactitem}
			
		% 258
		\item \textit{To Gethsemane}
			\begin{compactitem}
				\item[\textbf{Matthew}] \hyperref[MATTHEW-26-30]{26:30}
				\item[\textbf{Mark}] \hyperref[MARK-14-26]{14:26}
				\item[\textbf{Luke}] \hyperref[LUKE-22-39]{22:39}
				\item[\textbf{John}] \hyperref[JOHN-18-1]{18:1}
			\end{compactitem}
			
		% 259
		\item \textit{In Gethsemane} 
			\begin{compactitem}
				\item[\textbf{Matthew}] \hyperref[MATTHEW-26-36]{26:36--46}
				\item[\textbf{Mark}] \hyperref[MARK-14-32]{14:32--42}
				\item[\textbf{Luke}] \hyperref[LUKE-22-39]{22:39--46}
				\item[\textbf{John}] ---
			\end{compactitem}
			
		% 260
		\item \textit{Jesus Arrested} 
			\begin{compactitem}
				\item[\textbf{Matthew}] \hyperref[MATTHEW-26-47]{26:47--56}
				\item[\textbf{Mark}] \hyperref[LUKE-14-43]{14:43--50}
				\item[\textbf{Luke}] \hyperref[LUKE-22-47]{22:47--53}
				\item[\textbf{John}] \hyperref[JOHN-18-3]{18:3--12}
			\end{compactitem}
	
	\columnbreak
	
		% 261
		\item The Naked Young Man
			\begin{compactitem}
				\item[\textbf{Matthew}] ---
				\item[\textbf{Mark}] \hyperref[MARK-14-51]{14:51--52}
				\item[\textbf{Luke}] ---
				\item[\textbf{John}] ---
			\end{compactitem}
			
		% 262
		\item \textit{Jesus before the Sanhedrin}
			\begin{compactitem}
				\item[\textbf{Matthew}] \hyperref[MATTHEW-26-57]{26:57--68}
				\item[\textbf{Mark}] \hyperref[MARK-14-53]{14:53--65}
				\item[\textbf{Luke}] \hyperref[LUKE-22-54]{22:54--55}, \hyperref[LUKE-22-63]{63--71}
				\item[\textbf{John}] \hyperref[JOHN-18-13]{18:13--14}, \hyperref[JOHN-18-19]{19--24}
			\end{compactitem}
			
		% 263
		\item \textit{Peter's Denial of Christ} 
			\begin{compactitem}
				\item[\textbf{Matthew}] \hyperref[MATTHEW-26-69]{26:69--75}
				\item[\textbf{Mark}] \hyperref[MARK-14-66]{14:66--72}
				\item[\textbf{Luke}] \hyperref[LUKE-22-56]{22:56--62}
				\item[\textbf{John}] \hyperref[JOHN-18-15]{18:15--18}, \hyperref[JOHN-18-25]{25--27}
			\end{compactitem}
			
		% 264
		\item Jesus before the High Priest
			\begin{compactitem}
				\item[\textbf{Matthew}] ---
				\item[\textbf{Mark}] ---
				\item[\textbf{Luke}] ---
				\item[\textbf{John}] \hyperref[JOHN-18-19]{18:19--24}
			\end{compactitem}
			
		% 265
		\item Peter's Second and Third Denial
			\begin{compactitem}
				\item[\textbf{Matthew}] ---
				\item[\textbf{Mark}] ---
				\item[\textbf{Luke}] ---
				\item[\textbf{John}] \hyperref[JOHN-18-25]{18:25--27}
			\end{compactitem}
			
		% 266
		\item \textit{Mockery and Beating of Jesus}
			\begin{compactitem}
				\item[\textbf{Matthew}] \hyperref[MATTHEW-26-67]{26:67--68}
				\item[\textbf{Mark}] \hyperref[MARK-14-65]{14:65}
				\item[\textbf{Luke}] \hyperref[LUKE-22-63]{22:63--65}
				\item[\textbf{John}] ---
			\end{compactitem}
			
		% 267
		\item \textit{Christ Delivered to Pilate} 
			\begin{compactitem}
				\item[\textbf{Matthew}] \hyperref[MATTHEW-27-1]{27:1--2}
				\item[\textbf{Mark}] \hyperref[MARK-15-1]{15:1}
				\item[\textbf{Luke}] \hyperref[LUKE-23-1]{23:1--2}
				\item[\textbf{John}] \hyperref[JOHN-18-28]{18:28--32}
			\end{compactitem}
			
		% 268
		\item Death of Judas
			\begin{compactitem}
				\item[\textbf{Matthew}] \hyperref[MATTHEW-27-3]{27:3--10}
				\item[\textbf{Mark}] ---
				\item[\textbf{Luke}] ---
				\item[\textbf{John}] ---
			\end{compactitem}
			
		% 269
		\item \textit{Trial before Pilate}
			\begin{compactitem}
				\item[\textbf{Matthew}] \hyperref[MATTHEW-27-11]{27:11--14}
				\item[\textbf{Mark}] \hyperref[MARK-15-2]{15:2--5}
				\item[\textbf{Luke}] \hyperref[LUKE-23-3]{23:3--5}
				\item[\textbf{John}] \hyperref[JOHN-18-33]{18:33--38}
			\end{compactitem}
			
		% 270
		\item Christ before Herod
			\begin{compactitem}
				\item[\textbf{Matthew}] ---
				\item[\textbf{Mark}] ---
				\item[\textbf{Luke}] \hyperref[LUKE-23-6]{23:6--12}
				\item[\textbf{John}] ---
			\end{compactitem}
	
	\columnbreak
	
		% 271
		\item \textit{Jesus Sentenced to Die} 
			\begin{compactitem}
				\item[\textbf{Matthew}] \hyperref[MATTHEW-27-15]{27:15--26}
				\item[\textbf{Mark}] \hyperref[MARK-15-6]{15:6--15}
				\item[\textbf{Luke}] \hyperref[LUKE-23-13]{23:13--25}
				\item[\textbf{John}] \hyperref[JOHN-18-39]{18:399--19:16}
			\end{compactitem}
			
		% 272
		\item Mocking by Soldiers
			\begin{compactitem}
				\item[\textbf{Matthew}] \hyperref[MATTHEW-27-27]{27:27--31}
				\item[\textbf{Mark}] \hyperref[MARK-15-16]{15:16--20}
				\item[\textbf{Luke}] ---
				\item[\textbf{John}] \hyperref[JOHN-19-2]{19:2--3}
			\end{compactitem}
	
		% 273
		\item \textit{Road to Calvary}
			\begin{compactitem}
				\item[\textbf{Matthew}] \hyperref[MATTHEW-27-32]{27:32}
				\item[\textbf{Mark}] \hyperref[MARK-15-21]{15:21}
				\item[\textbf{Luke}] \hyperref[LUKE-23-26]{23:26--32}
				\item[\textbf{John}] ---
			\end{compactitem}
			
		% 274
		\item \textit{The Crucifixion}
			\begin{compactitem}
				\item[\textbf{Matthew}] \hyperref[MATTHEW-27-32]{27:32--44}
				\item[\textbf{Mark}] \hyperref[MARK-15-21]{15:21--32}
				\item[\textbf{Luke}] \hyperref[LUKE-23-26]{23:26--43}
				\item[\textbf{John}] \hyperref[JOHN-19-16]{19:16--27}
			\end{compactitem}
			
		% 275
		\item \textit{Death on the Cross}
			\begin{compactitem}
				\item[\textbf{Matthew}] \hyperref[MATTHEW-27-45]{27:45--56}
				\item[\textbf{Mark}] \hyperref[MARK-15-33]{15:33--41}
				\item[\textbf{Luke}] \hyperref[LUKE-23-44]{23:44--49}
				\item[\textbf{John}] \hyperref[JOHN-19-28]{19:28--30}
			\end{compactitem}
			
		% 276
		\item Piercing Jesus' Side
			\begin{compactitem}
				\item[\textbf{Matthew}] ---
				\item[\textbf{Mark}] ---
				\item[\textbf{Luke}] ---
				\item[\textbf{John}] \hyperref[JOHN-19-31]{19:31--37}
			\end{compactitem}
			
		% 277
		\item \textit{Burial of Jesus}
			\begin{compactitem}
				\item[\textbf{Matthew}] \hyperref[MATTHEW-27-57]{27:57--61}
				\item[\textbf{Mark}] \hyperref[MARK-15-42]{15:42--47}
				\item[\textbf{Luke}] \hyperref[LUKE-23-50]{23:50--56}
				\item[\textbf{John}] \hyperref[JOHN-19-38]{19:38--42}
			\end{compactitem}
			
		% 278
		\item Guard at the Tomb
			\begin{compactitem}
				\item[\textbf{Matthew}] \hyperref[MATTHEW-27-62]{27:62--66}
				\item[\textbf{Mark}] ---
				\item[\textbf{Luke}] ---
				\item[\textbf{John}] ---
			\end{compactitem}
			
		% 279
		\item \textit{The Resurrection}
			\begin{compactitem}
				\item[\textbf{Matthew}] \hyperref[MATTHEW-28-1]{28:1--10}
				\item[\textbf{Mark}] \hyperref[MARK-16-1]{16:1--8}
				\item[\textbf{Luke}] \hyperref[LUKE-24-1]{24:1--12}
				\item[\textbf{John}] \hyperref[JOHN-20:1]{20:1--10}
			\end{compactitem}
			
		% 280
		\item Jesus Appears to Mary
			\begin{compactitem}
				\item[\textbf{Matthew}] \hyperref[MATTHEW-28-9]{28:9--10}
				\item[\textbf{Mark}] \hyperref[MARK-16-9]{16:9--11}
				\item[\textbf{Luke}] ---
				\item[\textbf{John}] \hyperref[JOHN-20-11]{20:11--18}
			\end{compactitem}
	
	\columnbreak
	
		% 281
		\item The Bribed Guard's Report
			\begin{compactitem}
				\item[\textbf{Matthew}] \hyperref[MATTHEW-28-11]{28:11--15}
				\item[\textbf{Mark}] ---
				\item[\textbf{Luke}] ---
				\item[\textbf{John}] ---
			\end{compactitem}
			
		% 282
		\item Appearance on the Road to Emmaus
			\begin{compactitem}
				\item[\textbf{Matthew}] ---
				\item[\textbf{Mark}] \hyperref[MARK-16-12]{16:12--13}
				\item[\textbf{Luke}] \hyperref[LUKE-24-13]{24:13--35}
				\item[\textbf{John}] ---
			\end{compactitem}
			
		% 283
		\item \textit{Appearence to Ten Disciples}
			\begin{compactitem}
				\item[\textbf{Matthew}] \hyperref[MATTHEW-28-16]{28:16--20}
				\item[\textbf{Mark}] \hyperref[MARK-16-14]{16:14--18}
				\item[\textbf{Luke}] \hyperref[LUKE-24-36]{24:36--49}
				\item[\textbf{John}] \hyperref[JOHN-20-19]{20:19--23}
			\end{compactitem}
			
		% 284
		\item Jesus and Thomas
			\begin{compactitem}
				\item[\textbf{Matthew}] ---
				\item[\textbf{Mark}] ---
				\item[\textbf{Luke}] ---
				\item[\textbf{John}] \hyperref[JOHN-20-24]{20:24--29}
			\end{compactitem}
			
		% 285
		\item Appearence to Eleven Disciples 
			\begin{compactitem}
				\item[\textbf{Matthew}] ---
				\item[\textbf{Mark}] ---
				\item[\textbf{Luke}] ---
				\item[\textbf{John}] \hyperref[JOHN-20-24]{20:24--26}
			\end{compactitem}
			
		% 286
		\item Appearance to Seven Disciples
			\begin{compactitem}
				\item[\textbf{Matthew}] ---
				\item[\textbf{Mark}] ---
				\item[\textbf{Luke}] ---
				\item[\textbf{John}] \hyperref[JOHN-21-1]{21:1--14}
			\end{compactitem}
			
		% 287
		\item Jesus and Peter
			\begin{compactitem}
				\item[\textbf{Matthew}] ---
				\item[\textbf{Mark}] ---
				\item[\textbf{Luke}] ---
				\item[\textbf{John}] \hyperref[JOHN-21-15]{21:15--19}
			\end{compactitem}
			
		% 288
		\item Jesus and John
			\begin{compactitem}
				\item[\textbf{Matthew}] ---
				\item[\textbf{Mark}] ---
				\item[\textbf{Luke}] ---
				\item[\textbf{John}] \hyperref[JOHN-21-20]{21:20--25}
			\end{compactitem}
			
		% 289
		\item The Ascension of Jesus
			\begin{compactitem}
				\item[\textbf{Matthew}] ---
				\item[\textbf{Mark}] \hyperref[MARK-16-19]{16:19--20}
				\item[\textbf{Luke}] \hyperref[LUKE-24-50]{24:50--53}
				\item[\textbf{John}] ---
			\end{compactitem}
			
		% 290
		\item Purpose of this Book 
			\begin{compactitem}
				\item[\textbf{Matthew}] ---
				\item[\textbf{Mark}] ---
				\item[\textbf{Luke}] ---
				\item[\textbf{John}] \hyperref[JOHN-20-30]{20:30--31}
			\end{compactitem}
	
	\columnbreak
	
		% % 
		% \item 
			% \begin{compactitem}
				% \item[\textbf{Matthew}] ---
				% \item[\textbf{Mark}] ---
				% \item[\textbf{Luke}] ---
				% \item[\textbf{John}] ---
			% \end{compactitem}
		
		% You need this phantom material below in order to get the columns to align right after you compile.
		\phantom{Lorem ipsum dolor sit amet, consectetur adipiscing elit, sed do eiusmod tempor incididunt ut labore et dolore magna aliqua. Ut enim ad minim veniam, quis nostrud exercitation ullamco laboris nisi ut aliquip ex ea commodo consequat. Duis aute irure dolor in reprehenderit in voluptate velit esse cillum dolore eu fugiat nulla pariatur. Excepteur sint occaecat cupidatat non proident, sunt in culpa qui officia deserunt mollit anim id est laborum.}
	
	\end{enumerate}
	
	\end{multicols}
		
% % -----------------------------------------------------------
% \section{38 | The Beatitudes}
% % -----------------------------------------------------------
	% \label{LCHRIST-038}
	
	% % ----------------------
	% \subsection{Matthew 5:3--12}
	% % ----------------------
		% \label{LCHRIST-038-MT}
		
		% \gr{
		% 3 Μακάριοι}%
			% \footnote{%
				% \label{LCHRIST-038-makarioi}%
				% See \hyperref[MAKARIOS]{\grl{μακάριος}} for this word. The two important definitions here are \#1 and \#2, ``pertaining to being fortunate or happy because of circumstances, \textit{fortunate, happy},'' and ``pertaining to being especially favored, \textit{blessed, fortunate, happy, privileged}.'' I think the second definition is probably more operative in these statements, since it is the meaning which involves God taking a more active role in bestowing favor or in causing something to happen in someone's life.
				% }
		% \gr{οἱ πτωχοὶ τῷ πνεύματι,}%
			% \footnote{%
				% \label{LCHRIST-038-ptochoi}%
				% See the \hyperref[PTOCHOS]{definitions} for \gr{πτωχός}, especially definition \#3. After reviewing that material, I would say that I don't think being \gr{πτωχός τῷ πνεύματι} is a very desirable state to be in; I don't know why anyone would \textit{want} to be in a state of dependency on others for their spiritual insight, feeling, and other aspects of their ``human inner life.'' I used to read this verse as giving a reward to those who had some desirable spiritual quality, but I no longer see it that way; I think instead Christ is promising a \textit{recompense} or \textit{compensation} and \textit{amends} to people who were at an unfair disadvantage in some aspects of their spiritual life while on Earth. For me, it also reinforces the importance of \textit{being a person that others can depend on for that spiritual support}---for some people, I will be the one on whom they are dependent for some aspect of their spiritual lives.
				% }
		% \gr{ὅτι αὐτῶν ἐστιν ἡ βασιλεία τῶν οὐρανῶν.
		% }
		
		% \gr{
		% 4 μακάριοι οἱ πενθοῦντες, ὅτι αὐτοὶ παρακληθήσονται.
		% }
		
		% \gr{
		% 5 μακάριοι οἱ πραεῖς, ὅτι αὐτοὶ κληρονομήσουσι τὴν γῆν.
		% }
		
		% \gr{
		% 6 μακάριοι οἱ πεινῶντες καὶ διψῶντες τὴν δικαιοσύνην, ὅτι αὐτοὶ χορτασθήσονται.
		% }
		
		% \gr{
		% 7 μακάριοι οἱ ἐλεήμονες, ὅτι αὐτοὶ ἐλεηθήσονται.
		% }
		
		% \gr{
		% 8 μακάριοι οἱ καθαροὶ τῇ καρδίᾳ, ὅτι αὐτοὶ τὸν θεὸν ὄψονται.
		% }
		
		% \gr{
		% 9 μακάριοι οἱ εἰρηνοποιοί, ὅτι αὐτοὶ υἱοὶ θεοῦ κληθήσονται.
		% }
		
		% \gr{
		% 10 μακάριοι οἱ δεδιωγμένοι ἕνεκεν δικαιοσύνης, ὅτι αὐτῶν ἐστιν ἡ βασιλεία τῶν οὐρανῶν.}%
			% \footnote{%
				% \label{LCHRIST-038-pattern}%
				% If you group these verses in pairs, one from the top and one from the bottom, you see an interesting pattern:
					% \begin{compactdesc}
						% \item[3, 10] Both end with \gr{ὅτι αὐτῶν ἐστιν ἡ βασιλεία τῶν οὐρανῶν}.
						% \item[4, 9] Both finish with a passive verb.
						% \item[5, 8] Both have an active verb with an object. 
						% \item[6, 7] Both finish with a passive verb.
					% \end{compactdesc}
				% No idea whether this pattern is significant in any way, I just thought it was interesting.
				
				% Notice also, as I say below in another note, that the Luke version of these statements all have verbs in the second-person plural, as though Christ is speaking directly to his disciples and not making general statements that would apply to everyone. I guess if we're his disciples, though, they would still apply to us.
				% }
		
		% \gr{
		% 11 μακάριοί ἐστε ὅταν ὀνειδίσωσιν ὑμᾶς καὶ διώξωσιν καὶ εἴπωσιν πᾶν πονηρὸν καθ᾽ ὑμῶν ψευδόμενοι ἕνεκεν ἐμοῦ.
		% 12 χαίρετε καὶ ἀγαλλιᾶσθε, ὅτι ὁ μισθὸς ὑμῶν πολὺς ἐν τοῖς οὐρανοῖς· οὕτως γὰρ ἐδίωξαν τοὺς προφήτας τοὺς πρὸ ὑμῶν.
		% }
	
	% % ----------------------
	% \subsection{Luke 6:20--23}
	% % ----------------------
		% \label{LCHRIST-038-L}
		
		% \gr{
		% 20 Καὶ αὐτὸς ἐπάρας τοὺς ὀφθαλμοὺς αὐτοῦ εἰς τοὺς μαθητὰς αὐτοῦ ἔλεγεν· Μακάριοι οἱ πτωχοί, ὅτι ὑμετέρα ἐστὶν ἡ βασιλεία τοῦ θεοῦ.
		% 21 μακάριοι οἱ πεινῶντες νῦν, ὅτι χορτασθήσεσθε. μακάριοι οἱ κλαίοντες νῦν, ὅτι γελάσετε.}%
			% \footnote{%
				% \label{LCHRIST-038-verbs}%
				% In these two verses and below in the next two, we see some important differences between the account of Christ's words in Matthew and those here in Luke. To me the most important difference is that all the verbs here in the Luke account are in the second-person plural. There are a few like that in verses 11 and 12 in Matthew, but all of them in Luke are that way.
				% }
		
		% \gr{
		% 22 Μακάριοί ἐστε ὅταν μισήσωσιν ὑμᾶς οἱ ἄνθρωποι, καὶ ὅταν ἀφορίσωσιν ὑμᾶς καὶ ὀνειδίσωσιν καὶ ἐκβάλωσιν τὸ ὄνομα ὑμῶν ὡς πονηρὸν ἕνεκα τοῦ υἱοῦ τοῦ ἀνθρώπου·
		% 23 χάρητε ἐν ἐκείνῃ τῇ ἡμέρᾳ καὶ σκιρτήσατε, ἰδοὺ γὰρ ὁ μισθὸς ὑμῶν πολὺς ἐν τῷ οὐρανῷ· κατὰ τὰ αὐτὰ γὰρ ἐποίουν τοῖς προφήταις οἱ πατέρες αὐτῶν.
		% }
		
% % -----------------------------------------------------------
% \section{39 | Salt of the Earth, Light of the World}
% % -----------------------------------------------------------
	% \label{LCHRIST-039}

	% % ----------------------
	% \subsection{Matthew 5:13--16; 6:22--23}
	% % ----------------------
		% \label{LCHRIST-039-MT}
		
		% \gr{
		% 13 Ὑμεῖς ἐστε τὸ ἅλας τῆς γῆς· ἐὰν δὲ τὸ ἅλας μωρανθῇ, ἐν τίνι ἁλισθήσεται; εἰς οὐδὲν ἰσχύει ἔτι εἰ μὴ βληθὲν ἔξω καταπατεῖσθαι ὑπὸ τῶν ἀνθρώπων.
		% }
		
		% \gr{
		% 14 Ὑμεῖς ἐστε τὸ φῶς τοῦ κόσμου. οὐ δύναται πόλις κρυβῆναι ἐπάνω ὄρους κειμένη·
		% 15 οὐδὲ καίουσιν λύχνον καὶ τιθέασιν αὐτὸν ὑπὸ τὸν μόδιον ἀλλ᾽ ἐπὶ τὴν λυχνίαν, καὶ λάμπει πᾶσιν τοῖς ἐν τῇ οἰκίᾳ.
		% 16 οὕτως λαμψάτω τὸ φῶς ὑμῶν ἔμπροσθεν τῶν ἀνθρώπων, ὅπως ἴδωσιν ὑμῶν τὰ καλὰ ἔργα καὶ δοξάσωσιν τὸν πατέρα ὑμῶν τὸν ἐν τοῖς οὐρανοῖς.
		% }
		
		% \gr{
		% 22 Ὁ λύχνος τοῦ σώματός ἐστιν ὁ ὀφθαλμός. ἐὰν οὖν ᾖ ὁ ὀφθαλμός σου} \hyperref[HAPLOUS]{\grl{ἁπλοῦς}}\gr{, ὅλον τὸ σῶμά σου φωτεινὸν ἔσται·
		% 23 ἐὰν δὲ ὁ ὀφθαλμός σου πονηρὸς ᾖ, ὅλον τὸ σῶμά σου σκοτεινὸν ἔσται.}%
		% \footnote{%
			% \label{LCHRIST-039-haplous}%
			% See \hyperref[HAPLOUS]{\grl{ἁπλοῦς}} for that Greek word; in the KJV it's translated as ``single,'' and see \hyperref[SINGLE-1828]{here} for that word (definition \#8 is the one used to translate \gr{ἁπλοῦς}). I find the contrast made between \gr{ἁπλοῦς} on the one hand and \gr{πονηρός} to be interesting. You can be \gr{ἁπλοῦς} to others, but you can also be (or not be) \gr{ἁπλοῦς} with yourself; the latter is much more important. Part of this passage's message is about not lying to yourself, being honest about what you really want, in the way you live.
			
			% See also note \hyperref[LCHRIST-052-luchnos]{this note} on the same passage below.
			% } 
		% \gr{εἰ οὖν τὸ φῶς τὸ ἐν σοὶ σκότος ἐστίν, τὸ σκότος πόσον.
		% }
	
	% % ----------------------
	% \subsection{Mark 9:50}
	% % ----------------------
		% \label{LCHRIST-039-MK}
		
		% \gr{
		% 50 καλὸν τὸ ἅλας· ἐὰν δὲ τὸ ἅλας ἄναλον γένηται, ἐν τίνι αὐτὸ ἀρτύσετε; ἔχετε ἐν ἑαυτοῖς ἅλα, καὶ εἰρηνεύετε ἐν ἀλλήλοις.
		% }
	
	% % ----------------------
	% \subsection{Luke 14:34--35}
	% % ----------------------
		% \label{LCHRIST-039-L}
		
		% \gr{
		% 34 Καλὸν οὖν τὸ ἅλας· ἐὰν δὲ καὶ τὸ ἅλας μωρανθῇ, ἐν τίνι ἀρτυθήσεται;
		% 35 οὔτε εἰς γῆν οὔτε εἰς κοπρίαν εὔθετόν ἐστιν· ἔξω βάλλουσιν αὐτό. ὁ ἔχων ὦτα ἀκούειν ἀκουέτω.
		% }
	
% % -----------------------------------------------------------
% \section{40 | Jesus Came to Fulfil the Law}
% % -----------------------------------------------------------
	% \label{LCHRIST-040}

	% % ----------------------
	% \subsection{Matthew 5:17--20}
	% % ----------------------
		% \label{LCHRIST-040-MT}
		
		% \gr{
		% 17 Μὴ νομίσητε ὅτι ἦλθον καταλῦσαι τὸν νόμον ἢ τοὺς προφήτας· οὐκ ἦλθον καταλῦσαι ἀλλὰ πληρῶσαι·
		% 18 ἀμὴν γὰρ λέγω ὑμῖν, ἕως ἂν παρέλθῃ ὁ οὐρανὸς καὶ ἡ γῆ, ἰῶτα ἓν ἢ μία κεραία οὐ μὴ παρέλθῃ ἀπὸ τοῦ νόμου, ἕως ἂν πάντα γένηται.
		% 19 ὃς ἐὰν οὖν λύσῃ μίαν τῶν ἐντολῶν τούτων τῶν ἐλαχίστων καὶ διδάξῃ οὕτως τοὺς ἀνθρώπους, ἐλάχιστος κληθήσεται ἐν τῇ βασιλείᾳ τῶν οὐρανῶν· ὃς δ᾽ ἂν ποιήσῃ καὶ διδάξῃ, οὗτος μέγας κληθήσεται ἐν τῇ βασιλείᾳ τῶν οὐρανῶν.
		% 20 λέγω γὰρ ὑμῖν ὅτι ἐὰν μὴ περισσεύσῃ ὑμῶν ἡ δικαιοσύνη πλεῖον τῶν γραμματέων καὶ Φαρισαίων, οὐ μὴ εἰσέλθητε εἰς τὴν βασιλείαν τῶν οὐρανῶν.}%
			% \footnote{%
				% \label{LCHRIST-040-scribes}%
				% I think one key to understanding what Christ means here is in the first part of the next event, \hyperref[LCHRIST-041-MT]{event \#41}, in Christ's life, where he talks about murdering and anger. He goes through the old way of doing things and then says: \gr{ἐγὼ δὲ λέγω ὑμῖν ὅτι πᾶς ὁ ὀργιζόμενος τῷ ἀδελφῷ αὐτοῦ ἔνοχος ἔσται τῇ κρίσει·}. This is the same condemnation that was previously promised in the old way of doing things. So part of the way to have our righteousness overflow more than that of the scribes is to undergo internal changes to our patterns of behavior and habits in order to go beyond the law. We have to be better on the inside; we have to be cleansed from the inside out. 
				% }
	
% % -----------------------------------------------------------
% \section{41 | Be Not Angry, Agree with Adversary}
% % -----------------------------------------------------------
	% \label{LCHRIST-041}

	% % ----------------------
	% \subsection{Matthew 5:21--26}
	% % ----------------------
		% \label{LCHRIST-041-MT}
		
		% \gr{
		% 21 Ἠκούσατε ὅτι ἐρρέθη τοῖς ἀρχαίοις· Οὐ φονεύσεις· ὃς δ᾽ ἂν φονεύσῃ, ἔνοχος ἔσται τῇ κρίσει.
		% 22 ἐγὼ δὲ λέγω ὑμῖν ὅτι πᾶς ὁ ὀργιζόμενος τῷ ἀδελφῷ αὐτοῦ ἔνοχος ἔσται τῇ κρίσει· ὃς δ᾽ ἂν εἴπῃ τῷ ἀδελφῷ αὐτοῦ· Ῥακά, ἔνοχος ἔσται τῷ συνεδρίῳ· ὃς δ᾽ ἂν εἴπῃ· Μωρέ, ἔνοχος ἔσται εἰς τὴν γέενναν τοῦ πυρός.
		% 23 ἐὰν οὖν προσφέρῃς τὸ δῶρόν σου ἐπὶ τὸ θυσιαστήριον κἀκεῖ μνησθῇς ὅτι ὁ ἀδελφός σου ἔχει τι κατὰ σοῦ,
		% 24 ἄφες ἐκεῖ τὸ δῶρόν σου ἔμπροσθεν τοῦ θυσιαστηρίου καὶ ὕπαγε πρῶτον διαλλάγηθι τῷ ἀδελφῷ σου, καὶ τότε ἐλθὼν πρόσφερε τὸ δῶρόν σου.
		% 25 ἴσθι εὐνοῶν τῷ ἀντιδίκῳ σου ταχὺ ἕως ὅτου εἶ μετ᾽ αὐτοῦ ἐν τῇ ὁδῷ, μήποτέ σε παραδῷ ὁ ἀντίδικος τῷ κριτῇ, καὶ ὁ κριτὴς τῷ ὑπηρέτῃ, καὶ εἰς φυλακὴν βληθήσῃ·
		% 26 ἀμὴν λέγω σοι, οὐ μὴ ἐξέλθῃς ἐκεῖθεν ἕως ἂν ἀποδῷς τὸν ἔσχατον κοδράντην.
		% }
	
	% % ----------------------
	% \subsection{Luke 12:57--59}
	% % ----------------------
		% \label{LCHRIST-041-L}
		
		% \gr{
		% 57 Τί δὲ καὶ ἀφ᾽ ἑαυτῶν οὐ κρίνετε τὸ δίκαιον;
		% 58 ὡς γὰρ ὑπάγεις μετὰ τοῦ ἀντιδίκου σου ἐπ᾽ ἄρχοντα, ἐν τῇ ὁδῷ δὸς ἐργασίαν ἀπηλλάχθαι ἀπ᾽ αὐτοῦ, μήποτε κατασύρῃ σε πρὸς τὸν κριτήν, καὶ ὁ κριτής σε παραδώσει τῷ πράκτορι, καὶ ὁ πράκτωρ σε βαλεῖ εἰς φυλακήν.
		% 59 λέγω σοι, οὐ μὴ ἐξέλθῃς ἐκεῖθεν ἕως καὶ τὸ ἔσχατον λεπτὸν ἀποδῷς.
		% }
	
% % -----------------------------------------------------------
% \section{42 | The Higher Law on Adultery}
% % -----------------------------------------------------------
	% \label{LCHRIST-042}

	% % ----------------------
	% \subsection{Matthew 5:27--30}
	% % ----------------------
		% \label{LCHRIST-042-MT}
		
		% \gr{
		% 27 Ἠκούσατε ὅτι ἐρρέθη· Οὐ μοιχεύσεις.
		% 28 ἐγὼ δὲ λέγω ὑμῖν ὅτι πᾶς ὁ βλέπων γυναῖκα πρὸς τὸ ἐπιθυμῆσαι αὐτὴν ἤδη ἐμοίχευσεν αὐτὴν ἐν τῇ καρδίᾳ αὐτοῦ.
		% 29 εἰ δὲ ὁ ὀφθαλμός σου ὁ δεξιὸς σκανδαλίζει σε, ἔξελε αὐτὸν καὶ βάλε ἀπὸ σοῦ, συμφέρει γάρ σοι ἵνα ἀπόληται ἓν τῶν μελῶν σου καὶ μὴ ὅλον τὸ σῶμά σου βληθῇ εἰς γέενναν.
		% 30 καὶ εἰ ἡ δεξιά σου χεὶρ σκανδαλίζει σε, ἔκκοψον αὐτὴν καὶ βάλε ἀπὸ σοῦ, συμφέρει γάρ σοι ἵνα ἀπόληται ἓν τῶν μελῶν σου καὶ μὴ ὅλον τὸ σῶμά σου εἰς γέενναν ἀπέλθῃ.
		% }
	
% % -----------------------------------------------------------
% \section{43 | The Higher Law on Divorce}
% % -----------------------------------------------------------
	% \label{LCHRIST-043}

	% % ----------------------
	% \subsection{Matthew 5:31--32; 19:9}
	% % ----------------------
		% \label{LCHRIST-043-MT}
		
		% \gr{
		% 31 Ἐρρέθη δέ· Ὃς ἂν ἀπολύσῃ τὴν γυναῖκα αὐτοῦ, δότω αὐτῇ ἀποστάσιον.
		% 32 ἐγὼ δὲ λέγω ὑμῖν ὅτι πᾶς ὁ ἀπολύων τὴν γυναῖκα αὐτοῦ παρεκτὸς λόγου πορνείας ποιεῖ αὐτὴν μοιχευθῆναι, καὶ ὃς ἐὰν ἀπολελυμένην γαμήσῃ μοιχᾶται.
		% }
		
		% \gr{
		% 9 λέγω δὲ ὑμῖν ὅτι ὃς ἂν ἀπολύσῃ τὴν γυναῖκα αὐτοῦ μὴ ἐπὶ πορνείᾳ καὶ γαμήσῃ ἄλλην μοιχᾶται καὶ ὁ ἀπολελυμένην γαμήσας μοιχᾶται.
		% }
	
	% % ----------------------
	% \subsection{Mark 10:11--12}
	% % ----------------------
		% \label{LCHRIST-043-MK}
		
		% \gr{
		% 11 καὶ λέγει αὐτοῖς· Ὃς ἂν ἀπολύσῃ τὴν γυναῖκα αὐτοῦ καὶ γαμήσῃ ἄλλην μοιχᾶται ἐπ᾽ αὐτήν,
		% 12 καὶ ἐὰν αὐτὴ ἀπολύσασα τὸν ἄνδρα αὐτῆς γαμήσῃ ἄλλον μοιχᾶται.
		% }
	
	% % ----------------------
	% \subsection{Luke 16:18}
	% % ----------------------
		% \label{LCHRIST-043-L}
		
		% \gr{
		% 18 Πᾶς ὁ ἀπολύων τὴν γυναῖκα αὐτοῦ καὶ γαμῶν ἑτέραν μοιχεύει, καὶ ὁ ἀπολελυμένην ἀπὸ ἀνδρὸς γαμῶν μοιχεύει.
		% }
	
% % -----------------------------------------------------------
% \section{44 | Swear Not at All}
% % -----------------------------------------------------------
	% \label{LCHRIST-044}

	% % ----------------------
	% \subsection{Matthew 5:33--37}
	% % ----------------------
		% \label{LCHRIST-044-MT}
		
		% \gr{
		% 33 Πάλιν ἠκούσατε ὅτι ἐρρέθη τοῖς ἀρχαίοις· Οὐκ ἐπιορκήσεις, ἀποδώσεις δὲ τῷ κυρίῳ τοὺς ὅρκους σου.
		% 34 ἐγὼ δὲ λέγω ὑμῖν μὴ ὀμόσαι ὅλως· μήτε ἐν τῷ οὐρανῷ, ὅτι θρόνος ἐστὶν τοῦ θεοῦ·
		% 35 μήτε ἐν τῇ γῇ, ὅτι ὑποπόδιόν ἐστιν τῶν ποδῶν αὐτοῦ· μήτε εἰς Ἱεροσόλυμα, ὅτι πόλις ἐστὶν τοῦ μεγάλου βασιλέως·
		% 36 μήτε ἐν τῇ κεφαλῇ σου ὀμόσῃς, ὅτι οὐ δύνασαι μίαν τρίχα λευκὴν ποιῆσαι ἢ μέλαιναν.
		% 37 ἔστω δὲ ὁ λόγος ὑμῶν ναὶ ναί, οὒ οὔ· τὸ δὲ περισσὸν τούτων ἐκ τοῦ πονηροῦ ἐστιν.
		% }
	
% % -----------------------------------------------------------
% \section{45 | The Higher Law on Retaliation}
% % -----------------------------------------------------------
	% \label{LCHRIST-045}
	
	% % ----------------------
	% \subsection{Matthew 5:38--42}
	% % ----------------------
		% \label{LCHRIST-045-MT}
	
		% \gr{
		% 38 Ἠκούσατε ὅτι ἐρρέθη· Ὀφθαλμὸν ἀντὶ ὀφθαλμοῦ καὶ ὀδόντα ἀντὶ ὀδόντος.
		% 39 ἐγὼ δὲ λέγω ὑμῖν μὴ ἀντιστῆναι τῷ πονηρῷ· ἀλλ᾽ ὅστις σε ῥαπίζει εἰς τὴν δεξιὰν σιαγόνα, στρέψον αὐτῷ καὶ τὴν ἄλλην·
		% 40 καὶ τῷ θέλοντί σοι κριθῆναι καὶ τὸν χιτῶνά σου λαβεῖν, ἄφες αὐτῷ καὶ τὸ ἱμάτιον·
		% 41 καὶ ὅστις σε ἀγγαρεύσει μίλιον ἕν, ὕπαγε μετ᾽ αὐτοῦ δύο.
		% 42 τῷ αἰτοῦντί σε δός, καὶ τὸν θέλοντα ἀπὸ σοῦ δανίσασθαι μὴ ἀποστραφῇς.
		% }
		
	% % ----------------------
	% \subsection{Luke 6:29--30}
	% % ----------------------
		% \label{LCHRIST-045-L}
	
		% \gr{
		% 29 τῷ τύπτοντί σε ἐπὶ τὴν σιαγόνα πάρεχε καὶ τὴν ἄλλην, καὶ ἀπὸ τοῦ αἴροντός σου τὸ ἱμάτιον καὶ τὸν χιτῶνα μὴ κωλύσῃς.
		% 30 παντὶ αἰτοῦντί σε δίδου, καὶ ἀπὸ τοῦ αἴροντος τὰ σὰ μὴ ἀπαίτει.
		% }
	
% % -----------------------------------------------------------
% \section{46 | Love Your Enemies, Bless Them that Curse You}
% % -----------------------------------------------------------	
	% \label{LCHRIST-046}
	
	% % ----------------------
	% \subsection{Matthew 5:43--48; 7:12}
	% % ----------------------
		% \label{LCHRIST-046-MT}
		
		% \gr{
		% 43 Ἠκούσατε ὅτι ἐρρέθη· Ἀγαπήσεις τὸν πλησίον σου καὶ μισήσεις τὸν ἐχθρόν σου.
		% 44 ἐγὼ δὲ λέγω ὑμῖν, ἀγαπᾶτε τοὺς ἐχθροὺς ὑμῶν καὶ προσεύχεσθε ὑπὲρ τῶν διωκόντων ὑμᾶς·
		% 45 ὅπως γένησθε υἱοὶ τοῦ πατρὸς ὑμῶν τοῦ ἐν οὐρανοῖς, ὅτι τὸν ἥλιον αὐτοῦ ἀνατέλλει ἐπὶ πονηροὺς καὶ ἀγαθοὺς καὶ βρέχει ἐπὶ δικαίους καὶ ἀδίκους.
		% 46 ἐὰν γὰρ ἀγαπήσητε τοὺς ἀγαπῶντας ὑμᾶς, τίνα μισθὸν ἔχετε; οὐχὶ καὶ οἱ τελῶναι τὸ αὐτὸ ποιοῦσιν;
		% 47 καὶ ἐὰν ἀσπάσησθε τοὺς ἀδελφοὺς ὑμῶν μόνον, τί περισσὸν ποιεῖτε; οὐχὶ καὶ οἱ ἐθνικοὶ τὸ αὐτὸ ποιοῦσιν;
		% 48 Ἔσεσθε οὖν ὑμεῖς τέλειοι ὡς ὁ πατὴρ ὑμῶν ὁ οὐράνιος τέλειός ἐστιν.
		% }
		
		% \gr{
		% 12 Πάντα οὖν ὅσα ἐὰν θέλητε ἵνα ποιῶσιν ὑμῖν οἱ ἄνθρωποι, οὕτως καὶ ὑμεῖς ποιεῖτε αὐτοῖς· οὗτος γάρ ἐστιν ὁ νόμος καὶ οἱ προφῆται.
		% }
		
	% % ----------------------
	% \subsection{Luke 6:27--36}
	% % ----------------------
		% \label{LCHRIST-046-L}
		
		% \gr{
		% 27 Ἀλλὰ ὑμῖν λέγω τοῖς ἀκούουσιν, ἀγαπᾶτε τοὺς ἐχθροὺς ὑμῶν, καλῶς ποιεῖτε τοῖς μισοῦσιν ὑμᾶς,
		% 28 εὐλογεῖτε τοὺς καταρωμένους ὑμᾶς, προσεύχεσθε περὶ τῶν ἐπηρεαζόντων ὑμᾶς.
		% 29 τῷ τύπτοντί σε ἐπὶ τὴν σιαγόνα πάρεχε καὶ τὴν ἄλλην, καὶ ἀπὸ τοῦ αἴροντός σου τὸ ἱμάτιον καὶ τὸν χιτῶνα μὴ κωλύσῃς.
		% 30 παντὶ αἰτοῦντί σε δίδου, καὶ ἀπὸ τοῦ αἴροντος τὰ σὰ μὴ ἀπαίτει.
		% 31 καὶ καθὼς θέλετε ἵνα ποιῶσιν ὑμῖν οἱ ἄνθρωποι, ποιεῖτε αὐτοῖς ὁμοίως.
		% 32 Καὶ εἰ ἀγαπᾶτε τοὺς ἀγαπῶντας ὑμᾶς, ποία ὑμῖν χάρις ἐστίν; καὶ γὰρ οἱ ἁμαρτωλοὶ τοὺς ἀγαπῶντας αὐτοὺς ἀγαπῶσιν.}
			% \footnote{%
				% \label{LCHRIST-046-charis}%
				% Note that in this verse and in the next two, Christ speaks of a person getting or not getting \gr{χάρις} in exchange for doing or not doing something. This is a very important passage that helps us understand a crucial phrase appearing again in restoration scripture, ``grace for grace'' (see \hyperref[GRACE-PHRASES]{here}). The connection is not apparent in the KJV translation because all three instances of the question get translated as ``what thank have ye?'' It's also important to undestanding the process of becoming children of God (in some sense beyond what we already are).
				% }
		% \gr{33 καὶ ἐὰν ἀγαθοποιῆτε τοὺς ἀγαθοποιοῦντας ὑμᾶς, ποία ὑμῖν χάρις ἐστίν; καὶ οἱ ἁμαρτωλοὶ τὸ αὐτὸ ποιοῦσιν.
		% 34 καὶ ἐὰν δανίσητε παρ᾽ ὧν ἐλπίζετε λαβεῖν, ποία ὑμῖν χάρις ἐστίν; καὶ ἁμαρτωλοὶ ἁμαρτωλοῖς δανίζουσιν ἵνα ἀπολάβωσιν τὰ ἴσα.
		% 35 πλὴν ἀγαπᾶτε τοὺς ἐχθροὺς ὑμῶν καὶ ἀγαθοποιεῖτε καὶ δανίζετε μηδὲν ἀπελπίζοντες· καὶ ἔσται ὁ μισθὸς ὑμῶν πολύς, καὶ ἔσεσθε υἱοὶ Ὑψίστου, ὅτι αὐτὸς χρηστός ἐστιν ἐπὶ τοὺς ἀχαρίστους καὶ πονηρούς.
		% 36 γίνεσθε οἰκτίρμονες καθὼς ὁ πατὴρ ὑμῶν οἰκτίρμων ἐστίν·
		% }
		
% % -----------------------------------------------------------
% \section{47 | Give Alms in Secret}
% % -----------------------------------------------------------
	% \label{LCHRIST-047}
	
	% % ----------------------
	% \subsection{Matthew 6:1--4}
	% % ----------------------
		% \label{LCHRIST-047-MT}
		
		% \gr{
		% 1 Προσέχετε δὲ τὴν δικαιοσύνην}%
			% \footnote{%
				% \label{LCHRIST-047-dikaiosunen}%
				% For the object of the verb, this Greek text has \gr{δικαιοσύνην}, while the source the KJV goes with has \gr{ἐλεημοσύνην}. I think the former is more interesting because it tells us that Christ's advice here applies to any instance of us trying to do an act of \gr{δικαιοσύνη}. It is not and cannot be about us being seen by others, and if that's our concern, then we get no \gr{μισθός} from God. 
				% }
		% \gr{ὑμῶν μὴ ποιεῖν ἔμπροσθεν τῶν ἀνθρώπων πρὸς τὸ θεαθῆναι αὐτοῖς· εἰ δὲ μή γε, μισθὸν οὐκ ἔχετε παρὰ τῷ πατρὶ ὑμῶν τῷ ἐν τοῖς οὐρανοῖς.
		% }
		
		% \gr{
		% 2 Ὅταν οὖν ποιῇς ἐλεημοσύνην, μὴ σαλπίσῃς ἔμπροσθέν σου, ὥσπερ οἱ ὑποκριταὶ ποιοῦσιν ἐν ταῖς συναγωγαῖς καὶ ἐν ταῖς ῥύμαις, ὅπως δοξασθῶσιν ὑπὸ τῶν ἀνθρώπων· ἀμὴν λέγω ὑμῖν, ἀπέχουσιν τὸν μισθὸν αὐτῶν.}
			% \footnote{%
				% \label{LCHRIST-047-apechousin}%
				% Note that the KJV has ``They have their reward'' for the end of this verse. That doesn't really capture the force of \gr{ἀπέχω}, whose first definition is ``to receive in full what is due, \textit{to be paid in full, receive in full}.'' You don't get the ``in full'' idea, nor the idea of actually receiving it. A better translation would be ``I say to you, they receive their reward in full,'' or ``they are receiving their reward in full.''
				% }
		% \gr{3 σοῦ δὲ ποιοῦντος ἐλεημοσύνην μὴ γνώτω ἡ ἀριστερά σου τί ποιεῖ ἡ δεξιά σου,
		% 4 ὅπως ᾖ σου ἡ ἐλεημοσύνη ἐν τῷ κρυπτῷ· καὶ ὁ πατήρ σου ὁ βλέπων ἐν τῷ κρυπτῷ ἀποδώσει σοι.}%
			% \footnote{%
				% \label{LCHRIST-047-apodosei}%
				% Note that there's no ``openly'' here, as in the KJV.
				% }
		
% % -----------------------------------------------------------
% \section{48 | On Prayer}
% % -----------------------------------------------------------
	% \label{LCHRIST-048}
	
	% % ----------------------
	% \subsection{Matthew 6:5--15; 7:7--11}
	% % ----------------------
		% \label{LCHRIST-048-MT}
		
		% \gr{
		% 5 Καὶ ὅταν προσεύχησθε, οὐκ ἔσεσθε ὡς οἱ ὑποκριταί· ὅτι φιλοῦσιν ἐν ταῖς συναγωγαῖς καὶ ἐν ταῖς γωνίαις τῶν πλατειῶν ἑστῶτες προσεύχεσθαι, ὅπως φανῶσιν τοῖς ἀνθρώποις· ἀμὴν λέγω ὑμῖν, ἀπέχουσι τὸν μισθὸν αὐτῶν.
		% 6 σὺ δὲ ὅταν προσεύχῃ, εἴσελθε εἰς τὸ ταμεῖόν σου καὶ κλείσας τὴν θύραν σου πρόσευξαι τῷ πατρί σου τῷ ἐν τῷ κρυπτῷ· καὶ ὁ πατήρ σου ὁ βλέπων ἐν τῷ κρυπτῷ ἀποδώσει σοι.
		% }
		
		% \gr{
		% 7 Προσευχόμενοι δὲ μὴ βατταλογήσητε ὥσπερ οἱ ἐθνικοί, δοκοῦσιν γὰρ ὅτι ἐν τῇ πολυλογίᾳ αὐτῶν εἰσακουσθήσονται·
		% 8 μὴ οὖν ὁμοιωθῆτε αὐτοῖς, οἶδεν γὰρ ὁ πατὴρ ὑμῶν ὧν χρείαν ἔχετε πρὸ τοῦ ὑμᾶς αἰτῆσαι αὐτόν.
		% }
		
		% \gr{
		% 9 Οὕτως οὖν προσεύχεσθε ὑμεῖς· Πάτερ ἡμῶν ὁ ἐν τοῖς οὐρανοῖς· ἁγιασθήτω τὸ ὄνομά σου,
		% 10 ἐλθέτω ἡ βασιλεία σου, γενηθήτω τὸ θέλημά σου, ὡς ἐν οὐρανῷ καὶ ἐπὶ γῆς·
		% 11 τὸν ἄρτον ἡμῶν τὸν ἐπιούσιον δὸς ἡμῖν σήμερον·
		% 12 καὶ ἄφες ἡμῖν τὰ ὀφειλήματα ἡμῶν, ὡς καὶ ἡμεῖς ἀφήκαμεν τοῖς ὀφειλέταις ἡμῶν·
		% 13 καὶ μὴ εἰσενέγκῃς ἡμᾶς εἰς πειρασμόν, ἀλλὰ ῥῦσαι ἡμᾶς ἀπὸ τοῦ πονηροῦ.
		% 14 ἐὰν γὰρ ἀφῆτε τοῖς ἀνθρώποις τὰ παραπτώματα αὐτῶν, ἀφήσει καὶ ὑμῖν ὁ πατὴρ ὑμῶν ὁ οὐράνιος·
		% 15 ἐὰν δὲ μὴ ἀφῆτε τοῖς ἀνθρώποις, οὐδὲ ὁ πατὴρ ὑμῶν ἀφήσει τὰ παραπτώματα ὑμῶν.
		% }
		
		% \gr{
		% 7 Αἰτεῖτε, καὶ δοθήσεται ὑμῖν· ζητεῖτε, καὶ εὑρήσετε· κρούετε, καὶ ἀνοιγήσεται ὑμῖν.}
			% \footnote{%
				% \label{LCHRIST-048-participles}%
				% Note that all these participles are present, meaning they carry the idea that someone \textit{continues to perform} these actions---``Be asking, and it will be given to you; be seeking, and you will find; be knocking, and it will be opened to you. For each person who keeps asking receives, and the one who keeps seeking finds, and to the one who keeps knocking, it will be opened.''
				% }
		% \gr{8 πᾶς γὰρ ὁ αἰτῶν λαμβάνει καὶ ὁ ζητῶν εὑρίσκει καὶ τῷ κρούοντι ἀνοιγήσεται.
		% 9 ἢ τίς ἐστιν ἐξ ὑμῶν ἄνθρωπος, ὃν αἰτήσει ὁ υἱὸς αὐτοῦ ἄρτον--- μὴ λίθον ἐπιδώσει αὐτῷ;
		% 10 ἢ καὶ ἰχθὺν αἰτήσει--- μὴ ὄφιν ἐπιδώσει αὐτῷ;
		% 11 εἰ οὖν ὑμεῖς πονηροὶ ὄντες οἴδατε δόματα ἀγαθὰ διδόναι τοῖς τέκνοις ὑμῶν, πόσῳ μᾶλλον ὁ πατὴρ ὑμῶν ὁ ἐν τοῖς οὐρανοῖς δώσει ἀγαθὰ τοῖς αἰτοῦσιν αὐτόν.
		% }
		
	% % ----------------------
	% \subsection{Luke 11:1--13}
	% % ----------------------
		% \label{LCHRIST-048-L}
		
		% \gr{
		% 1 Καὶ ἐγένετο ἐν τῷ εἶναι αὐτὸν ἐν τόπῳ τινὶ προσευχόμενον, ὡς ἐπαύσατο, εἶπέν τις τῶν μαθητῶν αὐτοῦ πρὸς αὐτόν· Κύριε, δίδαξον ἡμᾶς προσεύχεσθαι, καθὼς καὶ Ἰωάννης ἐδίδαξεν τοὺς μαθητὰς αὐτοῦ.
		% 2 εἶπεν δὲ αὐτοῖς· Ὅταν προσεύχησθε, λέγετε· Πάτερ, ἁγιασθήτω τὸ ὄνομά σου· ἐλθέτω ἡ βασιλεία σου·
		% 3 τὸν ἄρτον ἡμῶν τὸν ἐπιούσιον δίδου ἡμῖν τὸ καθ᾽ ἡμέραν·
		% 4 καὶ ἄφες ἡμῖν τὰς ἁμαρτίας ἡμῶν, καὶ γὰρ αὐτοὶ ἀφίομεν παντὶ ὀφείλοντι ἡμῖν· καὶ μὴ εἰσενέγκῃς ἡμᾶς εἰς πειρασμόν.
		% }
		
		% \gr{
		% 5 Καὶ εἶπεν πρὸς αὐτούς· Τίς ἐξ ὑμῶν ἕξει φίλον καὶ πορεύσεται πρὸς αὐτὸν μεσονυκτίου καὶ εἴπῃ αὐτῷ· Φίλε, χρῆσόν μοι τρεῖς ἄρτους,
		% 6 ἐπειδὴ φίλος μου παρεγένετο ἐξ ὁδοῦ πρός με καὶ οὐκ ἔχω ὃ παραθήσω αὐτῷ·
		% 7 κἀκεῖνος ἔσωθεν ἀποκριθεὶς εἴπῃ· Μή μοι κόπους πάρεχε· ἤδη ἡ θύρα κέκλεισται, καὶ τὰ παιδία μου μετ᾽ ἐμοῦ εἰς τὴν κοίτην εἰσίν· οὐ δύναμαι ἀναστὰς δοῦναί σοι.
		% 8 λέγω ὑμῖν, εἰ καὶ οὐ δώσει αὐτῷ ἀναστὰς διὰ τὸ εἶναι φίλον αὐτοῦ, διά γε τὴν ἀναίδειαν αὐτοῦ ἐγερθεὶς δώσει αὐτῷ ὅσων χρῄζει.
		% }
		
		% \gr{
		% 9 Κἀγὼ ὑμῖν λέγω, αἰτεῖτε, καὶ δοθήσεται ὑμῖν· ζητεῖτε, καὶ εὑρήσετε· κρούετε, καὶ ἀνοιγήσεται ὑμῖν·
		% 10 πᾶς γὰρ ὁ αἰτῶν λαμβάνει, καὶ ὁ ζητῶν εὑρίσκει, καὶ τῷ κρούοντι ἀνοιγήσεται.
		% 11 τίνα δὲ ἐξ ὑμῶν τὸν πατέρα αἰτήσει ὁ υἱὸς ἰχθύν, καὶ ἀντὶ ἰχθύος ὄφιν αὐτῷ ἐπιδώσει;
		% 12 ἢ καὶ αἰτήσει ᾠόν, ἐπιδώσει αὐτῷ σκορπίον;
		% 13 εἰ οὖν ὑμεῖς πονηροὶ ὑπάρχοντες οἴδατε δόματα ἀγαθὰ διδόναι τοῖς τέκνοις ὑμῶν, πόσῳ μᾶλλον ὁ πατὴρ ὁ ἐξ οὐρανοῦ δώσει πνεῦμα ἅγιον τοῖς αἰτοῦσιν αὐτόν.
		% }
		
% % -----------------------------------------------------------
% \section{49 | The Lord's Prayer}
% % -----------------------------------------------------------
	% \label{LCHRIST-049}
	
	% % ----------------------
	% \subsection{Matthew 6:9--15}
	% % ---------------------- 
		% \label{LCHRIST-049-MT}
		
		% \gr{
		% 9 Οὕτως οὖν προσεύχεσθε ὑμεῖς· Πάτερ ἡμῶν ὁ ἐν τοῖς οὐρανοῖς· ἁγιασθήτω τὸ ὄνομά σου,
		% 10 ἐλθέτω ἡ βασιλεία σου, γενηθήτω τὸ θέλημά σου, ὡς ἐν οὐρανῷ καὶ ἐπὶ γῆς·
		% 11 τὸν ἄρτον ἡμῶν τὸν ἐπιούσιον δὸς ἡμῖν σήμερον·
		% 12 καὶ ἄφες ἡμῖν τὰ ὀφειλήματα ἡμῶν, ὡς καὶ ἡμεῖς ἀφήκαμεν τοῖς ὀφειλέταις ἡμῶν·
		% 13 καὶ μὴ εἰσενέγκῃς ἡμᾶς εἰς πειρασμόν, ἀλλὰ ῥῦσαι ἡμᾶς ἀπὸ τοῦ πονηροῦ.
		% 14 ἐὰν γὰρ ἀφῆτε τοῖς ἀνθρώποις τὰ παραπτώματα αὐτῶν, ἀφήσει καὶ ὑμῖν ὁ πατὴρ ὑμῶν ὁ οὐράνιος·
		% 15 ἐὰν δὲ μὴ ἀφῆτε τοῖς ἀνθρώποις, οὐδὲ ὁ πατὴρ ὑμῶν ἀφήσει τὰ παραπτώματα ὑμῶν.
		% }
		
	% % ----------------------
	% \subsection{Luke 11:2--4}
	% % ----------------------
		% \label{LCHRIST-049-L}
		
		% \gr{
		% 2 εἶπεν δὲ αὐτοῖς· Ὅταν προσεύχησθε, λέγετε· Πάτερ, ἁγιασθήτω τὸ ὄνομά σου· ἐλθέτω ἡ βασιλεία σου·
		% 3 τὸν ἄρτον ἡμῶν τὸν ἐπιούσιον δίδου ἡμῖν τὸ καθ᾽ ἡμέραν·
		% 4 καὶ ἄφες ἡμῖν τὰς ἁμαρτίας ἡμῶν, καὶ γὰρ αὐτοὶ ἀφίομεν παντὶ ὀφείλοντι ἡμῖν· καὶ μὴ εἰσενέγκῃς ἡμᾶς εἰς πειρασμόν.
		% }
		
% % -----------------------------------------------------------
% \section{50 | On Fasting}
% % -----------------------------------------------------------
	% \label{LCHRIST-050}
	
	% % ----------------------
	% \subsection{Matthew 6:16--18}
	% % ----------------------
		% \label{LCHRIST-050-MT}
		
		% \gr{
		% 16 Ὅταν δὲ νηστεύητε, μὴ γίνεσθε ὡς οἱ ὑποκριταὶ σκυθρωποί, ἀφανίζουσιν γὰρ τὰ πρόσωπα αὐτῶν ὅπως φανῶσιν τοῖς ἀνθρώποις νηστεύοντες· ἀμὴν λέγω ὑμῖν, ἀπέχουσιν τὸν μισθὸν αὐτῶν.
		% 17 σὺ δὲ νηστεύων ἄλειψαί σου τὴν κεφαλὴν καὶ τὸ πρόσωπόν σου νίψαι,
		% 18 ὅπως μὴ φανῇς τοῖς ἀνθρώποις νηστεύων ἀλλὰ τῷ πατρί σου τῷ ἐν τῷ κρυφαίῳ· καὶ ὁ πατήρ σου ὁ βλέπων ἐν τῷ κρυφαίῳ ἀποδώσει σοι.
		% }
		
% % -----------------------------------------------------------
% \section{51 | Lay Up Treasures in Heaven}
% % -----------------------------------------------------------
	% \label{LCHRIST-051}
	
	% % ----------------------
	% \subsection{Matthew 6:19--21}
	% % ---------------------- 
		% \label{LCHRIST-051-MT}
		
		% \gr{
		% 19 Μὴ θησαυρίζετε ὑμῖν θησαυροὺς ἐπὶ τῆς γῆς, ὅπου σὴς καὶ βρῶσις ἀφανίζει, καὶ ὅπου κλέπται διορύσσουσιν καὶ κλέπτουσιν·
		% 20 θησαυρίζετε δὲ ὑμῖν θησαυροὺς ἐν οὐρανῷ, ὅπου οὔτε σὴς οὔτε βρῶσις ἀφανίζει, καὶ ὅπου κλέπται οὐ διορύσσουσιν οὐδὲ κλέπτουσιν·
		% 21 ὅπου γάρ ἐστιν ὁ θησαυρός σου, ἐκεῖ ἔσται καὶ ἡ καρδία σου.
		% }
		
	% % ----------------------
	% \subsection{Luke 12:33--34}
	% % ----------------------
		% \label{LCHRIST-051-L}
		
		% \gr{
		% 33 πωλήσατε τὰ ὑπάρχοντα ὑμῶν καὶ δότε ἐλεημοσύνην· ποιήσατε ἑαυτοῖς βαλλάντια μὴ παλαιούμενα, θησαυρὸν ἀνέκλειπτον ἐν τοῖς οὐρανοῖς, ὅπου κλέπτης οὐκ ἐγγίζει οὐδὲ σὴς διαφθείρει·
		% 34 ὅπου γάρ ἐστιν ὁ θησαυρὸς ὑμῶν, ἐκεῖ καὶ ἡ καρδία ὑμῶν ἔσται.
		% }
		
% % -----------------------------------------------------------
% \section{52 | The Eye is the Light of the Body}
% % -----------------------------------------------------------
	% \label{LCHRIST-052}
	
	% % ----------------------
	% \subsection{Matthew 6:22--23}
	% % ----------------------
		% \label{LCHRIST-052-MT}
		
		% \gr{
		% 22 Ὁ λύχνος}%
			% \footnote{%
				% \label{LCHRIST-052-luchnos}%
				% \gr{λύχνος} means ``lamp'' (BDAG 606b, PDF 344), or ``portable light'' (LSJ). It's the source of light that you can take to a certain place to light it up, not the light itself. I wonder if there's also something here about the extramission theory of visual perception, where you see by beams of light emitted from the eye, rather than received by the eye. The extramission theory was very common in antiquity, including during this time, and it gets even bigger with Galen later on, so it's possible that the writers of the gospels had this view of visual perception as well. This is worth looking into in more detail at some other time. Notice the similarity you get between visual perception and \textit{touch}, when you think of vision as occurring when something goes out of the eye. I found a book called \textit{The Sense of Sight in Rabbinic Culture}, by Rachel Neis, which may be relevant; it definitely has some useful discussions in various places. Notice also the connection between this and Christ's comments on \hyperref[LCHRIST-042]{lust} earlier in this sermon; if something of you is going out to interact with a body, then looking at someone with that intent may seem more sinful. In her book, Neis writes, ``In rabbinic culture, a woman’s body was, at least theoretically, governed from head to hem by a thick web of laws prescribing the extent to which visibility was permissible'' (136). There may be other connections to other NT things as well.
				
				% For see \gr{ἁπλοῦς} in this verse, see \hyperref[LCHRIST-039-haplous]{this note}.
				% }
		% \gr{τοῦ σώματός ἐστιν ὁ ὀφθαλμός. ἐὰν οὖν ᾖ ὁ ὀφθαλμός σου ἁπλοῦς, ὅλον τὸ σῶμά σου φωτεινὸν ἔσται·
		% 23 ἐὰν δὲ ὁ ὀφθαλμός σου πονηρὸς ᾖ, ὅλον τὸ σῶμά σου σκοτεινὸν ἔσται. εἰ οὖν τὸ φῶς τὸ ἐν σοὶ σκότος ἐστίν, τὸ σκότος πόσον.
		% }
		
	% % ----------------------
	% \subsection{Luke 11:34--36}
	% % ----------------------
		% \label{LCHRIST-052-L}
		
		% \gr{
		% 34 ὁ λύχνος τοῦ σώματός ἐστιν ὁ ὀφθαλμός σου. ὅταν ὁ ὀφθαλμός σου ἁπλοῦς ᾖ, καὶ ὅλον τὸ σῶμά σου φωτεινόν ἐστιν· ἐπὰν δὲ πονηρὸς ᾖ, καὶ τὸ σῶμά σου σκοτεινόν.
		% 35 σκόπει οὖν μὴ τὸ φῶς τὸ ἐν σοὶ σκότος ἐστίν.
		% 36 εἰ οὖν τὸ σῶμά σου ὅλον φωτεινόν, μὴ ἔχον μέρος τι σκοτεινόν, ἔσται φωτεινὸν ὅλον ὡς ὅταν ὁ λύχνος τῇ ἀστραπῇ φωτίζῃ σε.
		% }
		
% % -----------------------------------------------------------
% \section{53 | No Man Can Serve Both God and Mammon}
% % -----------------------------------------------------------
	% \label{LCHRIST-053}
	
	% % ----------------------
	% \subsection{Matthew 6:24}
	% % ----------------------
		% \label{LCHRIST-053-MT}
		
		% \gr{
		% 24 Οὐδεὶς δύναται δυσὶ κυρίοις δουλεύειν· ἢ γὰρ τὸν ἕνα μισήσει καὶ τὸν ἕτερον ἀγαπήσει, ἢ ἑνὸς ἀνθέξεται καὶ τοῦ ἑτέρου καταφρονήσει. οὐ δύνασθε θεῷ δουλεύειν καὶ μαμωνᾷ.
		% }
		
	% % ----------------------
	% \subsection{Luke 16:13}
	% % ----------------------
		% \label{LCHRIST-053-L}
		
		% \gr{
		% 13 οὐδεὶς οἰκέτης δύναται δυσὶ κυρίοις δουλεύειν· ἢ γὰρ τὸν ἕνα μισήσει καὶ τὸν ἕτερον ἀγαπήσει, ἢ ἑνὸς ἀνθέξεται καὶ τοῦ ἑτέρου καταφρονήσει. οὐ δύνασθε θεῷ δουλεύειν καὶ μαμωνᾷ.
		% }

% % -----------------------------------------------------------
% \section{54 | Take No Thought for the Morrow}
% % -----------------------------------------------------------
	% \label{LCHRIST-054}
	
	% % ----------------------
	% \subsection{Matthew 6:25--34; 19--21}
	% % ----------------------
		% \label{LCHRIST-054-MT}
		
		% \gr{
		% 25 Διὰ τοῦτο λέγω ὑμῖν· μὴ μεριμνᾶτε τῇ ψυχῇ ὑμῶν τί φάγητε, μηδὲ τῷ σώματι ὑμῶν τί ἐνδύσησθε· οὐχὶ ἡ ψυχὴ πλεῖόν ἐστι τῆς τροφῆς καὶ τὸ σῶμα τοῦ ἐνδύματος;
		% 26 ἐμβλέψατε εἰς τὰ πετεινὰ τοῦ οὐρανοῦ ὅτι οὐ σπείρουσιν οὐδὲ θερίζουσιν οὐδὲ συνάγουσιν εἰς ἀποθήκας, καὶ ὁ πατὴρ ὑμῶν ὁ οὐράνιος τρέφει αὐτά· οὐχ ὑμεῖς μᾶλλον διαφέρετε αὐτῶν;
		% 27 τίς δὲ ἐξ ὑμῶν μεριμνῶν δύναται προσθεῖναι ἐπὶ τὴν ἡλικίαν αὐτοῦ πῆχυν ἕνα;
		% 28 καὶ περὶ ἐνδύματος τί μεριμνᾶτε; καταμάθετε τὰ κρίνα τοῦ ἀγροῦ πῶς αὐξάνουσιν· οὐ κοπιῶσιν οὐδὲ νήθουσιν·
		% 29 λέγω δὲ ὑμῖν ὅτι οὐδὲ Σολομὼν ἐν πάσῃ τῇ δόξῃ αὐτοῦ περιεβάλετο ὡς ἓν τούτων.
		% 30 εἰ δὲ τὸν χόρτον τοῦ ἀγροῦ σήμερον ὄντα καὶ αὔριον εἰς κλίβανον βαλλόμενον ὁ θεὸς οὕτως ἀμφιέννυσιν, οὐ πολλῷ μᾶλλον ὑμᾶς, ὀλιγόπιστοι;
		% 31 μὴ οὖν μεριμνήσητε λέγοντες· Τί φάγωμεν; ἤ· Τί πίωμεν; ἤ· Τί περιβαλώμεθα;
		% 32 πάντα γὰρ ταῦτα τὰ ἔθνη ἐπιζητοῦσιν· οἶδεν γὰρ ὁ πατὴρ ὑμῶν ὁ οὐράνιος ὅτι χρῄζετε τούτων ἁπάντων.
		% 33 ζητεῖτε δὲ πρῶτον τὴν βασιλείαν καὶ τὴν δικαιοσύνην αὐτοῦ, καὶ ταῦτα πάντα προστεθήσεται ὑμῖν.
		% 34 μὴ οὖν μεριμνήσητε εἰς τὴν αὔριον, ἡ γὰρ αὔριον μεριμνήσει αὑτῆς· ἀρκετὸν τῇ ἡμέρᾳ ἡ κακία αὐτῆς.
		% }
		
		% \gr{
		% 19 Μὴ θησαυρίζετε ὑμῖν θησαυροὺς ἐπὶ τῆς γῆς, ὅπου σὴς καὶ βρῶσις ἀφανίζει, καὶ ὅπου κλέπται διορύσσουσιν καὶ κλέπτουσιν·
		% 20 θησαυρίζετε δὲ ὑμῖν θησαυροὺς ἐν οὐρανῷ, ὅπου οὔτε σὴς οὔτε βρῶσις ἀφανίζει, καὶ ὅπου κλέπται οὐ διορύσσουσιν οὐδὲ κλέπτουσιν·
		% 21 ὅπου γάρ ἐστιν ὁ θησαυρός σου, ἐκεῖ ἔσται καὶ ἡ καρδία σου.
		% }
		
	% % ----------------------
	% \subsection{Luke 12:22--34}
	% % ----------------------
		% \label{LCHRIST-054-L}
		
		% \gr{
		% 22 Εἶπεν δὲ πρὸς τοὺς μαθητὰς αὐτοῦ· Διὰ τοῦτο λέγω ὑμῖν, μὴ μεριμνᾶτε τῇ ψυχῇ τί φάγητε, μηδὲ τῷ σώματι τί ἐνδύσησθε.
		% 23 ἡ γὰρ ψυχὴ πλεῖόν ἐστιν τῆς τροφῆς καὶ τὸ σῶμα τοῦ ἐνδύματος.
		% 24 κατανοήσατε τοὺς κόρακας ὅτι οὐ σπείρουσιν οὐδὲ θερίζουσιν, οἷς οὐκ ἔστιν ταμεῖον οὐδὲ ἀποθήκη, καὶ ὁ θεὸς τρέφει αὐτούς· πόσῳ μᾶλλον ὑμεῖς διαφέρετε τῶν πετεινῶν.
		% 25 τίς δὲ ἐξ ὑμῶν μεριμνῶν δύναται ἐπὶ τὴν ἡλικίαν αὐτοῦ προσθεῖναι πῆχυν;
		% 26 εἰ οὖν οὐδὲ ἐλάχιστον δύνασθε, τί περὶ τῶν λοιπῶν μεριμνᾶτε;
		% 27 κατανοήσατε τὰ κρίνα πῶς αὐξάνει· οὐ κοπιᾷ οὐδὲ νήθει· λέγω δὲ ὑμῖν, οὐδὲ Σολομὼν ἐν πάσῃ τῇ δόξῃ αὐτοῦ περιεβάλετο ὡς ἓν τούτων.
		% 28 εἰ δὲ ἐν ἀγρῷ τὸν χόρτον ὄντα σήμερον καὶ αὔριον εἰς κλίβανον βαλλόμενον ὁ θεὸς οὕτως ἀμφιέζει, πόσῳ μᾶλλον ὑμᾶς, ὀλιγόπιστοι.
		% 29 καὶ ὑμεῖς μὴ ζητεῖτε τί φάγητε καὶ τί πίητε, καὶ μὴ μετεωρίζεσθε,
		% 30 ταῦτα γὰρ πάντα τὰ ἔθνη τοῦ κόσμου ἐπιζητοῦσιν, ὑμῶν δὲ ὁ πατὴρ οἶδεν ὅτι χρῄζετε τούτων.
		% 31 πλὴν ζητεῖτε τὴν βασιλείαν αὐτοῦ, καὶ ταῦτα προστεθήσεται ὑμῖν.
		% }
		
		% \gr{
		% 32 Μὴ φοβοῦ, τὸ μικρὸν ποίμνιον, ὅτι εὐδόκησεν ὁ πατὴρ ὑμῶν δοῦναι ὑμῖν τὴν βασιλείαν.
		% 33 πωλήσατε τὰ ὑπάρχοντα ὑμῶν καὶ δότε ἐλεημοσύνην· ποιήσατε ἑαυτοῖς βαλλάντια μὴ παλαιούμενα, θησαυρὸν ἀνέκλειπτον}%
			% \footnote{%
				% \label{LCHRIST-054-anekleipton}%
				% This is the only use of this word in the NT. BDAG 76b (PDF 79) gives it as ``unfaiing, inexhaustible,'' whereas CGL 120a gives the first definition as ``not coming to an end, \textit{unceasing, perpetual},'' used especially of wars and acts of violence; the next definition is for ``of a treasure in heaven,'' and gives ``inexhaustible.'' The KJV gives a relative clause, ``that faileth not.'' I like ``inexhaustible'' here. The treasure you are storing up is abundant; it is overflowing; it will never run out. It goes toward the fountain of living waters in John 4.
				% }
		% \gr{ἐν τοῖς οὐρανοῖς, ὅπου κλέπτης οὐκ ἐγγίζει οὐδὲ σὴς διαφθείρει·
		% 34 ὅπου γάρ ἐστιν ὁ θησαυρὸς ὑμῶν, ἐκεῖ καὶ ἡ καρδία ὑμῶν ἔσται.
		% }
		
% % -----------------------------------------------------------
% \section{55 | Ye Shall Be Judged with What Judgment ye Judge}
% % -----------------------------------------------------------
	% \label{LCHRIST-055}
	
	% % ----------------------
	% \subsection{Matthew 7:1--5}
	% % ----------------------
		% \label{LCHRIST-055-MT}
		
		% \gr{
		% 1 Μὴ κρίνετε, ἵνα μὴ κριθῆτε·
		% 2 ἐν ᾧ γὰρ κρίματι κρίνετε κριθήσεσθε, καὶ ἐν ᾧ μέτρῳ μετρεῖτε μετρηθήσεται ὑμῖν.
		% 3 τί δὲ βλέπεις τὸ κάρφος τὸ ἐν τῷ ὀφθαλμῷ τοῦ ἀδελφοῦ σου, τὴν δὲ ἐν τῷ σῷ ὀφθαλμῷ δοκὸν οὐ κατανοεῖς;
		% 4 ἢ πῶς ἐρεῖς τῷ ἀδελφῷ σου· Ἄφες ἐκβάλω τὸ κάρφος ἐκ τοῦ ὀφθαλμοῦ σου, καὶ ἰδοὺ ἡ δοκὸς ἐν τῷ ὀφθαλμῷ σοῦ;
		% 5 ὑποκριτά, ἔκβαλε πρῶτον ἐκ τοῦ ὀφθαλμοῦ σοῦ τὴν δοκόν, καὶ τότε διαβλέψεις ἐκβαλεῖν τὸ κάρφος ἐκ τοῦ ὀφθαλμοῦ τοῦ ἀδελφοῦ σου.}%
			% \footnote{%
				% \label{LCHRIST-055-beam}%
				% There's something about the beam in our eye which obstructs our vision of other people---we cannot see them clearly. Notice the difference between the verbs---\gr{βλέπω} first, when we have the beam, and then later \gr{διαβλέπω}, distinguishing the way we are able to see without the beam from the way we could see before that. \hyperref[1CORINTHIANS-13-12]{1 Corinthians 13:12} is another place that we see this altered or unclear perception which will change later---that verse also uses \gr{βλέπω}, but that word is very common so the similarity there isn't that remarkable.
				% }
		
	% % ----------------------
	% \subsection{Luke 6:37--42}
	% % ----------------------
		% \label{LCHRIST-055-L}
		
		% \gr{
		% 37 Καὶ μὴ κρίνετε, καὶ οὐ μὴ κριθῆτε· καὶ μὴ καταδικάζετε, καὶ οὐ μὴ καταδικασθῆτε. ἀπολύετε, καὶ ἀπολυθήσεσθε·
		% 38 δίδοτε, καὶ δοθήσεται ὑμῖν· μέτρον καλὸν πεπιεσμένον σεσαλευμένον ὑπερεκχυννόμενον δώσουσιν εἰς τὸν κόλπον ὑμῶν· ᾧ γὰρ μέτρῳ μετρεῖτε ἀντιμετρηθήσεται ὑμῖν.
		% }
		
		% \gr{
		% 39 Εἶπεν δὲ καὶ παραβολὴν αὐτοῖς· Μήτι δύναται τυφλὸς τυφλὸν ὁδηγεῖν; οὐχὶ ἀμφότεροι εἰς βόθυνον ἐμπεσοῦνται;
		% 40 οὐκ ἔστιν μαθητὴς ὑπὲρ τὸν διδάσκαλον, κατηρτισμένος δὲ πᾶς ἔσται ὡς ὁ διδάσκαλος αὐτοῦ.
		% 41 τί δὲ βλέπεις τὸ κάρφος τὸ ἐν τῷ ὀφθαλμῷ τοῦ ἀδελφοῦ σου, τὴν δὲ δοκὸν τὴν ἐν τῷ ἰδίῳ ὀφθαλμῷ οὐ κατανοεῖς;
		% 42 πῶς δύνασαι λέγειν τῷ ἀδελφῷ σου· Ἀδελφέ, ἄφες ἐκβάλω τὸ κάρφος τὸ ἐν τῷ ὀφθαλμῷ σου, αὐτὸς τὴν ἐν τῷ ὀφθαλμῷ σοῦ δοκὸν οὐ βλέπων; ὑποκριτά, ἔκβαλε πρῶτον τὴν δοκὸν ἐκ τοῦ ὀφθαλμοῦ σοῦ, καὶ τότε διαβλέψεις τὸ κάρφος τὸ ἐν τῷ ὀφθαλμῷ τοῦ ἀδελφοῦ σου ἐκβαλεῖν.
		% }
		
% % -----------------------------------------------------------
% \section{56 | Cast Not Your Pearls before Swine}
% % -----------------------------------------------------------
	% \label{LCHRIST-056}
	
	% % ----------------------
	% \subsection{Matthew 7:6}
	% % ----------------------
		% \label{LCHRIST-056-MT}
		
		% \gr{
		% 6 Μὴ δῶτε τὸ ἅγιον τοῖς κυσίν, μηδὲ βάλητε τοὺς μαργαρίτας ὑμῶν ἔμπροσθεν τῶν χοίρων, μήποτε καταπατήσουσιν αὐτοὺς ἐν τοῖς ποσὶν αὐτῶν καὶ στραφέντες ῥήξωσιν ὑμᾶς.
		% }
	
% % -----------------------------------------------------------
% \section{57 | Ask, Seek, Knock}
% % -----------------------------------------------------------
	% \label{LCHRIST-057}
	
	% % ----------------------
	% \subsection{Matthew 7:7--12}
	% % ----------------------
		% \label{LCHRIST-057-MT}
		
		% \gr{
		% 7 Αἰτεῖτε, καὶ δοθήσεται ὑμῖν· ζητεῖτε, καὶ εὑρήσετε· κρούετε, καὶ ἀνοιγήσεται ὑμῖν.
		% 8 πᾶς γὰρ ὁ αἰτῶν λαμβάνει καὶ ὁ ζητῶν εὑρίσκει καὶ τῷ κρούοντι ἀνοιγήσεται.
		% 9 ἢ τίς ἐστιν ἐξ ὑμῶν ἄνθρωπος, ὃν αἰτήσει ὁ υἱὸς αὐτοῦ ἄρτον— μὴ λίθον ἐπιδώσει αὐτῷ;
		% 10 ἢ καὶ ἰχθὺν αἰτήσει— μὴ ὄφιν ἐπιδώσει αὐτῷ;
		% 11 εἰ οὖν ὑμεῖς πονηροὶ ὄντες οἴδατε δόματα ἀγαθὰ διδόναι τοῖς τέκνοις ὑμῶν, πόσῳ μᾶλλον ὁ πατὴρ ὑμῶν ὁ ἐν τοῖς οὐρανοῖς δώσει ἀγαθὰ τοῖς αἰτοῦσιν αὐτόν.
		% }
		
		% \gr{
		% 12 Πάντα οὖν ὅσα ἐὰν θέλητε ἵνα ποιῶσιν ὑμῖν οἱ ἄνθρωποι, οὕτως καὶ ὑμεῖς ποιεῖτε αὐτοῖς· οὗτος γάρ ἐστιν ὁ νόμος καὶ οἱ προφῆται.
		% }
		
	% % ----------------------
	% \subsection{Luke 11:9--13}
	% % ----------------------
		% \label{LCHRIST-057-L}
		
		% \gr{
		% 9 Κἀγὼ ὑμῖν λέγω, αἰτεῖτε, καὶ δοθήσεται ὑμῖν· ζητεῖτε, καὶ εὑρήσετε· κρούετε, καὶ ἀνοιγήσεται ὑμῖν·
		% 10 πᾶς γὰρ ὁ αἰτῶν λαμβάνει, καὶ ὁ ζητῶν εὑρίσκει, καὶ τῷ κρούοντι ἀνοιγήσεται.
		% 11 τίνα δὲ ἐξ ὑμῶν τὸν πατέρα αἰτήσει ὁ υἱὸς ἰχθύν, καὶ ἀντὶ ἰχθύος ὄφιν αὐτῷ ἐπιδώσει;
		% 12 ἢ καὶ αἰτήσει ᾠόν, ἐπιδώσει αὐτῷ σκορπίον;
		% 13 εἰ οὖν ὑμεῖς πονηροὶ ὑπάρχοντες οἴδατε δόματα ἀγαθὰ διδόναι τοῖς τέκνοις ὑμῶν, πόσῳ μᾶλλον ὁ πατὴρ ὁ ἐξ οὐρανοῦ δώσει πνεῦμα ἅγιον τοῖς αἰτοῦσιν αὐτόν.
		% }
		
% % -----------------------------------------------------------
% \section{58 | The Golden Rule}
% % -----------------------------------------------------------
	% \label{LCHRIST-058}
	
	% % ----------------------
	% \subsection{Matthew 7:12}
	% % ----------------------
		% \label{LCHRIST-058-MT}
		
		% \gr{
		% 12 Πάντα οὖν ὅσα ἐὰν θέλητε ἵνα ποιῶσιν ὑμῖν οἱ ἄνθρωποι, οὕτως καὶ ὑμεῖς ποιεῖτε αὐτοῖς· οὗτος γάρ ἐστιν ὁ νόμος καὶ οἱ προφῆται.
		% }
		
	% % ----------------------
	% \subsection{Luke 6:31}
	% % ----------------------
		% \label{LCHRIST-058-L}
		
		% \gr{
		% 31 καὶ καθὼς θέλετε ἵνα ποιῶσιν ὑμῖν οἱ ἄνθρωποι, ποιεῖτε αὐτοῖς ὁμοίως.
		% }
		
% % -----------------------------------------------------------
% \section{59 | Strait is the Gate, Narrow the Way}
% % -----------------------------------------------------------
	% \label{LCHRIST-059}
	
	% % ----------------------
	% \subsection{Matthew 7:13--14, 21--23}
	% % ----------------------
		% \label{LCHRIST-059-MT}
		
		% \gr{
		% 13 Εἰσέλθατε διὰ τῆς στενῆς πύλης· ὅτι πλατεῖα ἡ πύλη καὶ εὐρύχωρος ἡ ὁδὸς ἡ ἀπάγουσα εἰς τὴν ἀπώλειαν, καὶ πολλοί εἰσιν οἱ εἰσερχόμενοι δι᾽ αὐτῆς·
		% 14 ὅτι στενὴ ἡ πύλη καὶ τεθλιμμένη ἡ ὁδὸς ἡ ἀπάγουσα εἰς τὴν ζωήν, καὶ ὀλίγοι εἰσὶν οἱ εὑρίσκοντες αὐτήν.
		% }
		
		% \gr{
		% 21 Οὐ πᾶς ὁ λέγων μοι· Κύριε κύριε εἰσελεύσεται εἰς τὴν βασιλείαν τῶν οὐρανῶν, ἀλλ᾽ ὁ ποιῶν τὸ θέλημα τοῦ πατρός μου τοῦ ἐν τοῖς οὐρανοῖς.
		% 22 πολλοὶ ἐροῦσίν μοι ἐν ἐκείνῃ τῇ ἡμέρᾳ· Κύριε κύριε, οὐ τῷ σῷ ὀνόματι ἐπροφητεύσαμεν, καὶ τῷ σῷ ὀνόματι δαιμόνια ἐξεβάλομεν, καὶ τῷ σῷ ὀνόματι δυνάμεις πολλὰς ἐποιήσαμεν;
		% 23 καὶ τότε ὁμολογήσω αὐτοῖς ὅτι Οὐδέποτε ἔγνων ὑμᾶς· ἀποχωρεῖτε ἀπ᾽ ἐμοῦ οἱ ἐργαζόμενοι τὴν ἀνομίαν.
		% }
		
	% % ----------------------
	% \subsection{Luke 13:22--23}
	% % ----------------------
		% \label{LCHRIST-059-L}
		
		% \gr{
		% 22 Καὶ διεπορεύετο κατὰ πόλεις καὶ κώμας διδάσκων καὶ πορείαν ποιούμενος εἰς Ἱεροσόλυμα.
		% 23 εἶπεν δέ τις αὐτῷ· Κύριε, εἰ ὀλίγοι οἱ σῳζόμενοι; ὁ δὲ εἶπεν πρὸς αὐτούς·
		% }
		
% % -----------------------------------------------------------
% \section{60 | A Tree Bears Only One Kind of Fruit}
% % -----------------------------------------------------------
	% \label{LCHRIST-060}
	
	% % ----------------------
	% \subsection{Matthew 7:15--20; 12:34--35}
	% % ----------------------
		% \label{LCHRIST-060-MT}
		
		% \gr{
		% 15 Προσέχετε ἀπὸ τῶν ψευδοπροφητῶν, οἵτινες ἔρχονται πρὸς ὑμᾶς ἐν ἐνδύμασι προβάτων ἔσωθεν δέ εἰσιν λύκοι ἅρπαγες.
		% 16 ἀπὸ τῶν καρπῶν αὐτῶν ἐπιγνώσεσθε αὐτούς. μήτι συλλέγουσιν ἀπὸ ἀκανθῶν σταφυλὰς ἢ ἀπὸ τριβόλων σῦκα;
		% 17 οὕτως πᾶν δένδρον ἀγαθὸν καρποὺς καλοὺς ποιεῖ, τὸ δὲ σαπρὸν δένδρον καρποὺς πονηροὺς ποιεῖ·
		% 18 οὐ δύναται δένδρον ἀγαθὸν καρποὺς πονηροὺς ποιεῖν, οὐδὲ δένδρον σαπρὸν καρποὺς καλοὺς ποιεῖν.
		% 19 πᾶν δένδρον μὴ ποιοῦν καρπὸν καλὸν ἐκκόπτεται καὶ εἰς πῦρ βάλλεται.
		% 20 ἄρα γε ἀπὸ τῶν καρπῶν αὐτῶν ἐπιγνώσεσθε αὐτούς.
		% }
		
		% \gr{
		% 34 γεννήματα ἐχιδνῶν, πῶς δύνασθε ἀγαθὰ λαλεῖν πονηροὶ ὄντες; ἐκ γὰρ τοῦ περισσεύματος τῆς καρδίας τὸ στόμα λαλεῖ.
		% 35 ὁ ἀγαθὸς ἄνθρωπος ἐκ τοῦ ἀγαθοῦ θησαυροῦ ἐκβάλλει ἀγαθά, καὶ ὁ πονηρὸς ἄνθρωπος ἐκ τοῦ πονηροῦ θησαυροῦ ἐκβάλλει πονηρά.
		% }
		
	% % ----------------------
	% \subsection{Luke 6:43--45}
	% % ----------------------
		% \label{LCHRIST-060-L}
		
		% \gr{
		% 43 Οὐ γάρ ἐστιν δένδρον καλὸν ποιοῦν καρπὸν σαπρόν, οὐδὲ πάλιν δένδρον σαπρὸν ποιοῦν καρπὸν καλόν.
		% 44 ἕκαστον γὰρ δένδρον ἐκ τοῦ ἰδίου καρποῦ γινώσκεται· οὐ γὰρ ἐξ ἀκανθῶν συλλέγουσιν σῦκα, οὐδὲ ἐκ βάτου σταφυλὴν τρυγῶσιν.
		% 45 ὁ ἀγαθὸς ἄνθρωπος ἐκ τοῦ ἀγαθοῦ θησαυροῦ τῆς καρδίας αὐτοῦ προφέρει τὸ ἀγαθόν, καὶ ὁ πονηρὸς ἐκ τοῦ πονηροῦ προφέρει τὸ πονηρόν· ἐκ γὰρ περισσεύματος καρδίας λαλεῖ τὸ στόμα αὐτοῦ.
		% }
		
% % -----------------------------------------------------------
% \section{61 | I Never Knew You}
% % -----------------------------------------------------------
	% \label{LCHRIST-061}
	
	% % ----------------------
	% \subsection{Matthew 7:21--23}
	% % ----------------------
		% \label{LCHRIST-061-MT}
		
		% \gr{
		% 21 Οὐ πᾶς ὁ λέγων μοι· Κύριε κύριε εἰσελεύσεται εἰς τὴν βασιλείαν τῶν οὐρανῶν, ἀλλ᾽ ὁ ποιῶν τὸ θέλημα τοῦ πατρός μου τοῦ ἐν τοῖς οὐρανοῖς.
		% 22 πολλοὶ ἐροῦσίν μοι ἐν ἐκείνῃ τῇ ἡμέρᾳ· Κύριε κύριε, οὐ τῷ σῷ ὀνόματι ἐπροφητεύσαμεν, καὶ τῷ σῷ ὀνόματι δαιμόνια ἐξεβάλομεν, καὶ τῷ σῷ ὀνόματι δυνάμεις πολλὰς ἐποιήσαμεν;
		% 23 καὶ τότε ὁμολογήσω αὐτοῖς ὅτι Οὐδέποτε ἔγνων ὑμᾶς· ἀποχωρεῖτε ἀπ᾽ ἐμοῦ οἱ ἐργαζόμενοι τὴν ἀνομίαν.
		% }
		
	% % ----------------------
	% \subsection{Luke 13:25--27}
	% % ----------------------
		% \label{LCHRIST-061-L}
		
		% \gr{
		% 25 ἀφ᾽ οὗ ἂν ἐγερθῇ ὁ οἰκοδεσπότης καὶ ἀποκλείσῃ τὴν θύραν, καὶ ἄρξησθε ἔξω ἑστάναι καὶ κρούειν τὴν θύραν λέγοντες· Κύριε, ἄνοιξον ἡμῖν· καὶ ἀποκριθεὶς ἐρεῖ ὑμῖν· Οὐκ οἶδα ὑμᾶς πόθεν ἐστέ.
		% 26 τότε ἄρξεσθε λέγειν· Ἐφάγομεν ἐνώπιόν σου καὶ ἐπίομεν, καὶ ἐν ταῖς πλατείαις ἡμῶν ἐδίδαξας·
		% 27 καὶ ἐρεῖ λέγων ὑμῖν· Οὐκ οἶδα πόθεν ἐστέ· ἀπόστητε ἀπ᾽ ἐμοῦ, πάντες ἐργάται ἀδικίας.
		% }
		
% % -----------------------------------------------------------
% \section{62 | Wise Man Builds House on Rock}
% % -----------------------------------------------------------
	% \label{LCHRIST-062}

	% % ----------------------
	% \subsection{Matthew 7:24--28}
	% % ----------------------
		% \label{LCHRIST-062-MT}
		
		% \gr{
		% 24 Πᾶς οὖν ὅστις ἀκούει μου τοὺς λόγους τούτους καὶ ποιεῖ αὐτούς, ὁμοιωθήσεται ἀνδρὶ φρονίμῳ, ὅστις ᾠκοδόμησεν αὐτοῦ τὴν οἰκίαν ἐπὶ τὴν πέτραν.
		% 25 καὶ κατέβη ἡ βροχὴ καὶ ἦλθον οἱ ποταμοὶ καὶ ἔπνευσαν οἱ ἄνεμοι καὶ προσέπεσαν τῇ οἰκίᾳ ἐκείνῃ, καὶ οὐκ ἔπεσεν, τεθεμελίωτο γὰρ ἐπὶ τὴν πέτραν.
		% 26 καὶ πᾶς ὁ ἀκούων μου τοὺς λόγους τούτους καὶ μὴ ποιῶν αὐτοὺς ὁμοιωθήσεται ἀνδρὶ μωρῷ, ὅστις ᾠκοδόμησεν αὐτοῦ τὴν οἰκίαν ἐπὶ τὴν ἄμμον.
		% 27 καὶ κατέβη ἡ βροχὴ καὶ ἦλθον οἱ ποταμοὶ καὶ ἔπνευσαν οἱ ἄνεμοι καὶ προσέκοψαν τῇ οἰκίᾳ ἐκείνῃ, καὶ ἔπεσεν, καὶ ἦν ἡ πτῶσις αὐτῆς μεγάλη.
		% }%
			% \footnote{%
				% \label{LCHRIST-062-MT-sand}%
				% Why would someone build a house on sand in the first place? The Luke account is helpful here, saying that the second house is built \gr{ἐπὶ τὴν γῆν χωρὶς θεμελίου}---on the earth or ground \textit{without a foundation}. I think the ``without a foundation'' is the key part here, and so the question would be, why would someone want to build a house without a foundation? Why would they choose an area, like one covered in sand, on which it would not be possible to put a foundation? My first thought is that it would just be easier to build your house without a foundation; it would take less time and less effort. Maybe you also don't know how? Maybe you find the ground to be more comfortable---you prefer a sandy floor instead of a harder one? I'm not totally sure. See \hyperref[HELAMAN-5-12]{Helaman 5:12}---\textit{Christ} is the rock on which \textit{we} build the foundation. What is the foundation? Our habits? Not sure yet about that one either, though see sentences 7--9 of Openshaw's blessing, and compare to Helaman 5:12---together they suggest habits are a good way of thinking about it.
				% }
		
		% \gr{
		% 28 Καὶ ἐγένετο ὅτε ἐτέλεσεν ὁ Ἰησοῦς τοὺς λόγους τούτους, ἐξεπλήσσοντο οἱ ὄχλοι ἐπὶ τῇ διδαχῇ αὐτοῦ·
		% }
		
	% % ----------------------
	% \subsection{Luke 6:46--49}
	% % ----------------------
		% \label{LCHRIST-062-L}
		
		% \gr{
		% 46 Τί δέ με καλεῖτε· Κύριε κύριε, καὶ οὐ ποιεῖτε ἃ λέγω;
		% 47 πᾶς ὁ ἐρχόμενος πρός με καὶ ἀκούων μου τῶν λόγων καὶ ποιῶν αὐτούς, ὑποδείξω ὑμῖν τίνι ἐστὶν ὅμοιος·
		% 48 ὅμοιός ἐστιν ἀνθρώπῳ οἰκοδομοῦντι οἰκίαν ὃς ἔσκαψεν καὶ ἐβάθυνεν καὶ ἔθηκεν θεμέλιον ἐπὶ τὴν πέτραν· πλημμύρης δὲ γενομένης προσέρηξεν ὁ ποταμὸς τῇ οἰκίᾳ ἐκείνῃ, καὶ οὐκ ἴσχυσεν σαλεῦσαι αὐτὴν διὰ τὸ καλῶς οἰκοδομῆσθαι αὐτήν.
		% 49 ὁ δὲ ἀκούσας καὶ μὴ ποιήσας ὅμοιός ἐστιν ἀνθρώπῳ οἰκοδομήσαντι οἰκίαν ἐπὶ τὴν γῆν χωρὶς θεμελίου, ᾗ προσέρηξεν ὁ ποταμός, καὶ εὐθὺς συνέπεσεν, καὶ ἐγένετο τὸ ῥῆγμα τῆς οἰκίας ἐκείνης μέγα.
		% }
		
% % -----------------------------------------------------------
% \section{63 | Jesus Taught as One with Authority}
% % -----------------------------------------------------------
	% \label{LCHRIST-063}
	
	% % ----------------------
	% \subsection{Matthew 7:28--29}
	% % ----------------------
		% \label{LCHRIST-063-MT}
		
		% \gr{
		% 28 Καὶ ἐγένετο ὅτε ἐτέλεσεν ὁ Ἰησοῦς τοὺς λόγους τούτους, ἐξεπλήσσοντο οἱ ὄχλοι ἐπὶ τῇ διδαχῇ αὐτοῦ·
		% 29 ἦν γὰρ διδάσκων αὐτοὺς ὡς} \hyperref[EXOUSIA]{\grl{ἐξουσίαν}} \gr{ἔχων καὶ οὐχ ὡς οἱ γραμματεῖς αὐτῶν.}
		
% % -----------------------------------------------------------
% \section{64 | Samaritan Woman at the Well}
% % -----------------------------------------------------------
	% \label{LCHRIST-064}
	
	% % ----------------------
	% \subsection{John 4:1--42}
	% % ----------------------
		% \label{LCHRIST-064-J}
		
		% \gr{
		% 1 Ὡς οὖν ἔγνω ὁ Ἰησοῦς ὅτι ἤκουσαν οἱ Φαρισαῖοι ὅτι Ἰησοῦς πλείονας μαθητὰς ποιεῖ καὶ βαπτίζει ἢ Ἰωάννης—
		% 2 καίτοιγε Ἰησοῦς αὐτὸς οὐκ ἐβάπτιζεν ἀλλ᾽ οἱ μαθηταὶ αὐτοῦ—
		% 3 ἀφῆκεν τὴν Ἰουδαίαν καὶ ἀπῆλθεν πάλιν εἰς τὴν Γαλιλαίαν.
		% 4 ἔδει δὲ αὐτὸν διέρχεσθαι διὰ τῆς Σαμαρείας.
		% 5 ἔρχεται οὖν εἰς πόλιν τῆς Σαμαρείας λεγομένην Συχὰρ πλησίον τοῦ χωρίου ὃ ἔδωκεν Ἰακὼβ τῷ Ἰωσὴφ τῷ υἱῷ αὐτοῦ·
		% 6 ἦν δὲ ἐκεῖ πηγὴ τοῦ Ἰακώβ. ὁ οὖν Ἰησοῦς κεκοπιακὼς ἐκ τῆς ὁδοιπορίας ἐκαθέζετο οὕτως ἐπὶ τῇ πηγῇ· ὥρα ἦν ὡς ἕκτη.
		% }
		
		% \gr{
		% 7 Ἔρχεται γυνὴ ἐκ τῆς Σαμαρείας ἀντλῆσαι ὕδωρ. λέγει αὐτῇ ὁ Ἰησοῦς· Δός μοι πεῖν·
		% 8 οἱ γὰρ μαθηταὶ αὐτοῦ ἀπεληλύθεισαν εἰς τὴν πόλιν, ἵνα τροφὰς ἀγοράσωσιν.
		% 9 λέγει οὖν αὐτῷ ἡ γυνὴ ἡ Σαμαρῖτις· Πῶς σὺ Ἰουδαῖος ὢν παρ᾽ ἐμοῦ πεῖν αἰτεῖς γυναικὸς Σαμαρίτιδος οὔσης; οὐ γὰρ συγχρῶνται Ἰουδαῖοι Σαμαρίταις.
		% 10 ἀπεκρίθη Ἰησοῦς καὶ εἶπεν αὐτῇ· Εἰ ᾔδεις τὴν δωρεὰν τοῦ θεοῦ καὶ τίς ἐστιν ὁ λέγων σοι· Δός μοι πεῖν, σὺ ἂν ᾔτησας αὐτὸν καὶ ἔδωκεν ἄν σοι ὕδωρ ζῶν.
		% 11 λέγει αὐτῷ ἡ γυνή· Κύριε, οὔτε ἄντλημα ἔχεις καὶ τὸ φρέαρ ἐστὶν βαθύ· πόθεν οὖν ἔχεις τὸ ὕδωρ τὸ ζῶν;
		% 12 μὴ σὺ μείζων εἶ τοῦ πατρὸς ἡμῶν Ἰακώβ, ὃς ἔδωκεν ἡμῖν τὸ φρέαρ καὶ αὐτὸς ἐξ αὐτοῦ ἔπιεν καὶ οἱ υἱοὶ αὐτοῦ καὶ τὰ θρέμματα αὐτοῦ;
		% 13 ἀπεκρίθη Ἰησοῦς καὶ εἶπεν αὐτῇ· Πᾶς ὁ πίνων ἐκ τοῦ ὕδατος τούτου διψήσει πάλιν·
		% 14 ὃς δ᾽ ἂν πίῃ ἐκ τοῦ ὕδατος οὗ ἐγὼ δώσω αὐτῷ, οὐ μὴ διψήσει εἰς τὸν αἰῶνα, ἀλλὰ τὸ ὕδωρ ὃ δώσω αὐτῷ γενήσεται ἐν αὐτῷ πηγὴ ὕδατος ἁλλομένου
		% εἰς ζωὴν αἰώνιον.
		% 15 λέγει πρὸς αὐτὸν ἡ γυνή· Κύριε, δός μοι τοῦτο τὸ ὕδωρ, ἵνα μὴ διψῶ μηδὲ διέρχωμαι ἐνθάδε ἀντλεῖν.
		% }
		
		% \gr{
		% 16 Λέγει αὐτῇ· Ὕπαγε φώνησον τὸν ἄνδρα σου καὶ ἐλθὲ ἐνθάδε.
		% 17 ἀπεκρίθη ἡ γυνὴ καὶ εἶπεν αὐτῷ· Οὐκ ἔχω ἄνδρα. λέγει αὐτῇ ὁ Ἰησοῦς· Καλῶς εἶπας ὅτι Ἄνδρα οὐκ ἔχω·
		% 18 πέντε γὰρ ἄνδρας ἔσχες, καὶ νῦν ὃν ἔχεις οὐκ ἔστιν σου ἀνήρ· τοῦτο ἀληθὲς εἴρηκας.
		% 19 λέγει αὐτῷ ἡ γυνή· Κύριε, θεωρῶ ὅτι προφήτης εἶ σύ.
		% 20 οἱ πατέρες ἡμῶν ἐν τῷ ὄρει τούτῳ προσεκύνησαν· καὶ ὑμεῖς λέγετε ὅτι ἐν Ἱεροσολύμοις ἐστὶν ὁ τόπος ὅπου προσκυνεῖν δεῖ.
		% 21 λέγει αὐτῇ ὁ Ἰησοῦς· Πίστευέ μοι, γύναι, ὅτι ἔρχεται ὥρα ὅτε οὔτε ἐν τῷ ὄρει τούτῳ οὔτε ἐν Ἱεροσολύμοις προσκυνήσετε τῷ πατρί.
		% 22 ὑμεῖς προσκυνεῖτε ὃ οὐκ οἴδατε, ἡμεῖς προσκυνοῦμεν ὃ οἴδαμεν,}
			% \footnote{%
				% \label{LCHRIST-064-know}%
				% See \hyperref[DC93-19]{DC 93:19}.
				% }
		% \gr{ὅτι ἡ σωτηρία ἐκ τῶν Ἰουδαίων ἐστίν·
		% 23 ἀλλὰ ἔρχεται ὥρα καὶ νῦν ἐστιν, ὅτε οἱ ἀληθινοὶ προσκυνηταὶ προσκυνήσουσιν τῷ πατρὶ ἐν πνεύματι καὶ ἀληθείᾳ, καὶ γὰρ ὁ πατὴρ τοιούτους ζητεῖ τοὺς προσκυνοῦντας αὐτόν·
		% 24 πνεῦμα ὁ θεός, καὶ τοὺς προσκυνοῦντας αὐτὸν ἐν πνεύματι καὶ ἀληθείᾳ δεῖ προσκυνεῖν.
		% 25 λέγει αὐτῷ ἡ γυνή· Οἶδα ὅτι Μεσσίας ἔρχεται, ὁ λεγόμενος χριστός· ὅταν ἔλθῃ ἐκεῖνος, ἀναγγελεῖ ἡμῖν ἅπαντα.
		% 26 λέγει αὐτῇ ὁ Ἰησοῦς· Ἐγώ εἰμι, ὁ λαλῶν σοι.
		% }
		
		% \gr{
		% 27 Καὶ ἐπὶ τούτῳ ἦλθαν οἱ μαθηταὶ αὐτοῦ, καὶ ἐθαύμαζον ὅτι μετὰ γυναικὸς ἐλάλει· οὐδεὶς μέντοι εἶπεν· Τί ζητεῖς; ἢ τί λαλεῖς μετ᾽ αὐτῆς;
		% 28 ἀφῆκεν οὖν τὴν ὑδρίαν αὐτῆς ἡ γυνὴ καὶ ἀπῆλθεν εἰς τὴν πόλιν καὶ λέγει τοῖς ἀνθρώποις·
		% 29 Δεῦτε ἴδετε ἄνθρωπον ὃς εἶπέ μοι πάντα ὅσα ἐποίησα· μήτι οὗτός ἐστιν ὁ χριστός;
		% 30 ἐξῆλθον ἐκ τῆς πόλεως καὶ ἤρχοντο πρὸς αὐτόν.
		% }
		
		% \gr{
		% 31 Ἐν τῷ μεταξὺ ἠρώτων αὐτὸν οἱ μαθηταὶ λέγοντες· Ῥαββί, φάγε.
		% 32 ὁ δὲ εἶπεν αὐτοῖς· Ἐγὼ βρῶσιν ἔχω φαγεῖν ἣν ὑμεῖς οὐκ οἴδατε.
		% 33 ἔλεγον οὖν οἱ μαθηταὶ πρὸς ἀλλήλους· Μή τις ἤνεγκεν αὐτῷ φαγεῖν;
		% 34 λέγει αὐτοῖς ὁ Ἰησοῦς· Ἐμὸν βρῶμά ἐστιν ἵνα ποιήσω τὸ θέλημα τοῦ πέμψαντός με καὶ τελειώσω αὐτοῦ τὸ ἔργον.
		% 35 οὐχ ὑμεῖς λέγετε ὅτι Ἔτι τετράμηνός ἐστιν καὶ ὁ θερισμὸς ἔρχεται; ἰδοὺ λέγω ὑμῖν, ἐπάρατε τοὺς ὀφθαλμοὺς ὑμῶν καὶ θεάσασθε τὰς χώρας ὅτι λευκαί εἰσιν πρὸς θερισμόν· ἤδη
		% 36 ὁ θερίζων μισθὸν λαμβάνει καὶ συνάγει καρπὸν εἰς ζωὴν αἰώνιον, ἵνα ὁ σπείρων ὁμοῦ χαίρῃ καὶ ὁ θερίζων.
		% 37 ἐν γὰρ τούτῳ ὁ λόγος ἐστὶν ἀληθινὸς ὅτι Ἄλλος ἐστὶν ὁ σπείρων καὶ ἄλλος ὁ θερίζων·
		% 38 ἐγὼ ἀπέστειλα ὑμᾶς θερίζειν ὃ οὐχ ὑμεῖς κεκοπιάκατε· ἄλλοι κεκοπιάκασιν, καὶ ὑμεῖς εἰς τὸν κόπον αὐτῶν εἰσεληλύθατε.
		% }
		
		% \gr{
		% 39 Ἐκ δὲ τῆς πόλεως ἐκείνης πολλοὶ ἐπίστευσαν εἰς αὐτὸν τῶν Σαμαριτῶν διὰ τὸν λόγον τῆς γυναικὸς μαρτυρούσης ὅτι Εἶπέν μοι πάντα ἃ ἐποίησα.
		% 40 ὡς οὖν ἦλθον πρὸς αὐτὸν οἱ Σαμαρῖται, ἠρώτων αὐτὸν μεῖναι παρ᾽ αὐτοῖς· καὶ ἔμεινεν ἐκεῖ δύο ἡμέρας.
		% 41 καὶ πολλῷ πλείους ἐπίστευσαν διὰ τὸν λόγον αὐτοῦ,
		% 42 τῇ τε γυναικὶ ἔλεγον ὅτι Οὐκέτι διὰ τὴν σὴν λαλιὰν πιστεύομεν· αὐτοὶ γὰρ ἀκηκόαμεν, καὶ οἴδαμεν ὅτι οὗτός ἐστιν ἀληθῶς ὁ σωτὴρ τοῦ κόσμου.
		% }
		
% % -----------------------------------------------------------
% \section{65 | Jesus Heals Faithful Leper}
% % -----------------------------------------------------------
	% \label{LCHRIST-065}
	
	% % ----------------------
	% \subsection{Matthew 8:1--4}
	% % ----------------------
		% \label{LCHRIST-065-MT}
		
		% \gr{
		% 1 Καταβάντος δὲ αὐτοῦ ἀπὸ τοῦ ὄρους ἠκολούθησαν αὐτῷ ὄχλοι πολλοί.
		% 2 καὶ ἰδοὺ λεπρὸς προσελθὼν προσεκύνει αὐτῷ λέγων· Κύριε, ἐὰν θέλῃς δύνασαί με καθαρίσαι.
		% 3 καὶ ἐκτείνας τὴν χεῖρα ἥψατο αὐτοῦ λέγων· Θέλω, καθαρίσθητι· καὶ εὐθέως ἐκαθαρίσθη αὐτοῦ ἡ λέπρα.
		% 4 καὶ λέγει αὐτῷ ὁ Ἰησοῦς· Ὅρα μηδενὶ εἴπῃς, ἀλλὰ ὕπαγε σεαυτὸν δεῖξον τῷ ἱερεῖ, καὶ προσένεγκον τὸ δῶρον ὃ προσέταξεν Μωϋσῆς εἰς μαρτύριον αὐτοῖς.
		% }
		
	% % ----------------------
	% \subsection{Mark 1:40--45}
	% % ----------------------
		% \label{LCHRIST-065-MK}
		
		% \gr{
		% 40 Καὶ ἔρχεται πρὸς αὐτὸν λεπρὸς παρακαλῶν αὐτὸν καὶ γονυπετῶν λέγων αὐτῷ ὅτι Ἐὰν θέλῃς δύνασαί με καθαρίσαι.
		% 41 καὶ ὀργισθεὶς ἐκτείνας τὴν χεῖρα αὐτοῦ ἥψατο καὶ λέγει αὐτῷ· Θέλω, καθαρίσθητι·
		% 42 καὶ εὐθὺς ἀπῆλθεν ἀπ᾽ αὐτοῦ ἡ λέπρα, καὶ ἐκαθαρίσθη.
		% 43 καὶ ἐμβριμησάμενος αὐτῷ εὐθὺς ἐξέβαλεν αὐτόν,
		% 44 καὶ λέγει αὐτῷ· Ὅρα μηδενὶ μηδὲν εἴπῃς, ἀλλὰ ὕπαγε σεαυτὸν δεῖξον τῷ ἱερεῖ καὶ προσένεγκε περὶ τοῦ καθαρισμοῦ σου ἃ προσέταξεν Μωϋσῆς εἰς μαρτύριον αὐτοῖς.
		% 45 ὁ δὲ ἐξελθὼν ἤρξατο κηρύσσειν πολλὰ καὶ διαφημίζειν τὸν λόγον, ὥστε μηκέτι αὐτὸν δύνασθαι φανερῶς εἰς πόλιν εἰσελθεῖν, ἀλλὰ ἔξω ἐπ᾽ ἐρήμοις τόποις ἦν· καὶ ἤρχοντο πρὸς αὐτὸν πάντοθεν.
		% }
		
	% % ----------------------
	% \subsection{Luke 5:12--16}
	% % ----------------------
		% \label{LCHRIST-065-L}
		
		% \gr{
		% 12 Καὶ ἐγένετο ἐν τῷ εἶναι αὐτὸν ἐν μιᾷ τῶν πόλεων καὶ ἰδοὺ ἀνὴρ πλήρης λέπρας· καὶ ἰδὼν τὸν Ἰησοῦν πεσὼν ἐπὶ πρόσωπον ἐδεήθη αὐτοῦ λέγων· Κύριε, ἐὰν θέλῃς δύνασαί με καθαρίσαι.
		% 13 καὶ ἐκτείνας τὴν χεῖρα ἥψατο αὐτοῦ εἰπών· Θέλω, καθαρίσθητι· καὶ εὐθέως ἡ λέπρα ἀπῆλθεν ἀπ᾽ αὐτοῦ.
		% 14 καὶ αὐτὸς παρήγγειλεν αὐτῷ μηδενὶ εἰπεῖν, ἀλλὰ ἀπελθὼν δεῖξον σεαυτὸν τῷ ἱερεῖ, καὶ προσένεγκε περὶ τοῦ καθαρισμοῦ σου καθὼς προσέταξεν Μωϋσῆς εἰς μαρτύριον αὐτοῖς.
		% 15 διήρχετο δὲ μᾶλλον ὁ λόγος περὶ αὐτοῦ, καὶ συνήρχοντο ὄχλοι πολλοὶ ἀκούειν καὶ θεραπεύεσθαι ἀπὸ τῶν ἀσθενειῶν αὐτῶν·
		% 16 αὐτὸς δὲ ἦν ὑποχωρῶν ἐν ταῖς ἐρήμοις καὶ προσευχόμενος.
		% }
		
% % -----------------------------------------------------------
% \section{66 | Jesus Heals Centurion's Servant}
% % -----------------------------------------------------------
	% \label{LCHRIST-066}
	
	% % ----------------------
	% \subsection{Matthew 8:5--13}
	% % ----------------------
		% \label{LCHRIST-066-MT}
		
		% \gr{
		% 5 Εἰσελθόντος δὲ αὐτοῦ εἰς Καφαρναοὺμ προσῆλθεν αὐτῷ ἑκατόνταρχος παρακαλῶν αὐτὸν
		% 6 καὶ λέγων· Κύριε, ὁ παῖς μου βέβληται ἐν τῇ οἰκίᾳ παραλυτικός, δεινῶς βασανιζόμενος.
		% 7 καὶ λέγει αὐτῷ· Ἐγὼ ἐλθὼν θεραπεύσω αὐτόν.
		% 8 καὶ ἀποκριθεὶς ὁ ἑκατόνταρχος ἔφη· Κύριε, οὐκ εἰμὶ ἱκανὸς ἵνα μου ὑπὸ τὴν στέγην εἰσέλθῃς· ἀλλὰ μόνον εἰπὲ λόγῳ, καὶ ἰαθήσεται ὁ παῖς μου·
		% 9 καὶ γὰρ ἐγὼ ἄνθρωπός εἰμι ὑπὸ ἐξουσίαν, ἔχων ὑπ᾽ ἐμαυτὸν στρατιώτας, καὶ λέγω τούτῳ· Πορεύθητι, καὶ πορεύεται, καὶ ἄλλῳ· Ἔρχου, καὶ ἔρχεται, καὶ τῷ δούλῳ μου· Ποίησον τοῦτο, καὶ ποιεῖ.
		% 10 ἀκούσας δὲ ὁ Ἰησοῦς ἐθαύμασεν καὶ εἶπεν τοῖς ἀκολουθοῦσιν· Ἀμὴν λέγω ὑμῖν, παρ᾽ οὐδενὶ τοσαύτην πίστιν ἐν τῷ Ἰσραὴλ εὗρον.
		% 11 λέγω δὲ ὑμῖν ὅτι πολλοὶ ἀπὸ ἀνατολῶν καὶ δυσμῶν ἥξουσιν καὶ ἀνακλιθήσονται μετὰ Ἀβραὰμ καὶ Ἰσαὰκ καὶ Ἰακὼβ ἐν τῇ βασιλείᾳ τῶν οὐρανῶν·
		% 12 οἱ δὲ υἱοὶ τῆς βασιλείας ἐκβληθήσονται εἰς τὸ σκότος τὸ ἐξώτερον· ἐκεῖ ἔσται ὁ κλαυθμὸς καὶ ὁ βρυγμὸς τῶν ὀδόντων.
		% 13 καὶ εἶπεν ὁ Ἰησοῦς τῷ ἑκατοντάρχῃ· Ὕπαγε, ὡς ἐπίστευσας γενηθήτω σοι· καὶ ἰάθη ὁ παῖς ἐν τῇ ὥρᾳ ἐκείνῃ.
		% }
		
	% % ----------------------
	% \subsection{Luke 7:1--10}
	% % ----------------------
		% \label{LCHRIST-066-L}
		
		% \gr{
		% 1 Ἐπειδὴ ἐπλήρωσεν πάντα τὰ ῥήματα αὐτοῦ εἰς τὰς ἀκοὰς τοῦ λαοῦ, εἰσῆλθεν εἰς Καφαρναούμ.
		% 2 Ἑκατοντάρχου δέ τινος δοῦλος κακῶς ἔχων ἤμελλεν τελευτᾶν, ὃς ἦν αὐτῷ ἔντιμος.
		% 3 ἀκούσας δὲ περὶ τοῦ Ἰησοῦ ἀπέστειλεν πρὸς αὐτὸν πρεσβυτέρους τῶν Ἰουδαίων, ἐρωτῶν αὐτὸν ὅπως ἐλθὼν διασώσῃ τὸν δοῦλον αὐτοῦ.
		% 4 οἱ δὲ παραγενόμενοι πρὸς τὸν Ἰησοῦν παρεκάλουν αὐτὸν σπουδαίως λέγοντες ὅτι Ἄξιός ἐστιν ᾧ παρέξῃ τοῦτο,
		% 5 ἀγαπᾷ γὰρ τὸ ἔθνος ἡμῶν καὶ τὴν συναγωγὴν αὐτὸς ᾠκοδόμησεν ἡμῖν.
		% 6 ὁ δὲ Ἰησοῦς ἐπορεύετο σὺν αὐτοῖς. ἤδη δὲ αὐτοῦ οὐ μακρὰν ἀπέχοντος ἀπὸ τῆς οἰκίας ἔπεμψεν φίλους ὁ ἑκατοντάρχης λέγων αὐτῷ· Κύριε, μὴ σκύλλου, οὐ γὰρ ἱκανός εἰμι ἵνα ὑπὸ τὴν στέγην μου εἰσέλθῃς·
		% 7 διὸ οὐδὲ ἐμαυτὸν ἠξίωσα πρὸς σὲ ἐλθεῖν· ἀλλὰ εἰπὲ λόγῳ, καὶ ἰαθήτω ὁ παῖς μου·
		% 8 καὶ γὰρ ἐγὼ ἄνθρωπός εἰμι ὑπὸ ἐξουσίαν τασσόμενος, ἔχων ὑπ᾽ ἐμαυτὸν στρατιώτας, καὶ λέγω τούτῳ· Πορεύθητι, καὶ πορεύεται, καὶ ἄλλῳ· Ἔρχου, καὶ ἔρχεται, καὶ τῷ δούλῳ μου· Ποίησον τοῦτο, καὶ ποιεῖ.
		% 9 ἀκούσας δὲ ταῦτα ὁ Ἰησοῦς ἐθαύμασεν αὐτόν, καὶ στραφεὶς τῷ ἀκολουθοῦντι αὐτῷ ὄχλῳ εἶπεν· Λέγω ὑμῖν, οὐδὲ ἐν τῷ Ἰσραὴλ τοσαύτην πίστιν εὗρον.
		% 10 καὶ ὑποστρέψαντες εἰς τὸν οἶκον οἱ πεμφθέντες εὗρον τὸν δοῦλον ὑγιαίνοντα.
		% }
		
	% % ----------------------
	% \subsection{John 4:43--54}
	% % ----------------------
		% \label{LCHRIST-066-J}
		
		% \gr{
		% 43 Μετὰ δὲ τὰς δύο ἡμέρας ἐξῆλθεν ἐκεῖθεν εἰς τὴν Γαλιλαίαν·
		% 44 αὐτὸς γὰρ Ἰησοῦς ἐμαρτύρησεν ὅτι προφήτης ἐν τῇ ἰδίᾳ πατρίδι τιμὴν οὐκ ἔχει.
		% 45 ὅτε οὖν ἦλθεν εἰς τὴν Γαλιλαίαν, ἐδέξαντο αὐτὸν οἱ Γαλιλαῖοι, πάντα ἑωρακότες ὅσα ἐποίησεν ἐν Ἱεροσολύμοις ἐν τῇ ἑορτῇ, καὶ αὐτοὶ γὰρ ἦλθον εἰς τὴν ἑορτήν.
		% }
		
		% \gr{
		% 46 Ἦλθεν οὖν πάλιν εἰς τὴν Κανὰ τῆς Γαλιλαίας, ὅπου ἐποίησεν τὸ ὕδωρ οἶνον. καὶ ἦν τις βασιλικὸς οὗ ὁ υἱὸς ἠσθένει ἐν Καφαρναούμ.
		% 47 οὗτος ἀκούσας ὅτι Ἰησοῦς ἥκει ἐκ τῆς Ἰουδαίας εἰς τὴν Γαλιλαίαν ἀπῆλθεν πρὸς αὐτὸν καὶ ἠρώτα ἵνα καταβῇ καὶ ἰάσηται αὐτοῦ τὸν υἱόν, ἤμελλεν γὰρ ἀποθνῄσκειν.
		% 48 εἶπεν οὖν ὁ Ἰησοῦς πρὸς αὐτόν· Ἐὰν μὴ σημεῖα καὶ τέρατα ἴδητε, οὐ μὴ πιστεύσητε.
		% 49 λέγει πρὸς αὐτὸν ὁ βασιλικός· Κύριε, κατάβηθι πρὶν ἀποθανεῖν τὸ παιδίον μου.
		% 50 λέγει αὐτῷ ὁ Ἰησοῦς· Πορεύου· ὁ υἱός σου ζῇ. ἐπίστευσεν ὁ ἄνθρωπος τῷ λόγῳ ὃν εἶπεν αὐτῷ ὁ Ἰησοῦς καὶ ἐπορεύετο.
		% 51 ἤδη δὲ αὐτοῦ καταβαίνοντος οἱ δοῦλοι αὐτοῦ ὑπήντησαν αὐτῷ λέγοντες ὅτι ὁ παῖς αὐτοῦ ζῇ.
		% 52 ἐπύθετο οὖν τὴν ὥραν παρ᾽ αὐτῶν ἐν ᾗ κομψότερον ἔσχεν· εἶπαν οὖν αὐτῷ ὅτι Ἐχθὲς ὥραν ἑβδόμην ἀφῆκεν αὐτὸν ὁ πυρετός.
		% 53 ἔγνω οὖν ὁ πατὴρ ὅτι ἐκείνῃ τῇ ὥρᾳ ἐν ᾗ εἶπεν αὐτῷ ὁ Ἰησοῦς· Ὁ υἱός σου ζῇ, καὶ ἐπίστευσεν αὐτὸς καὶ ἡ οἰκία αὐτοῦ ὅλη.
		% 54 τοῦτο δὲ πάλιν δεύτερον σημεῖον ἐποίησεν ὁ Ἰησοῦς ἐλθὼν ἐκ τῆς Ἰουδαίας εἰς τὴν Γαλιλαίαν.
		% }
		
% % -----------------------------------------------------------
% \section{67 | Jesus Raises Widow's Son at Nain}
% % -----------------------------------------------------------
	% \label{LCHRIST-067}
	
	% % ----------------------
	% \subsection{Luke 7:11--17}
	% % ----------------------
		% \label{LCHRIST-067-L}
		
		% \gr{
		% 11 Καὶ ἐγένετο ἐν τῷ ἑξῆς ἐπορεύθη εἰς πόλιν καλουμένην Ναΐν, καὶ συνεπορεύοντο αὐτῷ οἱ μαθηταὶ αὐτοῦ καὶ ὄχλος πολύς.
		% 12 ὡς δὲ ἤγγισεν τῇ πύλῃ τῆς πόλεως, καὶ ἰδοὺ ἐξεκομίζετο τεθνηκὼς μονογενὴς υἱὸς τῇ μητρὶ αὐτοῦ, καὶ αὐτὴ ἦν χήρα, καὶ ὄχλος τῆς πόλεως ἱκανὸς ἦν σὺν αὐτῇ.
		% 13 καὶ ἰδὼν αὐτὴν ὁ κύριος ἐσπλαγχνίσθη ἐπ᾽ αὐτῇ καὶ εἶπεν αὐτῇ· Μὴ κλαῖε.
		% 14 καὶ προσελθὼν ἥψατο τῆς σοροῦ, οἱ δὲ βαστάζοντες ἔστησαν, καὶ εἶπεν· Νεανίσκε, σοὶ λέγω, ἐγέρθητι.
		% 15 καὶ ἀνεκάθισεν ὁ νεκρὸς καὶ ἤρξατο λαλεῖν, καὶ ἔδωκεν αὐτὸν τῇ μητρὶ αὐτοῦ.
		% 16 ἔλαβεν δὲ φόβος πάντας, καὶ ἐδόξαζον τὸν θεὸν λέγοντες ὅτι Προφήτης μέγας ἠγέρθη ἐν ἡμῖν, καὶ ὅτι Ἐπεσκέψατο ὁ θεὸς τὸν λαὸν αὐτοῦ.
		% 17 καὶ ἐξῆλθεν ὁ λόγος οὗτος ἐν ὅλῃ τῇ Ἰουδαίᾳ περὶ αὐτοῦ καὶ πάσῃ τῇ περιχώρῳ.
		% }
		
% % -----------------------------------------------------------
% \section{68 | Jesus Heals the Sick, Casts Out Devils}
% % -----------------------------------------------------------
	% \label{LCHRIST-068}
	
	% % ----------------------
	% \subsection{Matthew 8:16--17}
	% % ----------------------
		% \label{LCHRIST-068-MT}
		
		% \gr{
		% 16 Ὀψίας δὲ γενομένης προσήνεγκαν αὐτῷ δαιμονιζομένους πολλούς· καὶ ἐξέβαλεν τὰ πνεύματα λόγῳ, καὶ πάντας τοὺς κακῶς ἔχοντας ἐθεράπευσεν·
		% 17 ὅπως πληρωθῇ τὸ ῥηθὲν διὰ Ἠσαΐου τοῦ προφήτου λέγοντος· Αὐτὸς τὰς ἀσθενείας ἡμῶν ἔλαβεν καὶ τὰς νόσους ἐβάστασεν.
		% }
		
	% % ----------------------
	% \subsection{Mark 1:32--34}
	% % ----------------------
		% \label{LCHRIST-068-MK}
		
		% \gr{
		% 32 Ὀψίας δὲ γενομένης, ὅτε ἔδυ ὁ ἥλιος, ἔφερον πρὸς αὐτὸν πάντας τοὺς κακῶς ἔχοντας καὶ τοὺς δαιμονιζομένους·
		% 33 καὶ ἦν ὅλη ἡ πόλις ἐπισυνηγμένη πρὸς τὴν θύραν.
		% 34 καὶ ἐθεράπευσεν πολλοὺς κακῶς ἔχοντας ποικίλαις νόσοις, καὶ δαιμόνια πολλὰ ἐξέβαλεν, καὶ οὐκ ἤφιεν λαλεῖν τὰ δαιμόνια, ὅτι ᾔδεισαν αὐτόν.
		% }
		
	% % ----------------------
	% \subsection{Luke 4:40--41}
	% % ----------------------
		% \label{LCHRIST-068-L}
		
		% \gr{
		% 40 Δύνοντος δὲ τοῦ ἡλίου ἅπαντες ὅσοι εἶχον ἀσθενοῦντας νόσοις ποικίλαις ἤγαγον αὐτοὺς πρὸς αὐτόν· ὁ δὲ ἑνὶ ἑκάστῳ αὐτῶν τὰς χεῖρας ἐπιτιθεὶς ἐθεράπευεν αὐτούς.
		% 41 ἐξήρχετο δὲ καὶ δαιμόνια ἀπὸ πολλῶν κραυγάζοντα καὶ λέγοντα ὅτι Σὺ εἶ ὁ υἱὸς τοῦ θεοῦ. καὶ ἐπιτιμῶν οὐκ εἴα αὐτὰ λαλεῖν, ὅτι ᾔδεισαν τὸν χριστὸν αὐτὸν εἶναι.
		% }
		
% % -----------------------------------------------------------
% \section{69 | Would-be Followers}
% % -----------------------------------------------------------
	% \label{LCHRIST-069}
	
	% % ----------------------
	% \subsection{Matthew 8:18--22}
	% % ----------------------
		% \label{LCHRIST-069-MT}
		
		% \gr{
		% 18 Ἰδὼν δὲ ὁ Ἰησοῦς πολλοὺς ὄχλους περὶ αὐτὸν ἐκέλευσεν ἀπελθεῖν εἰς τὸ πέραν.
		% 19 καὶ προσελθὼν εἷς γραμματεὺς εἶπεν αὐτῷ· Διδάσκαλε, ἀκολουθήσω σοι ὅπου ἐὰν ἀπέρχῃ.
		% 20 καὶ λέγει αὐτῷ ὁ Ἰησοῦς· Αἱ ἀλώπεκες φωλεοὺς ἔχουσιν καὶ τὰ πετεινὰ τοῦ οὐρανοῦ κατασκηνώσεις, ὁ δὲ υἱὸς τοῦ ἀνθρώπου οὐκ ἔχει ποῦ τὴν κεφαλὴν κλίνῃ.
		% 21 ἕτερος δὲ τῶν μαθητῶν εἶπεν αὐτῷ· Κύριε, ἐπίτρεψόν μοι πρῶτον ἀπελθεῖν καὶ θάψαι τὸν πατέρα μου.
		% 22 ὁ δὲ Ἰησοῦς λέγει αὐτῷ· Ἀκολούθει μοι, καὶ ἄφες τοὺς νεκροὺς θάψαι τοὺς ἑαυτῶν νεκρούς.
		% }
		
	% % ----------------------
	% \subsection{Luke 9:57--62}
	% % ----------------------
		% \label{LCHRIST-069-L}
		
		% \gr{
		% 57 Καὶ πορευομένων αὐτῶν ἐν τῇ ὁδῷ εἶπέν τις πρὸς αὐτόν· Ἀκολουθήσω σοι ὅπου ἐὰν ἀπέρχῃ.
		% 58 καὶ εἶπεν αὐτῷ ὁ Ἰησοῦς· Αἱ ἀλώπεκες φωλεοὺς ἔχουσιν καὶ τὰ πετεινὰ τοῦ οὐρανοῦ κατασκηνώσεις, ὁ δὲ υἱὸς τοῦ ἀνθρώπου οὐκ ἔχει ποῦ τὴν κεφαλὴν κλίνῃ.
		% 59 εἶπεν δὲ πρὸς ἕτερον· Ἀκολούθει μοι. ὁ δὲ εἶπεν· Κύριε, ἐπίτρεψόν μοι ἀπελθόντι πρῶτον θάψαι τὸν πατέρα μου.
		% 60 εἶπεν δὲ αὐτῷ· Ἄφες τοὺς νεκροὺς θάψαι τοὺς ἑαυτῶν νεκρούς, σὺ δὲ ἀπελθὼν διάγγελλε τὴν βασιλείαν τοῦ θεοῦ.
		% 61 εἶπεν δὲ καὶ ἕτερος· Ἀκολουθήσω σοι, κύριε· πρῶτον δὲ ἐπίτρεψόν μοι ἀποτάξασθαι τοῖς εἰς τὸν οἶκόν μου.
		% 62 εἶπεν δὲ ὁ Ἰησοῦς· Οὐδεὶς ἐπιβαλὼν τὴν χεῖρα ἐπ᾽ ἄροτρον καὶ βλέπων εἰς τὰ ὀπίσω εὔθετός ἐστιν τῇ βασιλείᾳ τοῦ θεοῦ.
		% }
		
% % -----------------------------------------------------------
% \section{70 | Stilling the Tempest}
% % -----------------------------------------------------------
	% \label{LCHRIST-070}
	
	% % ----------------------
	% \subsection{Matthew 8:23--27}
	% % ----------------------
		% \label{LCHRIST-070-MT}
		
		% \gr{
		% 23 Καὶ ἐμβάντι αὐτῷ εἰς πλοῖον ἠκολούθησαν αὐτῷ οἱ μαθηταὶ αὐτοῦ.
		% 24 καὶ ἰδοὺ σεισμὸς μέγας ἐγένετο ἐν τῇ θαλάσσῃ, ὥστε τὸ πλοῖον καλύπτεσθαι ὑπὸ τῶν κυμάτων, αὐτὸς δὲ ἐκάθευδεν.
		% 25 καὶ προσελθόντες ἤγειραν αὐτὸν λέγοντες· Κύριε, σῶσον, ἀπολλύμεθα.
		% 26 καὶ λέγει αὐτοῖς· Τί δειλοί ἐστε, ὀλιγόπιστοι; τότε ἐγερθεὶς ἐπετίμησεν τοῖς ἀνέμοις καὶ τῇ θαλάσσῃ, καὶ ἐγένετο γαλήνη μεγάλη.
		% 27 οἱ δὲ ἄνθρωποι ἐθαύμασαν λέγοντες· Ποταπός ἐστιν οὗτος ὅτι καὶ οἱ ἄνεμοι καὶ ἡ θάλασσα αὐτῷ ὑπακούουσιν;
		% }
		
	% % ----------------------
	% \subsection{Mark 4:36--41}
	% % ----------------------
		% \label{LCHRIST-070-MK}
		
		% \gr{
		% 36 καὶ ἀφέντες τὸν ὄχλον παραλαμβάνουσιν αὐτὸν ὡς ἦν ἐν τῷ πλοίῳ, καὶ ἄλλα πλοῖα ἦν μετ᾽ αὐτοῦ.
		% 37 καὶ γίνεται λαῖλαψ μεγάλη ἀνέμου, καὶ τὰ κύματα ἐπέβαλλεν εἰς τὸ πλοῖον, ὥστε ἤδη γεμίζεσθαι τὸ πλοῖον.
		% 38 καὶ αὐτὸς ἦν ἐν τῇ πρύμνῃ ἐπὶ τὸ προσκεφάλαιον καθεύδων· καὶ ἐγείρουσιν αὐτὸν καὶ λέγουσιν αὐτῷ· Διδάσκαλε, οὐ μέλει σοι ὅτι ἀπολλύμεθα;
		% 39 καὶ διεγερθεὶς ἐπετίμησεν τῷ ἀνέμῳ καὶ εἶπεν τῇ θαλάσσῃ· Σιώπα, πεφίμωσο. καὶ ἐκόπασεν ὁ ἄνεμος, καὶ ἐγένετο γαλήνη μεγάλη.
		% 40 καὶ εἶπεν αὐτοῖς· Τί δειλοί ἐστε; οὔπω ἔχετε πίστιν;
		% 41 καὶ ἐφοβήθησαν φόβον μέγαν, καὶ ἔλεγον πρὸς ἀλλήλους· Τίς ἄρα οὗτός ἐστιν ὅτι καὶ ὁ ἄνεμος καὶ ἡ θάλασσα ὑπακούει αὐτῷ;
		% }
		
	% % ----------------------
	% \subsection{Luke 8:22--25}
	% % ----------------------
		% \label{LCHRIST-070-L}
		
		% \gr{
		% 22 Ἐγένετο δὲ ἐν μιᾷ τῶν ἡμερῶν καὶ αὐτὸς ἐνέβη εἰς πλοῖον καὶ οἱ μαθηταὶ αὐτοῦ, καὶ εἶπεν πρὸς αὐτούς· Διέλθωμεν εἰς τὸ πέραν τῆς λίμνης, καὶ ἀνήχθησαν.
		% 23 πλεόντων δὲ αὐτῶν ἀφύπνωσεν. καὶ κατέβη λαῖλαψ ἀνέμου εἰς τὴν λίμνην, καὶ συνεπληροῦντο καὶ ἐκινδύνευον.
		% 24 προσελθόντες δὲ διήγειραν αὐτὸν λέγοντες· Ἐπιστάτα ἐπιστάτα, ἀπολλύμεθα· ὁ δὲ διεγερθεὶς ἐπετίμησεν τῷ ἀνέμῳ καὶ τῷ κλύδωνι τοῦ ὕδατος, καὶ ἐπαύσαντο, καὶ ἐγένετο γαλήνη.
		% 25 εἶπεν δὲ αὐτοῖς· Ποῦ ἡ πίστις ὑμῶν; φοβηθέντες δὲ ἐθαύμασαν, λέγοντες πρὸς ἀλλήλους· Τίς ἄρα οὗτός ἐστιν ὅτι καὶ τοῖς ἀνέμοις ἐπιτάσσει καὶ τῷ ὕδατι, καὶ ὑπακούουσιν αὐτῷ;
		% }
		
% % -----------------------------------------------------------
% \section{71 | Gadarene Demoniacs}
% % -----------------------------------------------------------
	% \label{LCHRIST-071}
	
	% % ----------------------
	% \subsection{Matthew 8:28--34}
	% % ----------------------
		% \label{LCHRIST-071-MT}
		
		% \gr{
		% 28 Καὶ ἐλθόντος αὐτοῦ εἰς τὸ πέραν εἰς τὴν χώραν τῶν Γαδαρηνῶν ὑπήντησαν αὐτῷ δύο δαιμονιζόμενοι ἐκ τῶν μνημείων ἐξερχόμενοι, χαλεποὶ λίαν ὥστε μὴ ἰσχύειν τινὰ παρελθεῖν διὰ τῆς ὁδοῦ ἐκείνης.
		% 29 καὶ ἰδοὺ ἔκραξαν λέγοντες· Τί ἡμῖν καὶ σοί, υἱὲ τοῦ θεοῦ; ἦλθες ὧδε πρὸ καιροῦ βασανίσαι ἡμᾶς;
		% 30 ἦν δὲ μακρὰν ἀπ᾽ αὐτῶν ἀγέλη χοίρων πολλῶν βοσκομένη.
		% 31 οἱ δὲ δαίμονες παρεκάλουν αὐτὸν λέγοντες· Εἰ ἐκβάλλεις ἡμᾶς, ἀπόστειλον ἡμᾶς εἰς τὴν ἀγέλην τῶν χοίρων.
		% 32 καὶ εἶπεν αὐτοῖς· Ὑπάγετε. οἱ δὲ ἐξελθόντες ἀπῆλθον εἰς τοὺς χοίρους· καὶ ἰδοὺ ὥρμησεν πᾶσα ἡ ἀγέλη κατὰ τοῦ κρημνοῦ εἰς τὴν θάλασσαν, καὶ ἀπέθανον ἐν τοῖς ὕδασιν.
		% 33 οἱ δὲ βόσκοντες ἔφυγον, καὶ ἀπελθόντες εἰς τὴν πόλιν ἀπήγγειλαν πάντα καὶ τὰ τῶν δαιμονιζομένων.
		% 34 καὶ ἰδοὺ πᾶσα ἡ πόλις ἐξῆλθεν εἰς ὑπάντησιν τῷ Ἰησοῦ, καὶ ἰδόντες αὐτὸν παρεκάλεσαν ὅπως μεταβῇ ἀπὸ τῶν ὁρίων αὐτῶν.
		% }
		
	% % ----------------------
	% \subsection{Mark 5:1--20}
	% % ----------------------
		% \label{LCHRIST-071-MK}
		
		% \gr{
		% 1 Καὶ ἦλθον εἰς τὸ πέραν τῆς θαλάσσης εἰς τὴν χώραν τῶν Γερασηνῶν.
		% 2 καὶ ἐξελθόντος αὐτοῦ ἐκ τοῦ πλοίου εὐθὺς ὑπήντησεν αὐτῷ ἐκ τῶν μνημείων ἄνθρωπος ἐν πνεύματι ἀκαθάρτῳ,
		% 3 ὃς τὴν κατοίκησιν εἶχεν ἐν τοῖς μνήμασιν, καὶ οὐδὲ ἁλύσει οὐκέτι οὐδεὶς ἐδύνατο αὐτὸν δῆσαι
		% 4 διὰ τὸ αὐτὸν πολλάκις πέδαις καὶ ἁλύσεσι δεδέσθαι καὶ διεσπάσθαι ὑπ᾽ αὐτοῦ τὰς ἁλύσεις καὶ τὰς πέδας συντετρῖφθαι, καὶ οὐδεὶς ἴσχυεν αὐτὸν δαμάσαι·
		% 5 καὶ διὰ παντὸς νυκτὸς καὶ ἡμέρας ἐν τοῖς μνήμασιν καὶ ἐν τοῖς ὄρεσιν ἦν κράζων καὶ κατακόπτων ἑαυτὸν λίθοις.
		% 6 καὶ ἰδὼν τὸν Ἰησοῦν ἀπὸ μακρόθεν ἔδραμεν καὶ προσεκύνησεν αὐτόν,
		% 7 καὶ κράξας φωνῇ μεγάλῃ λέγει· Τί ἐμοὶ καὶ σοί, Ἰησοῦ υἱὲ τοῦ θεοῦ τοῦ ὑψίστου; ὁρκίζω σε τὸν θεόν, μή με βασανίσῃς.
		% 8 ἔλεγεν γὰρ αὐτῷ· Ἔξελθε τὸ πνεῦμα τὸ ἀκάθαρτον ἐκ τοῦ ἀνθρώπου.
		% 9 καὶ ἐπηρώτα αὐτόν· Τί ὄνομά σοι; καὶ λέγει αὐτῷ· Λεγιὼν ὄνομά μοι, ὅτι πολλοί ἐσμεν·
		% 10 καὶ παρεκάλει αὐτὸν πολλὰ ἵνα μὴ αὐτὰ ἀποστείλῃ ἔξω τῆς χώρας.
		% 11 ἦν δὲ ἐκεῖ πρὸς τῷ ὄρει ἀγέλη χοίρων μεγάλη βοσκομένη·
		% 12 καὶ παρεκάλεσαν αὐτὸν λέγοντες· Πέμψον ἡμᾶς εἰς τοὺς χοίρους, ἵνα εἰς αὐτοὺς εἰσέλθωμεν.
		% 13 καὶ ἐπέτρεψεν αὐτοῖς. καὶ ἐξελθόντα τὰ πνεύματα τὰ ἀκάθαρτα εἰσῆλθον εἰς τοὺς χοίρους, καὶ ὥρμησεν ἡ ἀγέλη κατὰ τοῦ κρημνοῦ εἰς τὴν θάλασσαν, ὡς δισχίλιοι, καὶ ἐπνίγοντο ἐν τῇ θαλάσσῃ.
		% 14 Καὶ οἱ βόσκοντες αὐτοὺς ἔφυγον καὶ ἀπήγγειλαν εἰς τὴν πόλιν καὶ εἰς τοὺς ἀγρούς· καὶ ἦλθον ἰδεῖν τί ἐστιν τὸ γεγονός.
		% 15 καὶ ἔρχονται πρὸς τὸν Ἰησοῦν, καὶ θεωροῦσιν τὸν δαιμονιζόμενον καθήμενον ἱματισμένον καὶ σωφρονοῦντα, τὸν ἐσχηκότα τὸν λεγιῶνα, καὶ ἐφοβήθησαν.
		% 16 καὶ διηγήσαντο αὐτοῖς οἱ ἰδόντες πῶς ἐγένετο τῷ δαιμονιζομένῳ καὶ περὶ τῶν χοίρων.
		% 17 καὶ ἤρξαντο παρακαλεῖν αὐτὸν ἀπελθεῖν ἀπὸ τῶν ὁρίων αὐτῶν.
		% 18 καὶ ἐμβαίνοντος αὐτοῦ εἰς τὸ πλοῖον παρεκάλει αὐτὸν ὁ δαιμονισθεὶς ἵνα μετ᾽ αὐτοῦ ᾖ.
		% 19 καὶ οὐκ ἀφῆκεν αὐτόν, ἀλλὰ λέγει αὐτῷ· Ὕπαγε εἰς τὸν οἶκόν σου πρὸς τοὺς σούς, καὶ ἀπάγγειλον αὐτοῖς ὅσα ὁ κύριός σοι πεποίηκεν καὶ ἠλέησέν σε.
		% 20 καὶ ἀπῆλθεν καὶ ἤρξατο κηρύσσειν ἐν τῇ Δεκαπόλει ὅσα ἐποίησεν αὐτῷ ὁ Ἰησοῦς, καὶ πάντες ἐθαύμαζον.
		% }
		
	% % ----------------------
	% \subsection{Luke 8:26--39}
	% % ----------------------
		% \label{LCHRIST-071-L}
		
		% \gr{
		% 26 Καὶ κατέπλευσαν εἰς τὴν χώραν τῶν Γερασηνῶν, ἥτις ἐστὶν ἀντιπέρα τῆς Γαλιλαίας.
		% 27 ἐξελθόντι δὲ αὐτῷ ἐπὶ τὴν γῆν ὑπήντησεν ἀνήρ τις ἐκ τῆς πόλεως ἔχων δαιμόνια· καὶ χρόνῳ ἱκανῷ οὐκ ἐνεδύσατο ἱμάτιον, καὶ ἐν οἰκίᾳ οὐκ ἔμενεν ἀλλ᾽ ἐν τοῖς μνήμασιν.
		% 28 ἰδὼν δὲ τὸν Ἰησοῦν ἀνακράξας προσέπεσεν αὐτῷ καὶ φωνῇ μεγάλῃ εἶπεν· Τί ἐμοὶ καὶ σοί, Ἰησοῦ υἱὲ τοῦ θεοῦ τοῦ ὑψίστου; δέομαί σου, μή με βασανίσῃς·
		% 29  παρήγγελλεν γὰρ τῷ πνεύματι τῷ ἀκαθάρτῳ ἐξελθεῖν ἀπὸ τοῦ ἀνθρώπου. πολλοῖς γὰρ χρόνοις συνηρπάκει αὐτόν, καὶ ἐδεσμεύετο ἁλύσεσιν καὶ πέδαις φυλασσόμενος, καὶ διαρρήσσων τὰ δεσμὰ ἠλαύνετο ὑπὸ τοῦ δαιμονίου εἰς τὰς ἐρήμους.
		% 30 ἐπηρώτησεν δὲ αὐτὸν ὁ Ἰησοῦς· Τί σοι ὄνομά ἐστιν; ὁ δὲ εἶπεν· Λεγιών, ὅτι εἰσῆλθεν δαιμόνια πολλὰ εἰς αὐτόν.
		% 31 καὶ παρεκάλουν αὐτὸν ἵνα μὴ ἐπιτάξῃ αὐτοῖς εἰς τὴν ἄβυσσον ἀπελθεῖν.
		% }
		
		% \gr{
		% 32 Ἦν δὲ ἐκεῖ ἀγέλη χοίρων ἱκανῶν βοσκομένη ἐν τῷ ὄρει· καὶ παρεκάλεσαν αὐτὸν ἵνα ἐπιτρέψῃ αὐτοῖς εἰς ἐκείνους εἰσελθεῖν· καὶ ἐπέτρεψεν αὐτοῖς.
		% 33 ἐξελθόντα δὲ τὰ δαιμόνια ἀπὸ τοῦ ἀνθρώπου εἰσῆλθον εἰς τοὺς χοίρους, καὶ ὥρμησεν ἡ ἀγέλη κατὰ τοῦ κρημνοῦ εἰς τὴν λίμνην καὶ ἀπεπνίγη.
		% }
		
		% \gr{
		% 34 Ἰδόντες δὲ οἱ βόσκοντες τὸ γεγονὸς ἔφυγον καὶ ἀπήγγειλαν εἰς τὴν πόλιν καὶ εἰς τοὺς ἀγρούς.
		% 35 ἐξῆλθον δὲ ἰδεῖν τὸ γεγονὸς καὶ ἦλθον πρὸς τὸν Ἰησοῦν, καὶ εὗρον καθήμενον τὸν ἄνθρωπον ἀφ᾽ οὗ τὰ δαιμόνια ἐξῆλθεν ἱματισμένον καὶ σωφρονοῦντα παρὰ τοὺς πόδας τοῦ Ἰησοῦ, καὶ ἐφοβήθησαν.
		% 36 ἀπήγγειλαν δὲ αὐτοῖς οἱ ἰδόντες πῶς ἐσώθη ὁ δαιμονισθείς.
		% 37 καὶ ἠρώτησεν αὐτὸν ἅπαν τὸ πλῆθος τῆς περιχώρου τῶν Γερασηνῶν ἀπελθεῖν ἀπ᾽ αὐτῶν, ὅτι φόβῳ μεγάλῳ συνείχοντο· αὐτὸς δὲ ἐμβὰς εἰς πλοῖον ὑπέστρεψεν.
		% 38 ἐδεῖτο δὲ αὐτοῦ ὁ ἀνὴρ ἀφ᾽ οὗ ἐξεληλύθει τὰ δαιμόνια εἶναι σὺν αὐτῷ· ἀπέλυσεν δὲ αὐτὸν λέγων·
		% 39 Ὑπόστρεφε εἰς τὸν οἶκόν σου, καὶ διηγοῦ ὅσα σοι ἐποίησεν ὁ θεός. καὶ ἀπῆλθεν καθ᾽ ὅλην τὴν πόλιν κηρύσσων ὅσα ἐποίησεν αὐτῷ ὁ Ἰησοῦς.
		% }
		
% % -----------------------------------------------------------
% \section{72 | Easier to Heal Sick or Forgive Sins?}
% % -----------------------------------------------------------
	% \label{LCHRIST-072}
	
	% % ----------------------
	% \subsection{Matthew 9:1--8}
	% % ----------------------
		% \label{LCHRIST-072-MT}
		
		% \gr{
		% 1 Καὶ ἐμβὰς εἰς πλοῖον διεπέρασεν καὶ ἦλθεν εἰς τὴν ἰδίαν πόλιν.
		% }
		
		% \gr{
		% 2 Καὶ ἰδοὺ προσέφερον αὐτῷ παραλυτικὸν ἐπὶ κλίνης βεβλημένον. καὶ ἰδὼν ὁ Ἰησοῦς τὴν πίστιν αὐτῶν εἶπεν τῷ παραλυτικῷ· Θάρσει, τέκνον· ἀφίενταί σου αἱ ἁμαρτίαι.
		% 3 καὶ ἰδού τινες τῶν γραμματέων εἶπαν ἐν ἑαυτοῖς· Οὗτος βλασφημεῖ.
		% 4 καὶ εἰδὼς ὁ Ἰησοῦς τὰς ἐνθυμήσεις αὐτῶν εἶπεν· Ἱνατί ἐνθυμεῖσθε πονηρὰ ἐν ταῖς καρδίαις ὑμῶν;
		% 5 τί γάρ ἐστιν εὐκοπώτερον, εἰπεῖν· Ἀφίενταί σου αἱ ἁμαρτίαι, ἢ εἰπεῖν· Ἔγειρε καὶ περιπάτει;
		% 6 ἵνα δὲ εἰδῆτε ὅτι ἐξουσίαν ἔχει ὁ υἱὸς τοῦ ἀνθρώπου ἐπὶ τῆς γῆς ἀφιέναι ἁμαρτίας— τότε λέγει τῷ παραλυτικῷ· Ἐγερθεὶς ἆρόν σου τὴν κλίνην καὶ ὕπαγε εἰς τὸν οἶκόν σου.
		% 7 καὶ ἐγερθεὶς ἀπῆλθεν εἰς τὸν οἶκον αὐτοῦ.
		% 8 ἰδόντες δὲ οἱ ὄχλοι ἐφοβήθησαν καὶ ἐδόξασαν τὸν θεὸν τὸν δόντα ἐξουσίαν τοιαύτην τοῖς ἀνθρώποις.
		% }
		
	% % ----------------------
	% \subsection{Mark 2:1--12}
	% % ----------------------
		% \label{LCHRIST-072-MK}
		
		% \gr{
		% 1 Καὶ εἰσελθὼν πάλιν εἰς Καφαρναοὺμ δι᾽ ἡμερῶν ἠκούσθη ὅτι ἐν οἴκῳ ἐστίν·
		% 2 καὶ συνήχθησαν πολλοὶ ὥστε μηκέτι χωρεῖν μηδὲ τὰ πρὸς τὴν θύραν, καὶ ἐλάλει αὐτοῖς τὸν λόγον.
		% 3 καὶ ἔρχονται φέροντες πρὸς αὐτὸν παραλυτικὸν αἰρόμενον ὑπὸ τεσσάρων.
		% 4 καὶ μὴ δυνάμενοι προσενέγκαι αὐτῷ διὰ τὸν ὄχλον ἀπεστέγασαν τὴν στέγην ὅπου ἦν, καὶ ἐξορύξαντες χαλῶσι τὸν κράβαττον ὅπου ὁ παραλυτικὸς κατέκειτο.
		% 5 καὶ ἰδὼν ὁ Ἰησοῦς τὴν πίστιν αὐτῶν λέγει τῷ παραλυτικῷ· Τέκνον, ἀφίενταί σου αἱ ἁμαρτίαι.
		% 6 ἦσαν δέ τινες τῶν γραμματέων ἐκεῖ καθήμενοι καὶ διαλογιζόμενοι ἐν ταῖς καρδίαις αὐτῶν·
		% 7 Τί οὗτος οὕτως λαλεῖ; βλασφημεῖ· τίς δύναται ἀφιέναι ἁμαρτίας εἰ μὴ εἷς ὁ θεός;
		% 8 καὶ εὐθὺς ἐπιγνοὺς ὁ Ἰησοῦς τῷ πνεύματι αὐτοῦ ὅτι οὕτως διαλογίζονται ἐν ἑαυτοῖς λέγει αὐτοῖς· Τί ταῦτα διαλογίζεσθε ἐν ταῖς καρδίαις ὑμῶν;
		% 9 τί ἐστιν εὐκοπώτερον, εἰπεῖν τῷ παραλυτικῷ· Ἀφίενταί σου αἱ ἁμαρτίαι, ἢ εἰπεῖν· Ἔγειρε καὶ ἆρον τὸν κράβαττόν σου καὶ περιπάτει;
		% 10 ἵνα δὲ εἰδῆτε ὅτι ἐξουσίαν ἔχει ὁ υἱὸς τοῦ ἀνθρώπου ἐπὶ τῆς γῆς ἀφιέναι ἁμαρτίας— λέγει τῷ παραλυτικῷ·
		% 11 Σοὶ λέγω, ἔγειρε ἆρον τὸν κράβαττόν σου καὶ ὕπαγε εἰς τὸν οἶκόν σου.
		% 12 καὶ ἠγέρθη καὶ εὐθὺς ἄρας τὸν κράβαττον ἐξῆλθεν ἔμπροσθεν πάντων, ὥστε ἐξίστασθαι πάντας καὶ δοξάζειν τὸν θεὸν λέγοντας ὅτι Οὕτως οὐδέποτε εἴδομεν.
		% }
		
	% % ----------------------
	% \subsection{Luke 5:27--32}
	% % ----------------------
		% \label{LCHRIST-072-L}
		
		% \gr{
		% 27 Καὶ μετὰ ταῦτα ἐξῆλθεν καὶ ἐθεάσατο τελώνην ὀνόματι Λευὶν καθήμενον ἐπὶ τὸ τελώνιον, καὶ εἶπεν αὐτῷ· Ἀκολούθει μοι.
		% 28 καὶ καταλιπὼν πάντα ἀναστὰς ἠκολούθει αὐτῷ.
		% }
		
		% \gr{
		% 29 Καὶ ἐποίησεν δοχὴν μεγάλην Λευὶς αὐτῷ ἐν τῇ οἰκίᾳ αὐτοῦ· καὶ ἦν ὄχλος πολὺς τελωνῶν καὶ ἄλλων οἳ ἦσαν μετ᾽ αὐτῶν κατακείμενοι.
		% 30 καὶ ἐγόγγυζον οἱ Φαρισαῖοι καὶ οἱ γραμματεῖς αὐτῶν πρὸς τοὺς μαθητὰς αὐτοῦ λέγοντες· Διὰ τί μετὰ τῶν τελωνῶν καὶ ἁμαρτωλῶν ἐσθίετε καὶ πίνετε;
		% 31 καὶ ἀποκριθεὶς ὁ Ἰησοῦς εἶπεν πρὸς αὐτούς· Οὐ χρείαν ἔχουσιν οἱ ὑγιαίνοντες ἰατροῦ ἀλλὰ οἱ κακῶς ἔχοντες·
		% 32 οὐκ ἐλήλυθα καλέσαι δικαίους ἀλλὰ ἁμαρτωλοὺς εἰς μετάνοιαν.
		% }
		
% % -----------------------------------------------------------
% \section{73 | Matthew and Tax Collecting}
% % -----------------------------------------------------------
	% \label{LCHRIST-073}
	
	% % ----------------------
	% \subsection{Matthew 9:9--13}
	% % ----------------------
		% \label{LCHRIST-073-MT}
		
		% \gr{
		% 9 Καὶ παράγων ὁ Ἰησοῦς ἐκεῖθεν εἶδεν ἄνθρωπον καθήμενον ἐπὶ τὸ τελώνιον, Μαθθαῖον λεγόμενον, καὶ λέγει αὐτῷ· Ἀκολούθει μοι· καὶ ἀναστὰς ἠκολούθησεν αὐτῷ.
		% }
		
		% \gr{
		% 10 Καὶ ἐγένετο αὐτοῦ ἀνακειμένου ἐν τῇ οἰκίᾳ, καὶ ἰδοὺ πολλοὶ τελῶναι καὶ ἁμαρτωλοὶ}%
		% \footnote{%
			% \label{LCHRIST-073-hamartoloi}%
			% In these passages, BDAG suggests we take \gr{ἁμαρτωλός} as a word which focuses on the person's status as an outsider---``\textit{irreligious, unobservant people, outsiders} of those who did not observe the Law in detail and therefore were shunned by observers of traditional precepts'' (51b, PDF 66). Christ was with them because he came to call those sorts of people to repentance or to a change of heart and mind. It reminds me of something Brad was asking about yesterday (7/25/2021), in \hyperref[DC82-22]{DC 82:22} where Christ quotes \hyperref[LUKE-16-9]{Luke 16:9} in instructing his followers to ``make unto yourselves friends with the mammon of unrighteousness.'' There's work to be done among those people in helping them elevate their level of righteousness, and we have a lot to learn from them as well.
			% } 
		% \gr{ἐλθόντες συνανέκειντο τῷ Ἰησοῦ καὶ τοῖς μαθηταῖς αὐτοῦ.
		% 11 καὶ ἰδόντες οἱ Φαρισαῖοι ἔλεγον τοῖς μαθηταῖς αὐτοῦ· Διὰ τί μετὰ τῶν τελωνῶν καὶ ἁμαρτωλῶν ἐσθίει ὁ διδάσκαλος ὑμῶν;
		% 12 ὁ δὲ ἀκούσας εἶπεν· Οὐ χρείαν ἔχουσιν οἱ ἰσχύοντες ἰατροῦ ἀλλὰ οἱ κακῶς ἔχοντες.
		% 13 πορευθέντες δὲ μάθετε τί ἐστιν· Ἔλεος θέλω καὶ οὐ θυσίαν· οὐ γὰρ ἦλθον καλέσαι δικαίους ἀλλὰ ἁμαρτωλούς.
		% }
		
	% % ----------------------
	% \subsection{Mark 2:13--17}
	% % ---------------------- 
		% \label{LCHRIST-073-MK}
		
		% \gr{
		% 13 Καὶ ἐξῆλθεν πάλιν παρὰ τὴν θάλασσαν· καὶ πᾶς ὁ ὄχλος ἤρχετο πρὸς αὐτόν, καὶ ἐδίδασκεν αὐτούς.
		% 14 καὶ παράγων εἶδεν Λευὶν τὸν τοῦ Ἁλφαίου καθήμενον ἐπὶ τὸ τελώνιον, καὶ λέγει αὐτῷ· Ἀκολούθει μοι. καὶ ἀναστὰς ἠκολούθησεν αὐτῷ.
		% }
		
		% \gr{
		% 15 Καὶ γίνεται κατακεῖσθαι αὐτὸν ἐν τῇ οἰκίᾳ αὐτοῦ, καὶ πολλοὶ τελῶναι καὶ ἁμαρτωλοὶ συνανέκειντο τῷ Ἰησοῦ καὶ τοῖς μαθηταῖς αὐτοῦ, ἦσαν γὰρ πολλοὶ καὶ ἠκολούθουν αὐτῷ.
		% 16 καὶ οἱ γραμματεῖς τῶν Φαρισαίων ἰδόντες ὅτι ἐσθίει μετὰ τῶν ἁμαρτωλῶν καὶ τελωνῶν ἔλεγον τοῖς μαθηταῖς αὐτοῦ· Ὅτι μετὰ τῶν τελωνῶν καὶ ἁμαρτωλῶν ἐσθίει;
		% 17 καὶ ἀκούσας ὁ Ἰησοῦς λέγει αὐτοῖς ὅτι Οὐ χρείαν ἔχουσιν οἱ ἰσχύοντες ἰατροῦ ἀλλ᾽ οἱ κακῶς ἔχοντες· οὐκ ἦλθον καλέσαι δικαίους ἀλλὰ ἁμαρτωλούς.
		% }
	
	% % ----------------------
	% \subsection{Luke 5:27--32}
	% % ---------------------- 
		% \label{LCHRIST-073-L}
		
		% \gr{
		% 27 Καὶ μετὰ ταῦτα ἐξῆλθεν καὶ ἐθεάσατο τελώνην ὀνόματι Λευὶν καθήμενον ἐπὶ τὸ τελώνιον, καὶ εἶπεν αὐτῷ· Ἀκολούθει μοι.
		% 28 καὶ καταλιπὼν πάντα ἀναστὰς ἠκολούθει αὐτῷ.
		% }
		
		% \gr{
		% 29 Καὶ ἐποίησεν δοχὴν μεγάλην Λευὶς αὐτῷ ἐν τῇ οἰκίᾳ αὐτοῦ· καὶ ἦν ὄχλος πολὺς τελωνῶν καὶ ἄλλων οἳ ἦσαν μετ᾽ αὐτῶν κατακείμενοι.
		% 30 καὶ ἐγόγγυζον οἱ Φαρισαῖοι καὶ οἱ γραμματεῖς αὐτῶν πρὸς τοὺς μαθητὰς αὐτοῦ λέγοντες· Διὰ τί μετὰ τῶν τελωνῶν καὶ ἁμαρτωλῶν ἐσθίετε καὶ πίνετε;
		% 31 καὶ ἀποκριθεὶς ὁ Ἰησοῦς εἶπεν πρὸς αὐτούς· Οὐ χρείαν ἔχουσιν οἱ ὑγιαίνοντες ἰατροῦ ἀλλὰ οἱ κακῶς ἔχοντες·
		% 32 οὐκ ἐλήλυθα καλέσαι δικαίους ἀλλὰ ἁμαρτωλοὺς εἰς μετάνοιαν.
		% }
		
% % -----------------------------------------------------------
% \section{74 | Question about Fasting}
% % -----------------------------------------------------------
	% \label{LCHRIST-074}
	
	% % ----------------------
	% \subsection{Matthew 9:14--17}
	% % ----------------------
		% \label{LCHRIST-074-MT}
		
		% \gr{
		% 14 Τότε προσέρχονται αὐτῷ οἱ μαθηταὶ Ἰωάννου λέγοντες· Διὰ τί ἡμεῖς καὶ οἱ Φαρισαῖοι νηστεύομεν πολλά, οἱ δὲ μαθηταί σου οὐ νηστεύουσιν;
		% 15 καὶ εἶπεν αὐτοῖς ὁ Ἰησοῦς· Μὴ δύνανται οἱ υἱοὶ τοῦ νυμφῶνος πενθεῖν ἐφ᾽ ὅσον μετ᾽ αὐτῶν ἐστιν ὁ νυμφίος; ἐλεύσονται δὲ ἡμέραι ὅταν ἀπαρθῇ ἀπ᾽ αὐτῶν ὁ νυμφίος, καὶ τότε νηστεύσουσιν.
		% 16 οὐδεὶς δὲ ἐπιβάλλει ἐπίβλημα ῥάκους ἀγνάφου ἐπὶ ἱματίῳ παλαιῷ· αἴρει γὰρ τὸ πλήρωμα αὐτοῦ ἀπὸ τοῦ ἱματίου, καὶ χεῖρον σχίσμα γίνεται.
		% 17 οὐδὲ βάλλουσιν οἶνον νέον εἰς ἀσκοὺς παλαιούς· εἰ δὲ μή γε, ῥήγνυνται οἱ ἀσκοί, καὶ ὁ οἶνος ἐκχεῖται καὶ οἱ ἀσκοὶ ἀπόλλυνται· ἀλλὰ βάλλουσιν οἶνον νέον εἰς ἀσκοὺς καινούς, καὶ ἀμφότεροι συντηροῦνται.
		% }
		
	% % ----------------------
	% \subsection{Mark 2:18--22}
	% % ----------------------
		% \label{LCHRIST-074-MK}
		
		% \gr{
		% 18 Καὶ ἦσαν οἱ μαθηταὶ Ἰωάννου καὶ οἱ Φαρισαῖοι νηστεύοντες. καὶ ἔρχονται καὶ λέγουσιν αὐτῷ· Διὰ τί οἱ μαθηταὶ Ἰωάννου καὶ οἱ μαθηταὶ τῶν Φαρισαίων νηστεύουσιν, οἱ δὲ σοὶ μαθηταὶ οὐ νηστεύουσιν;
		% 19 καὶ εἶπεν αὐτοῖς ὁ Ἰησοῦς· Μὴ δύνανται οἱ υἱοὶ τοῦ νυμφῶνος ἐν ᾧ ὁ νυμφίος μετ᾽ αὐτῶν ἐστιν νηστεύειν; ὅσον χρόνον ἔχουσιν τὸν νυμφίον μετ᾽ αὐτῶν οὐ δύνανται νηστεύειν·
		% 20 ἐλεύσονται δὲ ἡμέραι ὅταν ἀπαρθῇ ἀπ᾽ αὐτῶν ὁ νυμφίος, καὶ τότε νηστεύσουσιν ἐν ἐκείνῃ τῇ ἡμέρᾳ.
		% }
		
		% \gr{
		% 21 Οὐδεὶς ἐπίβλημα ῥάκους ἀγνάφου ἐπιράπτει ἐπὶ ἱμάτιον παλαιόν· εἰ δὲ μή, αἴρει τὸ πλήρωμα ἀπ᾽ αὐτοῦ τὸ καινὸν τοῦ παλαιοῦ, καὶ χεῖρον σχίσμα γίνεται.
		% 22 καὶ οὐδεὶς βάλλει οἶνον νέον εἰς ἀσκοὺς παλαιούς· εἰ δὲ μή, ῥήξει ὁ οἶνος τοὺς ἀσκούς, καὶ ὁ οἶνος ἀπόλλυται καὶ οἱ ἀσκοί. ἀλλὰ οἶνον νέον εἰς ἀσκοὺς καινούς.
		% }
		
	% % ----------------------
	% \subsection{Luke 5:33--39}
	% % ---------------------- 
		% \label{LCHRIST-074-L}
		
		% \gr{
		% 33 Οἱ δὲ εἶπαν πρὸς αὐτόν· Οἱ μαθηταὶ Ἰωάννου νηστεύουσιν πυκνὰ καὶ δεήσεις ποιοῦνται, ὁμοίως καὶ οἱ τῶν Φαρισαίων, οἱ δὲ σοὶ ἐσθίουσιν καὶ πίνουσιν.
		% 34 ὁ δὲ εἶπεν πρὸς αὐτούς· Μὴ δύνασθε τοὺς υἱοὺς τοῦ νυμφῶνος ἐν ᾧ ὁ νυμφίος μετ᾽ αὐτῶν ἐστιν ποιῆσαι νηστεῦσαι;
		% 35 ἐλεύσονται δὲ ἡμέραι, καὶ ὅταν ἀπαρθῇ ἀπ᾽ αὐτῶν ὁ νυμφίος τότε νηστεύσουσιν ἐν ἐκείναις ταῖς ἡμέραις.
		% 36 ἔλεγεν δὲ καὶ παραβολὴν πρὸς αὐτοὺς ὅτι Οὐδεὶς ἐπίβλημα ἀπὸ ἱματίου καινοῦ σχίσας ἐπιβάλλει ἐπὶ ἱμάτιον παλαιόν· εἰ δὲ μήγε, καὶ τὸ καινὸν σχίσει καὶ τῷ παλαιῷ οὐ συμφωνήσει τὸ ἐπίβλημα τὸ ἀπὸ τοῦ καινοῦ.
		% 37 καὶ οὐδεὶς βάλλει οἶνον νέον εἰς ἀσκοὺς παλαιούς· εἰ δὲ μήγε, ῥήξει ὁ οἶνος ὁ νέος τοὺς ἀσκούς, καὶ αὐτὸς ἐκχυθήσεται καὶ οἱ ἀσκοὶ ἀπολοῦνται·
		% 38 ἀλλὰ οἶνον νέον εἰς ἀσκοὺς καινοὺς βλητέον.
		% 39 καὶ οὐδεὶς πιὼν παλαιὸν θέλει νέον· λέγει γάρ· Ὁ παλαιὸς χρηστός ἐστιν.
		% }
		
% % -----------------------------------------------------------
% \section{75 | Jairus' Daughter}
% % -----------------------------------------------------------
	% \label{LCHRIST-075}
	
	% % ----------------------
	% \subsection{Matthew 9:18--19, 23--26}
	% % ----------------------
		% \label{LCHRIST-075-MT}
		
		% \gr{
		% 18 Ταῦτα αὐτοῦ λαλοῦντος αὐτοῖς ἰδοὺ ἄρχων εἷς ἐλθὼν προσεκύνει αὐτῷ λέγων ὅτι Ἡ θυγάτηρ μου ἄρτι ἐτελεύτησεν· ἀλλὰ ἐλθὼν ἐπίθες τὴν χεῖρά σου ἐπ᾽ αὐτήν, καὶ ζήσεται.
		% 19 καὶ ἐγερθεὶς ὁ Ἰησοῦς ἠκολούθει αὐτῷ καὶ οἱ μαθηταὶ αὐτοῦ.
		% }
		
		% \gr{
		% 23 καὶ ἐλθὼν ὁ Ἰησοῦς εἰς τὴν οἰκίαν τοῦ ἄρχοντος καὶ ἰδὼν τοὺς αὐλητὰς καὶ τὸν ὄχλον θορυβούμενον
		% 24 ἔλεγεν· Ἀναχωρεῖτε, οὐ γὰρ ἀπέθανεν τὸ κοράσιον ἀλλὰ καθεύδει· καὶ κατεγέλων αὐτοῦ.
		% 25 ὅτε δὲ ἐξεβλήθη ὁ ὄχλος, εἰσελθὼν ἐκράτησεν τῆς χειρὸς αὐτῆς, καὶ ἠγέρθη τὸ κοράσιον.
		% 26 καὶ ἐξῆλθεν ἡ φήμη αὕτη εἰς ὅλην τὴν γῆν ἐκείνην.
		% }
		
	% % ----------------------
	% \subsection{Mark 5:21--24, 35--43}
	% % ----------------------
		% \label{LCHRIST-075-MK}
		
		% \gr{
		% 21 Καὶ διαπεράσαντος τοῦ Ἰησοῦ ἐν τῷ πλοίῳ πάλιν εἰς τὸ πέραν συνήχθη ὄχλος πολὺς ἐπ᾽ αὐτόν, καὶ ἦν παρὰ τὴν θάλασσαν.
		% 22 καὶ ἔρχεται εἷς τῶν ἀρχισυναγώγων, ὀνόματι Ἰάϊρος, καὶ ἰδὼν αὐτὸν πίπτει πρὸς τοὺς πόδας αὐτοῦ
		% 23 καὶ παρακαλεῖ αὐτὸν πολλὰ λέγων ὅτι Τὸ θυγάτριόν μου ἐσχάτως ἔχει, ἵνα ἐλθὼν ἐπιθῇς τὰς χεῖρας αὐτῇ ἵνα σωθῇ καὶ ζήσῃ.
		% 24 καὶ ἀπῆλθεν μετ᾽ αὐτοῦ.
		% }
		
		% \gr{
		% Καὶ ἠκολούθει αὐτῷ ὄχλος πολύς, καὶ συνέθλιβον αὐτόν.
		% }
		
		% \gr{
		% 35 Ἔτι αὐτοῦ λαλοῦντος ἔρχονται ἀπὸ τοῦ ἀρχισυναγώγου λέγοντες ὅτι Ἡ θυγάτηρ σου ἀπέθανεν· τί ἔτι σκύλλεις τὸν διδάσκαλον;
		% 36 ὁ δὲ Ἰησοῦς παρακούσας τὸν λόγον λαλούμενον λέγει τῷ ἀρχισυναγώγῳ· Μὴ φοβοῦ, μόνον πίστευε.
		% 37 καὶ οὐκ ἀφῆκεν οὐδένα μετ᾽ αὐτοῦ συνακολουθῆσαι εἰ μὴ τὸν Πέτρον καὶ Ἰάκωβον καὶ Ἰωάννην τὸν ἀδελφὸν Ἰακώβου.
		% 38 καὶ ἔρχονται εἰς τὸν οἶκον τοῦ ἀρχισυναγώγου, καὶ θεωρεῖ θόρυβον καὶ κλαίοντας καὶ ἀλαλάζοντας πολλά,
		% 39 καὶ εἰσελθὼν λέγει αὐτοῖς· Τί θορυβεῖσθε καὶ κλαίετε; τὸ παιδίον οὐκ ἀπέθανεν ἀλλὰ καθεύδει.
		% 40 καὶ κατεγέλων αὐτοῦ. αὐτὸς δὲ ἐκβαλὼν πάντας παραλαμβάνει τὸν πατέρα τοῦ παιδίου καὶ τὴν μητέρα καὶ τοὺς μετ᾽ αὐτοῦ, καὶ εἰσπορεύεται ὅπου ἦν τὸ παιδίον·
		% 41 καὶ κρατήσας τῆς χειρὸς τοῦ παιδίου λέγει αὐτῇ· Ταλιθα κουμ, ὅ ἐστιν μεθερμηνευόμενον· Τὸ κοράσιον, σοὶ λέγω, ἔγειρε.
		% 42 καὶ εὐθὺς ἀνέστη τὸ κοράσιον καὶ περιεπάτει, ἦν γὰρ ἐτῶν δώδεκα. καὶ ἐξέστησαν εὐθὺς ἐκστάσει μεγάλῃ.
		% 43 καὶ διεστείλατο αὐτοῖς πολλὰ ἵνα μηδεὶς γνοῖ τοῦτο, καὶ εἶπεν δοθῆναι αὐτῇ φαγεῖν.
		% }
		
	% % ----------------------
	% \subsection{Luke 8:41--42, 49--56}
	% % ----------------------
		% \label{LCHRIST-075-L}
		
		% \gr{
		% 41 καὶ ἰδοὺ ἦλθεν ἀνὴρ ᾧ ὄνομα Ἰάϊρος, καὶ οὗτος ἄρχων τῆς συναγωγῆς ὑπῆρχεν, καὶ πεσὼν παρὰ τοὺς πόδας τοῦ Ἰησοῦ παρεκάλει αὐτὸν εἰσελθεῖν εἰς τὸν οἶκον αὐτοῦ,
		% 42 ὅτι θυγάτηρ μονογενὴς ἦν αὐτῷ ὡς ἐτῶν δώδεκα καὶ αὐτὴ ἀπέθνῃσκεν.
		% }
		
		% \gr{
		% Ἐν δὲ τῷ ὑπάγειν αὐτὸν οἱ ὄχλοι συνέπνιγον αὐτόν.
		% }
		
		% \gr{
		% 49 Ἔτι αὐτοῦ λαλοῦντος ἔρχεταί τις παρὰ τοῦ ἀρχισυναγώγου λέγων ὅτι Τέθνηκεν ἡ θυγάτηρ σου, μηκέτι σκύλλε τὸν διδάσκαλον.
		% 50 ὁ δὲ Ἰησοῦς ἀκούσας ἀπεκρίθη αὐτῷ· Μὴ φοβοῦ, μόνον πίστευσον, καὶ σωθήσεται.
		% 51 ἐλθὼν δὲ εἰς τὴν οἰκίαν οὐκ ἀφῆκεν εἰσελθεῖν τινα σὺν αὐτῷ εἰ μὴ Πέτρον καὶ Ἰωάννην καὶ Ἰάκωβον καὶ τὸν πατέρα τῆς παιδὸς καὶ τὴν μητέρα.
		% 52 ἔκλαιον δὲ πάντες καὶ ἐκόπτοντο αὐτήν. ὁ δὲ εἶπεν· Μὴ κλαίετε, οὐ γὰρ ἀπέθανεν ἀλλὰ καθεύδει.
		% 53 καὶ κατεγέλων αὐτοῦ, εἰδότες ὅτι ἀπέθανεν.
		% 54 αὐτὸς δὲ κρατήσας τῆς χειρὸς αὐτῆς ἐφώνησεν λέγων· Ἡ παῖς, ἔγειρε.
		% 55 καὶ ἐπέστρεψεν τὸ πνεῦμα αὐτῆς, καὶ ἀνέστη παραχρῆμα, καὶ διέταξεν αὐτῇ δοθῆναι φαγεῖν.
		% 56 καὶ ἐξέστησαν οἱ γονεῖς αὐτῆς· ὁ δὲ παρήγγειλεν αὐτοῖς μηδενὶ εἰπεῖν τὸ γεγονός.
		% }
		
% % -----------------------------------------------------------
% \section{76 | Woman Who Touched Jesus}
% % -----------------------------------------------------------
	% \label{LCHRIST-076}
	
	% % ----------------------
	% \subsection{Matthew 9:20--22}
	% % ----------------------
		% \label{LCHRIST-076-MT}
		
		% \gr{
		% 20 Καὶ ἰδοὺ γυνὴ αἱμορροοῦσα δώδεκα ἔτη προσελθοῦσα ὄπισθεν ἥψατο τοῦ κρασπέδου τοῦ ἱματίου αὐτοῦ·
		% 21 ἔλεγεν γὰρ ἐν ἑαυτῇ· Ἐὰν μόνον ἅψωμαι τοῦ ἱματίου αὐτοῦ σωθήσομαι.
		% 22 ὁ δὲ Ἰησοῦς στραφεὶς καὶ ἰδὼν αὐτὴν εἶπεν· Θάρσει, θύγατερ· ἡ πίστις σου σέσωκέν σε. καὶ ἐσώθη ἡ γυνὴ ἀπὸ τῆς ὥρας ἐκείνης.
		% }
		
	% % ----------------------
	% \subsection{Mark 5:25--34}
	% % ---------------------- 
		% \label{LCHRIST-076-MK}
		
		% \gr{
		% 25 καὶ γυνὴ οὖσα ἐν ῥύσει αἵματος δώδεκα ἔτη
		% 26 καὶ πολλὰ παθοῦσα ὑπὸ πολλῶν ἰατρῶν καὶ δαπανήσασα τὰ παρ᾽ αὐτῆς πάντα καὶ μηδὲν ὠφεληθεῖσα ἀλλὰ μᾶλλον εἰς τὸ χεῖρον ἐλθοῦσα,
		% 27 ἀκούσασα περὶ τοῦ Ἰησοῦ, ἐλθοῦσα ἐν τῷ ὄχλῳ ὄπισθεν ἥψατο τοῦ ἱματίου αὐτοῦ·
		% 28 ἔλεγεν γὰρ ὅτι Ἐὰν ἅψωμαι κἂν τῶν ἱματίων αὐτοῦ σωθήσομαι.
		% 29 καὶ εὐθὺς ἐξηράνθη ἡ πηγὴ τοῦ αἵματος αὐτῆς, καὶ ἔγνω τῷ σώματι ὅτι ἴαται ἀπὸ τῆς μάστιγος.
		% 30 καὶ εὐθὺς ὁ Ἰησοῦς ἐπιγνοὺς ἐν ἑαυτῷ τὴν ἐξ αὐτοῦ δύναμιν ἐξελθοῦσαν ἐπιστραφεὶς ἐν τῷ ὄχλῳ ἔλεγεν· Τίς μου ἥψατο τῶν ἱματίων;
		% 31 καὶ ἔλεγον αὐτῷ οἱ μαθηταὶ αὐτοῦ· Βλέπεις τὸν ὄχλον συνθλίβοντά σε, καὶ λέγεις· Τίς μου ἥψατο;
		% 32 καὶ περιεβλέπετο ἰδεῖν τὴν τοῦτο ποιήσασαν.
		% 33 ἡ δὲ γυνὴ φοβηθεῖσα καὶ τρέμουσα, εἰδυῖα ὃ γέγονεν αὐτῇ, ἦλθεν καὶ προσέπεσεν αὐτῷ καὶ εἶπεν αὐτῷ πᾶσαν τὴν ἀλήθειαν.
		% 34 ὁ δὲ εἶπεν αὐτῇ· Θυγάτηρ, ἡ πίστις σου σέσωκέν σε· ὕπαγε εἰς εἰρήνην, καὶ ἴσθι ὑγιὴς ἀπὸ τῆς μάστιγός σου.
		% }
		
	% % ----------------------
	% \subsection{Luke 8:43--48}
	% % ----------------------
		% \label{LCHRIST-076-L}
		
		% \gr{
		% 43 καὶ γυνὴ οὖσα ἐν ῥύσει αἵματος ἀπὸ ἐτῶν δώδεκα, ἥτις ἰατροῖς προσαναλώσασα ὅλον τὸν βίον οὐκ ἴσχυσεν ἀπ᾽ οὐδενὸς θεραπευθῆναι,
		% 44 προσελθοῦσα ὄπισθεν ἥψατο τοῦ κρασπέδου τοῦ ἱματίου αὐτοῦ, καὶ παραχρῆμα ἔστη ἡ ῥύσις τοῦ αἵματος αὐτῆς.
		% 45 καὶ εἶπεν ὁ Ἰησοῦς· Τίς ὁ ἁψάμενός μου; ἀρνουμένων δὲ πάντων εἶπεν ὁ Πέτρος· Ἐπιστάτα, οἱ ὄχλοι συνέχουσίν σε καὶ ἀποθλίβουσιν.
		% 46 ὁ δὲ Ἰησοῦς εἶπεν· Ἥψατό μού τις, ἐγὼ γὰρ ἔγνων δύναμιν ἐξεληλυθυῖαν ἀπ᾽ ἐμοῦ.
		% 47 ἰδοῦσα δὲ ἡ γυνὴ ὅτι οὐκ ἔλαθεν τρέμουσα ἦλθεν καὶ προσπεσοῦσα αὐτῷ δι᾽ ἣν αἰτίαν ἥψατο αὐτοῦ ἀπήγγειλεν ἐνώπιον παντὸς τοῦ λαοῦ καὶ ὡς ἰάθη παραχρῆμα.
		% 48 ὁ δὲ εἶπεν αὐτῇ· Θυγάτηρ, ἡ πίστις σου σέσωκέν σε· πορεύου εἰς εἰρήνην.
		% }
		
% % -----------------------------------------------------------
% \section{77 | Two Blind Men Healed}
% % -----------------------------------------------------------
	% \label{LCHRIST-077}
	
	% % ----------------------
	% \subsection{Matthew 9:27--31}
	% % ----------------------
		% \label{LCHRIST-077-MT}
		
		% \gr{
		% 27 Καὶ παράγοντι ἐκεῖθεν τῷ Ἰησοῦ ἠκολούθησαν αὐτῷ δύο τυφλοὶ κράζοντες καὶ λέγοντες· Ἐλέησον ἡμᾶς, υἱὲ Δαυίδ.
		% 28 ἐλθόντι δὲ εἰς τὴν οἰκίαν προσῆλθον αὐτῷ οἱ τυφλοί, καὶ λέγει αὐτοῖς ὁ Ἰησοῦς· Πιστεύετε ὅτι δύναμαι τοῦτο ποιῆσαι; λέγουσιν αὐτῷ· Ναί, κύριε.
		% 29 τότε ἥψατο τῶν ὀφθαλμῶν αὐτῶν λέγων· Κατὰ τὴν πίστιν ὑμῶν γενηθήτω ὑμῖν.
		% 30 καὶ ἠνεῴχθησαν αὐτῶν οἱ ὀφθαλμοί. καὶ ἐνεβριμήθη αὐτοῖς ὁ Ἰησοῦς λέγων· Ὁρᾶτε μηδεὶς γινωσκέτω·
		% 31 οἱ δὲ ἐξελθόντες διεφήμισαν αὐτὸν ἐν ὅλῃ τῇ γῇ ἐκείνῃ.
		% }
		
% % -----------------------------------------------------------
% \section{78 | Dumb Man Healed}
% % -----------------------------------------------------------
	% \label{LCHRIST-078}
	
	% % ----------------------
	% \subsection{Matthew 9:32--34}
	% % ----------------------
		% \label{LCHRIST-078-MT}
		
		% \gr{
		% 32 Αὐτῶν δὲ ἐξερχομένων ἰδοὺ προσήνεγκαν αὐτῷ ἄνθρωπον κωφὸν δαιμονιζόμενον·
		% 33 καὶ ἐκβληθέντος τοῦ δαιμονίου ἐλάλησεν ὁ κωφός. καὶ ἐθαύμασαν οἱ ὄχλοι λέγοντες· Οὐδέποτε ἐφάνη οὕτως ἐν τῷ Ἰσραήλ.
		% 34 οἱ δὲ Φαρισαῖοι ἔλεγον· Ἐν τῷ ἄρχοντι τῶν δαιμονίων ἐκβάλλει τὰ δαιμόνια.
		% }
		
% % -----------------------------------------------------------
% \section{79 | Compassion of Jesus}
% % -----------------------------------------------------------
	% \label{LCHRIST-079}
	
	% % ----------------------
	% \subsection{Matthew 9:35--38}
	% % ----------------------
	
		% \gr{
		% 35 Καὶ περιῆγεν ὁ Ἰησοῦς τὰς πόλεις πάσας καὶ τὰς κώμας, διδάσκων ἐν ταῖς συναγωγαῖς αὐτῶν καὶ κηρύσσων τὸ εὐαγγέλιον τῆς βασιλείας καὶ θεραπεύων πᾶσαν νόσον καὶ πᾶσαν μαλακίαν.
		% 36 Ἰδὼν δὲ τοὺς ὄχλους ἐσπλαγχνίσθη περὶ αὐτῶν ὅτι ἦσαν ἐσκυλμένοι καὶ ἐρριμμένοι ὡσεὶ πρόβατα μὴ ἔχοντα ποιμένα.
		% 37 τότε λέγει τοῖς μαθηταῖς αὐτοῦ· Ὁ μὲν θερισμὸς πολύς, οἱ δὲ ἐργάται ὀλίγοι·
		% 38 δεήθητε οὖν τοῦ κυρίου τοῦ θερισμοῦ ὅπως ἐκβάλῃ ἐργάτας εἰς τὸν θερισμὸν αὐτοῦ.
		% }
	
% % -----------------------------------------------------------
% \section{80 | Calling the Twelve}
% % -----------------------------------------------------------
	% \label{LCHRIST-080}
	
	% % ----------------------
	% \subsection{Matthew 10:1--4}
	% % ----------------------
		% \label{LCHRIST-080-MT}
		
		% \gr{
		% 1 Καὶ προσκαλεσάμενος τοὺς δώδεκα μαθητὰς αὐτοῦ ἔδωκεν αὐτοῖς ἐξουσίαν}
			% \footnote{%
				% \label{LCHRIST-080-exousia}%
				% Note that both the Matthew and Mark accounts are pretty explicit about Christ giving the apostles \gr{ἐξουσία} to perform their work; see \hyperref[EXOUSIA]{\grl{ἐξουσία}}. Luke mentions this later in \hyperref[LCHRIST-081-L]{Luke 9:1}.
				% }
		% \gr{πνευμάτων ἀκαθάρτων ὥστε ἐκβάλλειν αὐτὰ καὶ θεραπεύειν πᾶσαν νόσον καὶ πᾶσαν μαλακίαν.
		% 2 τῶν δὲ δώδεκα ἀποστόλων τὰ ὀνόματά ἐστιν ταῦτα· πρῶτος Σίμων ὁ λεγόμενος Πέτρος καὶ Ἀνδρέας ὁ ἀδελφὸς αὐτοῦ, Ἰάκωβος ὁ τοῦ Ζεβεδαίου καὶ Ἰωάννης ὁ ἀδελφὸς αὐτοῦ,
		% 3 Φίλιππος καὶ Βαρθολομαῖος, Θωμᾶς καὶ Μαθθαῖος ὁ τελώνης, Ἰάκωβος ὁ τοῦ Ἁλφαίου καὶ Θαδδαῖος,
		% 4 Σίμων ὁ Καναναῖος καὶ Ἰούδας ὁ Ἰσκαριώτης ὁ καὶ παραδοὺς αὐτόν.
		% }
		
	% % ----------------------
	% \subsection{Mark 3:13--19}
	% % ----------------------
		% \label{LCHRIST-080-MK}
		
		% \gr{
		% 13 Καὶ ἀναβαίνει εἰς τὸ ὄρος καὶ προσκαλεῖται οὓς ἤθελεν αὐτός, καὶ ἀπῆλθον πρὸς αὐτόν.
		% 14 καὶ ἐποίησεν δώδεκα, ἵνα ὦσιν μετ᾽ αὐτοῦ καὶ ἵνα ἀποστέλλῃ αὐτοὺς κηρύσσειν
		% 15 καὶ ἔχειν ἐξουσίαν ἐκβάλλειν τὰ δαιμόνια·
		% 16 καὶ ἐποίησεν τοὺς δώδεκα, καὶ ἐπέθηκεν ὄνομα τῷ Σίμωνι Πέτρον,
		% 17 καὶ Ἰάκωβον τὸν τοῦ Ζεβεδαίου καὶ Ἰωάννην τὸν ἀδελφὸν τοῦ Ἰακώβου (καὶ ἐπέθηκεν αὐτοῖς ὀνόματα Βοανηργές, ὅ ἐστιν Υἱοὶ Βροντῆς),
		% 18 καὶ Ἀνδρέαν καὶ Φίλιππον καὶ Βαρθολομαῖον καὶ Μαθθαῖον καὶ Θωμᾶν καὶ Ἰάκωβον τὸν τοῦ Ἁλφαίου καὶ Θαδδαῖον καὶ Σίμωνα τὸν Καναναῖον
		% 19 καὶ Ἰούδαν Ἰσκαριώθ, ὃς καὶ παρέδωκεν αὐτόν.
		% }
		
	% % ----------------------
	% \subsection{Luke 9:12--16}
	% % ----------------------
		% \label{LCHRIST-080-L}
		
		% \gr{
		% 12 Ἐγένετο δὲ ἐν ταῖς ἡμέραις ταύταις ἐξελθεῖν αὐτὸν εἰς τὸ ὄρος προσεύξασθαι, καὶ ἦν διανυκτερεύων ἐν τῇ προσευχῇ τοῦ θεοῦ.
		% 13 καὶ ὅτε ἐγένετο ἡμέρα, προσεφώνησεν τοὺς μαθητὰς αὐτοῦ, καὶ ἐκλεξάμενος ἀπ᾽ αὐτῶν δώδεκα, οὓς καὶ ἀποστόλους ὠνόμασεν,
		% 14 Σίμωνα ὃν καὶ ὠνόμασεν Πέτρον καὶ Ἀνδρέαν τὸν ἀδελφὸν αὐτοῦ καὶ Ἰάκωβον καὶ Ἰωάννην καὶ Φίλιππον καὶ Βαρθολομαῖον
		% 15 καὶ Μαθθαῖον καὶ Θωμᾶν καὶ Ἰάκωβον Ἁλφαίου καὶ Σίμωνα τὸν καλούμενον Ζηλωτὴν
		% 16 καὶ Ἰούδαν Ἰακώβου καὶ Ἰούδαν Ἰσκαριὼθ ὃς ἐγένετο προδότης.
		% }
		
% % -----------------------------------------------------------
% \section{81 | Mission of the Twelve}
% % -----------------------------------------------------------
	% \label{LCHRIST-081}
	
	% % ----------------------
	% \subsection{Matthew 10:1--15}
	% % ----------------------
		% \label{LCHRIST-081-MT}
		
		% \gr{
		% 1 Καὶ προσκαλεσάμενος τοὺς δώδεκα μαθητὰς αὐτοῦ ἔδωκεν αὐτοῖς ἐξουσίαν πνευμάτων ἀκαθάρτων ὥστε ἐκβάλλειν αὐτὰ καὶ θεραπεύειν πᾶσαν νόσον καὶ πᾶσαν μαλακίαν.
		% 2 τῶν δὲ δώδεκα ἀποστόλων τὰ ὀνόματά ἐστιν ταῦτα· πρῶτος Σίμων ὁ λεγόμενος Πέτρος καὶ Ἀνδρέας ὁ ἀδελφὸς αὐτοῦ, Ἰάκωβος ὁ τοῦ Ζεβεδαίου καὶ Ἰωάννης ὁ ἀδελφὸς αὐτοῦ,
		% 3 Φίλιππος καὶ Βαρθολομαῖος, Θωμᾶς καὶ Μαθθαῖος ὁ τελώνης, Ἰάκωβος ὁ τοῦ Ἁλφαίου καὶ Θαδδαῖος,
		% 4 Σίμων ὁ Καναναῖος καὶ Ἰούδας ὁ Ἰσκαριώτης ὁ καὶ παραδοὺς αὐτόν.
		% }
		
		% \gr{
		% 5 Τούτους τοὺς δώδεκα ἀπέστειλεν ὁ Ἰησοῦς παραγγείλας αὐτοῖς λέγων· Εἰς ὁδὸν ἐθνῶν μὴ ἀπέλθητε καὶ εἰς πόλιν Σαμαριτῶν μὴ εἰσέλθητε·
		% 6 πορεύεσθε δὲ μᾶλλον πρὸς τὰ πρόβατα τὰ ἀπολωλότα οἴκου Ἰσραήλ.
		% 7 πορευόμενοι δὲ κηρύσσετε λέγοντες ὅτι Ἤγγικεν ἡ βασιλεία τῶν οὐρανῶν.
		% 8 ἀσθενοῦντας θεραπεύετε, νεκροὺς ἐγείρετε, λεπροὺς καθαρίζετε, δαιμόνια ἐκβάλλετε· δωρεὰν ἐλάβετε, δωρεὰν δότε.}
			% \footnote{%
				% \label{LCHRIST-081-dorean}%
				% This is a very powerful idea that Christ says to his apostles: \gr{δωρεὰν ἐλάβετε, δωρεὰν δότε}. He's saying, you received good things as gifts, so you should give your good things to the world as gifts. It makes me think of Deida and how a man has a gift to give the world, and the more of his gift he can shape his life to give, the more people benefit from it, and the richer that man's life actually is.
				% }
		% \gr{9 μὴ κτήσησθε χρυσὸν μηδὲ ἄργυρον μηδὲ χαλκὸν εἰς τὰς ζώνας ὑμῶν,
		% 10 μὴ πήραν εἰς ὁδὸν μηδὲ δύο χιτῶνας μηδὲ ὑποδήματα μηδὲ ῥάβδον· ἄξιος γὰρ ὁ ἐργάτης τῆς τροφῆς αὐτοῦ.
		% 11 εἰς ἣν δ᾽ ἂν πόλιν ἢ κώμην εἰσέλθητε, ἐξετάσατε τίς ἐν αὐτῇ ἄξιός ἐστιν· κἀκεῖ μείνατε ἕως ἂν ἐξέλθητε.
		% 12 εἰσερχόμενοι δὲ εἰς τὴν οἰκίαν ἀσπάσασθε αὐτήν·
		% 13 καὶ ἐὰν μὲν ᾖ ἡ οἰκία ἀξία, ἐλθάτω ἡ εἰρήνη ὑμῶν ἐπ᾽ αὐτήν· ἐὰν δὲ μὴ ᾖ ἀξία, ἡ εἰρήνη ὑμῶν πρὸς ὑμᾶς ἐπιστραφήτω.
		% 14 καὶ ὃς ἂν μὴ δέξηται ὑμᾶς μηδὲ ἀκούσῃ τοὺς λόγους ὑμῶν, ἐξερχόμενοι ἔξω τῆς οἰκίας ἢ τῆς πόλεως ἐκείνης ἐκτινάξατε τὸν κονιορτὸν τῶν ποδῶν ὑμῶν.
		% 15 ἀμὴν λέγω ὑμῖν, ἀνεκτότερον ἔσται γῇ Σοδόμων καὶ Γομόρρων ἐν ἡμέρᾳ κρίσεως ἢ τῇ πόλει ἐκείνῃ.
		% }
		
	% % ----------------------
	% \subsection{Mark 6:7--13}
	% % ----------------------
		% \label{LCHRIST-081-MK}
		
		% \gr{
		% 7 καὶ προσκαλεῖται τοὺς δώδεκα, καὶ ἤρξατο αὐτοὺς ἀποστέλλειν δύο δύο, καὶ ἐδίδου αὐτοῖς ἐξουσίαν τῶν πνευμάτων τῶν ἀκαθάρτων,
		% 8 καὶ παρήγγειλεν αὐτοῖς ἵνα μηδὲν αἴρωσιν εἰς ὁδὸν εἰ μὴ ῥάβδον μόνον, μὴ ἄρτον, μὴ πήραν, μὴ εἰς τὴν ζώνην χαλκόν,
		% 9 ἀλλὰ ὑποδεδεμένους σανδάλια, καὶ μὴ ἐνδύσησθε δύο χιτῶνας.
		% 10 καὶ ἔλεγεν αὐτοῖς· Ὅπου ἐὰν εἰσέλθητε εἰς οἰκίαν, ἐκεῖ μένετε ἕως ἂν ἐξέλθητε ἐκεῖθεν.
		% 11 καὶ ὃς ἂν τόπος μὴ δέξηται ὑμᾶς μηδὲ ἀκούσωσιν ὑμῶν, ἐκπορευόμενοι ἐκεῖθεν ἐκτινάξατε τὸν χοῦν τὸν ὑποκάτω τῶν ποδῶν ὑμῶν εἰς μαρτύριον αὐτοῖς.
		% 12 Καὶ ἐξελθόντες ἐκήρυξαν ἵνα μετανοῶσιν,
		% 13 καὶ δαιμόνια πολλὰ ἐξέβαλλον, καὶ ἤλειφον ἐλαίῳ πολλοὺς ἀρρώστους καὶ ἐθεράπευον.
		% }
		
	% % ----------------------
	% \subsection{Luke 9:1--6}
	% % ----------------------
		% \label{LCHRIST-081-L}
	
		% \gr{
		% 1 Συγκαλεσάμενος δὲ τοὺς δώδεκα ἔδωκεν αὐτοῖς δύναμιν καὶ ἐξουσίαν ἐπὶ πάντα τὰ δαιμόνια καὶ νόσους θεραπεύειν,
		% 2 καὶ ἀπέστειλεν αὐτοὺς κηρύσσειν τὴν βασιλείαν τοῦ θεοῦ καὶ ἰᾶσθαι τοὺς ἀσθενεῖς,
		% 3 καὶ εἶπεν πρὸς αὐτούς· Μηδὲν αἴρετε εἰς τὴν ὁδόν, μήτε ῥάβδον μήτε πήραν μήτε ἄρτον μήτε ἀργύριον, μήτε ἀνὰ δύο χιτῶνας ἔχειν.
		% 4 καὶ εἰς ἣν ἂν οἰκίαν εἰσέλθητε, ἐκεῖ μένετε καὶ ἐκεῖθεν ἐξέρχεσθε.
		% 5 καὶ ὅσοι ἂν μὴ δέχωνται ὑμᾶς, ἐξερχόμενοι ἀπὸ τῆς πόλεως ἐκείνης τὸν κονιορτὸν ἀπὸ τῶν ποδῶν ὑμῶν ἀποτινάσσετε εἰς μαρτύριον ἐπ᾽ αὐτούς.
		% 6 ἐξερχόμενοι δὲ διήρχοντο κατὰ τὰς κώμας εὐαγγελιζόμενοι καὶ θεραπεύοντες πανταχοῦ.
		% }
		
% % -----------------------------------------------------------
% \section{82 | Coming Persecutions}
% % -----------------------------------------------------------
	% \label{LCHRIST-082}
	
	% % ----------------------
	% \subsection{Matthew 10:16--25}
	% % ----------------------
		% \label{LCHRIST-082-MT}
		
		% \gr{
		% 16 Ἰδοὺ ἐγὼ ἀποστέλλω ὑμᾶς ὡς πρόβατα ἐν μέσῳ λύκων· γίνεσθε οὖν φρόνιμοι ὡς οἱ ὄφεις καὶ ἀκέραιοι ὡς αἱ περιστεραί.}%
			% \footnote{%
				% \label{LCHRIST-082-akeraioi}%
				% See \hyperref[ALMA-18-22-wise]{this note} at Alma 18:22 on the end of this verse.
				% }
		% \gr{17 προσέχετε δὲ ἀπὸ τῶν ἀνθρώπων· παραδώσουσιν γὰρ ὑμᾶς εἰς συνέδρια, καὶ ἐν ταῖς συναγωγαῖς αὐτῶν μαστιγώσουσιν ὑμᾶς·
		% 18 καὶ ἐπὶ ἡγεμόνας δὲ καὶ βασιλεῖς ἀχθήσεσθε ἕνεκεν ἐμοῦ εἰς μαρτύριον αὐτοῖς καὶ τοῖς ἔθνεσιν.
		% 19 ὅταν δὲ παραδῶσιν ὑμᾶς, μὴ μεριμνήσητε πῶς ἢ τί λαλήσητε· δοθήσεται γὰρ ὑμῖν ἐν ἐκείνῃ τῇ ὥρᾳ τί λαλήσητε·
		% 20 οὐ γὰρ ὑμεῖς ἐστε οἱ λαλοῦντες ἀλλὰ τὸ πνεῦμα τοῦ πατρὸς ὑμῶν τὸ λαλοῦν ἐν ὑμῖν.
		% 21 παραδώσει δὲ ἀδελφὸς ἀδελφὸν εἰς θάνατον καὶ πατὴρ τέκνον, καὶ ἐπαναστήσονται τέκνα ἐπὶ γονεῖς καὶ θανατώσουσιν αὐτούς.
		% 22 καὶ ἔσεσθε μισούμενοι ὑπὸ πάντων διὰ τὸ ὄνομά μου· ὁ δὲ} \hyperref[HUPOMENO]{\grl{ὑπομείνας}} \gr{εἰς τέλος οὗτος σωθήσεται.
		% 23 ὅταν δὲ διώκωσιν ὑμᾶς ἐν τῇ πόλει ταύτῃ, φεύγετε εἰς τὴν ἑτέραν· ἀμὴν γὰρ λέγω ὑμῖν, οὐ μὴ τελέσητε τὰς πόλεις τοῦ Ἰσραὴλ ἕως ἂν ἔλθῃ ὁ υἱὸς τοῦ ἀνθρώπου.
		% }
		
		% \gr{
		% 24 Οὐκ ἔστιν μαθητὴς ὑπὲρ τὸν διδάσκαλον οὐδὲ δοῦλος ὑπὲρ τὸν κύριον αὐτοῦ.
		% 25 ἀρκετὸν τῷ μαθητῇ ἵνα γένηται ὡς ὁ διδάσκαλος αὐτοῦ, καὶ ὁ δοῦλος ὡς ὁ κύριος αὐτοῦ. εἰ τὸν οἰκοδεσπότην Βεελζεβοὺλ ἐπεκάλεσαν, πόσῳ μᾶλλον τοὺς οἰκιακοὺς αὐτοῦ.
		% }
		
	% % ----------------------
	% \subsection{Mark 13:9--13}
	% % ----------------------
		% \label{LCHRIST-082-MK}
		
		% \gr{
		% 9 βλέπετε δὲ ὑμεῖς ἑαυτούς· παραδώσουσιν ὑμᾶς εἰς συνέδρια καὶ εἰς συναγωγὰς δαρήσεσθε καὶ ἐπὶ ἡγεμόνων καὶ βασιλέων σταθήσεσθε ἕνεκεν ἐμοῦ εἰς μαρτύριον αὐτοῖς.
		% 10 καὶ εἰς πάντα τὰ ἔθνη πρῶτον δεῖ κηρυχθῆναι τὸ εὐαγγέλιον.
		% 11 καὶ ὅταν ἄγωσιν ὑμᾶς παραδιδόντες, μὴ προμεριμνᾶτε τί λαλήσητε, ἀλλ᾽ ὃ ἐὰν δοθῇ ὑμῖν ἐν ἐκείνῃ τῇ ὥρᾳ τοῦτο λαλεῖτε, οὐ γάρ ἐστε ὑμεῖς οἱ λαλοῦντες ἀλλὰ τὸ πνεῦμα τὸ ἅγιον.
		% 12 καὶ παραδώσει ἀδελφὸς ἀδελφὸν εἰς θάνατον καὶ πατὴρ τέκνον, καὶ ἐπαναστήσονται τέκνα ἐπὶ γονεῖς καὶ θανατώσουσιν αὐτούς·
		% 13 καὶ ἔσεσθε μισούμενοι ὑπὸ πάντων διὰ τὸ ὄνομά μου. ὁ δὲ ὑπομείνας εἰς τέλος οὗτος σωθήσεται.
		% }
		
	% % ----------------------
	% \subsection{Luke 21:12--17}
	% % ----------------------
		% \label{LCHRIST-082-L}
		
		% \gr{
		% 12 Πρὸ δὲ τούτων πάντων ἐπιβαλοῦσιν ἐφ᾽ ὑμᾶς τὰς χεῖρας αὐτῶν καὶ διώξουσιν, παραδιδόντες εἰς τὰς συναγωγὰς καὶ φυλακάς, ἀπαγομένους ἐπὶ βασιλεῖς καὶ ἡγεμόνας ἕνεκεν τοῦ ὀνόματός μου·
		% 13 ἀποβήσεται ὑμῖν εἰς μαρτύριον.
		% 14 θέτε οὖν ἐν ταῖς καρδίαις ὑμῶν μὴ προμελετᾶν ἀπολογηθῆναι,
		% 15 ἐγὼ γὰρ δώσω ὑμῖν στόμα καὶ σοφίαν ᾗ οὐ δυνήσονται ἀντιστῆναι ἢ ἀντειπεῖν ἅπαντες οἱ ἀντικείμενοι ὑμῖν.
		% 16 παραδοθήσεσθε δὲ καὶ ὑπὸ γονέων καὶ ἀδελφῶν καὶ συγγενῶν καὶ φίλων, καὶ θανατώσουσιν ἐξ ὑμῶν,
		% 17 καὶ ἔσεσθε μισούμενοι ὑπὸ πάντων διὰ τὸ ὄνομά μου.
		% }
		
% % -----------------------------------------------------------
% \section{83 | Confessing Christ before Men}
% % -----------------------------------------------------------
	% \label{LCHRIST-083}
	
	% % ----------------------
	% \subsection{Matthew 10:32--33; 12:32}
	% % ----------------------
		% \label{LCHRIST-083-MT}
		
		% \gr{
		% 32 Πᾶς οὖν ὅστις ὁμολογήσει ἐν ἐμοὶ}
			% \footnote{%
				% \label{LCHRIST-083-MT-en-emoi}%
				% Here and below there is this odd \gr{ὁμολογέω ἐν ἐμοὶ} construction. According to BDAG, at meaning \#4b for \gr{ὁμολογέω} (708b, PDF 397), ``Instead of accusative of person we may have \gr{ἐν τινί} \textit{confess someone}, an Aramaism.'' The entry then glosses Matthew 10:32a as ``\textit{whoever confesses me before people}.
				% }
		% \gr{ἔμπροσθεν τῶν ἀνθρώπων, ὁμολογήσω κἀγὼ ἐν αὐτῷ ἔμπροσθεν τοῦ πατρός μου τοῦ ἐν οὐρανοῖς·
		% 33 ὅστις δ᾽ ἂν ἀρνήσηταί με ἔμπροσθεν τῶν ἀνθρώπων, ἀρνήσομαι κἀγὼ αὐτὸν ἔμπροσθεν τοῦ πατρός μου τοῦ ἐν οὐρανοῖς.
		% }
		
		% \gr{
		% 32 καὶ ὃς ἐὰν εἴπῃ λόγον κατὰ τοῦ υἱοῦ τοῦ ἀνθρώπου, ἀφεθήσεται αὐτῷ· ὃς δ᾽ ἂν εἴπῃ κατὰ τοῦ πνεύματος τοῦ ἁγίου, οὐκ ἀφεθήσεται αὐτῷ οὔτε ἐν τούτῳ τῷ αἰῶνι οὔτε ἐν τῷ μέλλοντι.
		% }
		
	% % ----------------------
	% \subsection{Luke 12:8--12}
	% % ----------------------
		% \label{LCHRIST-083-L}
		
		% \gr{
		% 8 Λέγω δὲ ὑμῖν, πᾶς ὃς ἂν ὁμολογήσῃ ἐν ἐμοὶ ἔμπροσθεν τῶν ἀνθρώπων, καὶ ὁ υἱὸς τοῦ ἀνθρώπου ὁμολογήσει ἐν αὐτῷ ἔμπροσθεν τῶν ἀγγέλων τοῦ θεοῦ·
		% 9 ὁ δὲ ἀρνησάμενός με ἐνώπιον τῶν ἀνθρώπων ἀπαρνηθήσεται ἐνώπιον τῶν ἀγγέλων τοῦ θεοῦ.
		% 10 καὶ πᾶς ὃς ἐρεῖ λόγον εἰς τὸν υἱὸν τοῦ ἀνθρώπου, ἀφεθήσεται αὐτῷ τῷ δὲ εἰς τὸ ἅγιον πνεῦμα βλασφημήσαντι οὐκ ἀφεθήσεται.
		% 11 ὅταν δὲ εἰσφέρωσιν ὑμᾶς ἐπὶ τὰς συναγωγὰς καὶ τὰς ἀρχὰς καὶ τὰς ἐξουσίας, μὴ μεριμνήσητε πῶς ἢ τί ἀπολογήσησθε ἢ τί εἴπητε·
		% 12 τὸ γὰρ ἅγιον πνεῦμα διδάξει ὑμᾶς ἐν αὐτῇ τῇ ὥρᾳ ἃ δεῖ εἰπεῖν.
		% }
		
% % -----------------------------------------------------------
% \section{84 | Whom to Fear}
% % -----------------------------------------------------------
	% \label{LCHRIST-084}
	
	% % ----------------------
	% \subsection{Matthew 10:26--33}
	% % ----------------------
		% \label{LCHRIST-084-MT}
		
		% \gr{
		% 26 Μὴ οὖν φοβηθῆτε αὐτούς· οὐδὲν γάρ ἐστιν κεκαλυμμένον ὃ οὐκ ἀποκαλυφθήσεται, καὶ κρυπτὸν ὃ οὐ γνωσθήσεται.
		% 27 ὃ λέγω ὑμῖν ἐν τῇ σκοτίᾳ, εἴπατε ἐν τῷ φωτί· καὶ ὃ εἰς τὸ οὖς ἀκούετε, κηρύξατε ἐπὶ τῶν δωμάτων.
		% 28 καὶ μὴ φοβεῖσθε ἀπὸ τῶν ἀποκτεννόντων τὸ σῶμα τὴν δὲ ψυχὴν μὴ δυναμένων ἀποκτεῖναι· φοβεῖσθε δὲ μᾶλλον τὸν δυνάμενον καὶ ψυχὴν καὶ σῶμα ἀπολέσαι ἐν γεέννῃ.
		% 29 οὐχὶ δύο στρουθία ἀσσαρίου πωλεῖται; καὶ ἓν ἐξ αὐτῶν οὐ πεσεῖται ἐπὶ τὴν γῆν ἄνευ τοῦ πατρὸς ὑμῶν.
		% 30 ὑμῶν δὲ καὶ αἱ τρίχες τῆς κεφαλῆς πᾶσαι ἠριθμημέναι εἰσίν.
		% 31 μὴ οὖν φοβεῖσθε· πολλῶν στρουθίων διαφέρετε ὑμεῖς.
		% }
		
		% \gr{
		% 32 Πᾶς οὖν ὅστις ὁμολογήσει ἐν ἐμοὶ ἔμπροσθεν τῶν ἀνθρώπων, ὁμολογήσω κἀγὼ ἐν αὐτῷ ἔμπροσθεν τοῦ πατρός μου τοῦ ἐν οὐρανοῖς·
		% 33 ὅστις δ᾽ ἂν ἀρνήσηταί με ἔμπροσθεν τῶν ἀνθρώπων, ἀρνήσομαι κἀγὼ αὐτὸν ἔμπροσθεν τοῦ πατρός μου τοῦ ἐν οὐρανοῖς.
		% }
		
	% % ----------------------
	% \subsection{Luke 12:4--7}
	% % ----------------------
		% \label{LCHRIST-084-L}
		
		% \gr{
		% 4 Λέγω δὲ ὑμῖν τοῖς φίλοις μου, μὴ φοβηθῆτε ἀπὸ τῶν ἀποκτεινόντων τὸ σῶμα καὶ μετὰ ταῦτα μὴ ἐχόντων περισσότερόν τι ποιῆσαι.
		% 5 ὑποδείξω δὲ ὑμῖν τίνα φοβηθῆτε· φοβήθητε τὸν μετὰ τὸ ἀποκτεῖναι ἔχοντα ἐξουσίαν ἐμβαλεῖν εἰς τὴν γέενναν· ναί, λέγω ὑμῖν, τοῦτον φοβήθητε.
		% 6 οὐχὶ πέντε στρουθία πωλοῦνται ἀσσαρίων δύο; καὶ ἓν ἐξ αὐτῶν οὐκ ἔστιν ἐπιλελησμένον ἐνώπιον τοῦ θεοῦ.
		% 7 ἀλλὰ καὶ αἱ τρίχες τῆς κεφαλῆς ὑμῶν πᾶσαι ἠρίθμηνται· μὴ φοβεῖσθε· πολλῶν στρουθίων διαφέρετε.
		% }
		
% % -----------------------------------------------------------
% \section{85 | Not Peace But a Sword}
% % -----------------------------------------------------------
	% \label{LCHRIST-085}
	
	% % ----------------------
	% \subsection{Matthew 10:34--39}
	% % ----------------------
		% \label{LCHRIST-085-MT}
		
		% \gr{
		% 34 Μὴ νομίσητε ὅτι ἦλθον βαλεῖν εἰρήνην ἐπὶ τὴν γῆν· οὐκ ἦλθον βαλεῖν εἰρήνην ἀλλὰ μάχαιραν.
		% 35 ἦλθον γὰρ διχάσαι ἄνθρωπον κατὰ τοῦ πατρὸς αὐτοῦ καὶ θυγατέρα κατὰ τῆς μητρὸς αὐτῆς καὶ νύμφην κατὰ τῆς πενθερᾶς αὐτῆς,}
			% \footnote{%
				% \label{LCHRIST-085-dichasai}%
				% Some translations of this verse:
					% \begin{compactdesc}
						% \item[ASV] For I came to set a man at variance against his father, and the daughter against her mother, and the daughter in law against her mother in law:
						% \item[NIV] For I have come to turn `` `a man against his father, a daughter against her mother, a daughter-in-law against her mother-in-law---
						% \item[NRS] For I have come to set a man against his father, and a daughter against her mother, and a daughter-in-law against her mother-in-law;
					% \end{compactdesc}
				% These translations emphasize that Christ is turning a person against members of his family, and then verse 36 follows up on that, but I think it's also possible to take a more internalist interpretation of these verses, where you can understand the members of the family---father, mother, mother-in-law---as traditions or beliefs or practices that you ought to reject or change. Christ comes to set us in opposition to those things, and the long-term residents inside our heads, like beliefs or cherished actions or something, might be our enemies. Some of them will be, at least.
				
				% I think this interpretation goes well with verse 39 in particular: Christ is saying, you have to choose which life you want, the one you've always had and are comfortable with, or the one you could have by following me. A person who \gr{ὁ ἀπολέσας τὴν ψυχὴν αὐτοῦ} breaks down the weaknesses and sins in his life and tries to turn them into strengths through Christ and the atonement; that person does not hold on to the past or to what is comfortable or easy, but instead seeks for growth.
				% }
		% \gr{36 καὶ ἐχθροὶ τοῦ ἀνθρώπου οἱ οἰκιακοὶ αὐτοῦ.
		% 37 ὁ φιλῶν πατέρα ἢ μητέρα ὑπὲρ ἐμὲ οὐκ ἔστιν μου ἄξιος· καὶ ὁ φιλῶν υἱὸν ἢ θυγατέρα ὑπὲρ ἐμὲ οὐκ ἔστιν μου ἄξιος·}
			% \footnote{%
				% \label{LCHRIST-085-huper}%
				% The emphasis here seems to be on our loving something else \gr{ὑπὲρ}---more than or beyond---Christ.
				% }
		% \gr{38 καὶ ὃς οὐ λαμβάνει τὸν σταυρὸν αὐτοῦ καὶ ἀκολουθεῖ ὀπίσω μου, οὐκ ἔστιν μου ἄξιος.
		% 39 ὁ εὑρὼν τὴν ψυχὴν αὐτοῦ ἀπολέσει αὐτήν, καὶ ὁ ἀπολέσας τὴν ψυχὴν αὐτοῦ ἕνεκεν ἐμοῦ εὑρήσει αὐτήν.
		% }
		
	% % ----------------------
	% \subsection{Luke 12:49--53; 14:26--27}
	% % ----------------------
		% \label{LCHRIST-085-L}
		
		% \gr{
		% 49 Πῦρ ἦλθον βαλεῖν ἐπὶ τὴν γῆν, καὶ τί θέλω εἰ ἤδη ἀνήφθη;
		% 50 βάπτισμα δὲ ἔχω βαπτισθῆναι, καὶ πῶς συνέχομαι ἕως ὅτου τελεσθῇ.
		% 51 δοκεῖτε ὅτι εἰρήνην παρεγενόμην δοῦναι ἐν τῇ γῇ; οὐχί, λέγω ὑμῖν, ἀλλ᾽ ἢ διαμερισμόν.
		% 52 ἔσονται γὰρ ἀπὸ τοῦ νῦν πέντε ἐν ἑνὶ οἴκῳ διαμεμερισμένοι, τρεῖς ἐπὶ δυσὶν καὶ δύο ἐπὶ τρισίν,
		% 53 διαμερισθήσονται πατὴρ ἐπὶ υἱῷ καὶ υἱὸς ἐπὶ πατρί, μήτηρ ἐπὶ θυγατέρα καὶ θυγάτηρ ἐπὶ τὴν μητέρα, πενθερὰ ἐπὶ τὴν νύμφην αὐτῆς καὶ νύμφη ἐπὶ τὴν πενθεράν.
		% }
		
		% \gr{
		% 26 Εἴ τις ἔρχεται πρός με καὶ οὐ μισεῖ τὸν πατέρα ἑαυτοῦ καὶ τὴν μητέρα καὶ τὴν γυναῖκα καὶ τὰ τέκνα καὶ τοὺς ἀδελφοὺς καὶ τὰς ἀδελφάς, ἔτι τε καὶ τὴν ψυχὴν ἑαυτοῦ, οὐ δύναται εἶναί μου μαθητής.
		% 27 ὅστις οὐ βαστάζει τὸν σταυρὸν ἑαυτοῦ καὶ ἔρχεται ὀπίσω μου, οὐ δύναται εἶναί μου μαθητής.
		% }
		
% % -----------------------------------------------------------
% \section{86 | Conditions of Discipleship}
% % -----------------------------------------------------------
	% \label{LCHRIST-086}
	
	% % ----------------------
	% \subsection{Matthew 10:37--38}
	% % ----------------------
		% \label{LCHRIST-086-MT}
		
		% \gr{
		% 37 ὁ φιλῶν πατέρα ἢ μητέρα ὑπὲρ ἐμὲ οὐκ ἔστιν μου ἄξιος· καὶ ὁ φιλῶν υἱὸν ἢ θυγατέρα ὑπὲρ ἐμὲ οὐκ ἔστιν μου ἄξιος·
		% 38 καὶ ὃς οὐ λαμβάνει τὸν σταυρὸν αὐτοῦ καὶ ἀκολουθεῖ ὀπίσω μου, οὐκ ἔστιν μου ἄξιος.
		% }
		
	% % ----------------------
	% \subsection{Luke 14:26--33}
	% % ----------------------
		% \label{LCHRIST-086-L}
		
		% \gr{
		% 26 Εἴ τις ἔρχεται πρός με καὶ οὐ μισεῖ τὸν πατέρα ἑαυτοῦ καὶ τὴν μητέρα καὶ τὴν γυναῖκα καὶ τὰ τέκνα καὶ τοὺς ἀδελφοὺς καὶ τὰς ἀδελφάς, ἔτι τε καὶ τὴν ψυχὴν ἑαυτοῦ, οὐ δύναται εἶναί μου μαθητής.
		% 27 ὅστις οὐ βαστάζει τὸν σταυρὸν ἑαυτοῦ καὶ ἔρχεται ὀπίσω μου, οὐ δύναται εἶναί μου μαθητής.
		% 28 τίς γὰρ ἐξ ὑμῶν θέλων πύργον οἰκοδομῆσαι οὐχὶ πρῶτον καθίσας ψηφίζει τὴν δαπάνην, εἰ ἔχει εἰς ἀπαρτισμόν;
		% 29 ἵνα μήποτε θέντος αὐτοῦ θεμέλιον καὶ μὴ ἰσχύοντος ἐκτελέσαι πάντες οἱ θεωροῦντες ἄρξωνται αὐτῷ ἐμπαίζειν
		% 30 λέγοντες ὅτι Οὗτος ὁ ἄνθρωπος ἤρξατο οἰκοδομεῖν καὶ οὐκ ἴσχυσεν ἐκτελέσαι.
		% 31 ἢ τίς βασιλεὺς πορευόμενος ἑτέρῳ βασιλεῖ συμβαλεῖν εἰς πόλεμον οὐχὶ καθίσας πρῶτον βουλεύσεται εἰ δυνατός ἐστιν ἐν δέκα χιλιάσιν ὑπαντῆσαι τῷ μετὰ εἴκοσι χιλιάδων ἐρχομένῳ ἐπ᾽ αὐτόν;
		% 32 εἰ δὲ μήγε, ἔτι αὐτοῦ πόρρω ὄντος πρεσβείαν ἀποστείλας ἐρωτᾷ τὰ πρὸς εἰρήνην.
		% 33 οὕτως οὖν πᾶς ἐξ ὑμῶν ὃς οὐκ ἀποτάσσεται πᾶσιν τοῖς ἑαυτοῦ ὑπάρχουσιν οὐ δύναται εἶναί μου μαθητής.
		% }
		
% % -----------------------------------------------------------
% \section{87 | Rewards of Discipleship}
% % -----------------------------------------------------------
	% \label{LCHRIST-087}
	
	% % ----------------------
	% \subsection{Matthew 10:40--11:1}
	% % ----------------------
		% \label{LCHRIST-087-MT}
		
		% \gr{
		% 40 Ὁ δεχόμενος ὑμᾶς ἐμὲ δέχεται, καὶ ὁ ἐμὲ δεχόμενος δέχεται τὸν ἀποστείλαντά με.
		% 41 ὁ δεχόμενος προφήτην εἰς ὄνομα προφήτου μισθὸν προφήτου λήμψεται, καὶ ὁ δεχόμενος δίκαιον εἰς ὄνομα δικαίου μισθὸν δικαίου λήμψεται.
		% 42 καὶ ὃς ἂν ποτίσῃ ἕνα τῶν μικρῶν τούτων ποτήριον ψυχροῦ μόνον εἰς ὄνομα μαθητοῦ, ἀμὴν λέγω ὑμῖν, οὐ μὴ ἀπολέσῃ τὸν μισθὸν αὐτοῦ.
		% }
		
		% \gr{
		% 1 Καὶ ἐγένετο ὅτε ἐτέλεσεν ὁ Ἰησοῦς διατάσσων τοῖς δώδεκα μαθηταῖς αὐτοῦ, μετέβη ἐκεῖθεν τοῦ διδάσκειν καὶ κηρύσσειν ἐν ταῖς πόλεσιν αὐτῶν.
		% }
		
	% % ----------------------
	% \subsection{Mark 9:41}
	% % ----------------------
		% \label{LCHRIST-087-MK}
		
		% \gr{
		% 41 Ὃς γὰρ ἂν ποτίσῃ ὑμᾶς ποτήριον ὕδατος ἐν ὀνόματι ὅτι χριστοῦ ἐστε, ἀμὴν λέγω ὑμῖν ὅτι οὐ μὴ ἀπολέσῃ τὸν μισθὸν αὐτοῦ.
		% }
		
% % -----------------------------------------------------------
% \section{88 | Messengers from John the Baptist}
% % -----------------------------------------------------------
	% \label{LCHRIST-088}
	
	% % ----------------------
	% \subsection{Matthew 11:2--19}
	% % ----------------------
		% \label{LCHRIST-088-MT}
		
		% \gr{
		% 2 Ὁ δὲ Ἰωάννης ἀκούσας ἐν τῷ δεσμωτηρίῳ τὰ ἔργα τοῦ χριστοῦ πέμψας διὰ τῶν μαθητῶν αὐτοῦ
		% 3 εἶπεν αὐτῷ· Σὺ εἶ ὁ ἐρχόμενος ἢ ἕτερον προσδοκῶμεν;}
			% \footnote{%
				% \label{LCHRIST-088-su}%
				% Why would John need to ask this question---why does he need to ask if Christ is the one that is coming? John had been present at the \hyperref[LCHRIST-025]{baptism of Christ}, after all, and it seems like he would have heard the voice of the Father speaking after the baptism---though see \hyperref[LCHRIST-025-su]{this note}, which points out that the Mark and Luke accounts both say ``\textit{You} are my son.'' Is it possible that John didn't hear the voice, and so we looking for some more information? According to \hyperref[DC93-6]{DC 93:6--18}, especially \hyperref[DC93-15]{verse 15}, John \textit{did} hear the voice. So I'm not sure why he would have needed to ask. It could be that he was just passing through a very difficult time in his life---he's in prison, many people hate him, and so on---and so he needed some reassurance; he needed some assurance that what he was doing was the right thing. Christ's response is, ``Look at all these evidences.'' Pretty powerful passage I think.
				% }
		% \gr{4 καὶ ἀποκριθεὶς ὁ Ἰησοῦς εἶπεν αὐτοῖς· Πορευθέντες ἀπαγγείλατε Ἰωάννῃ ἃ ἀκούετε καὶ βλέπετε·
		% 5 τυφλοὶ ἀναβλέπουσιν καὶ χωλοὶ περιπατοῦσιν, λεπροὶ καθαρίζονται καὶ κωφοὶ ἀκούουσιν, καὶ νεκροὶ ἐγείρονται καὶ πτωχοὶ εὐαγγελίζονται·
		% 6 καὶ μακάριός ἐστιν ὃς ἐὰν μὴ} \hyperref[SKANDALIZO]{\grl{σκανδαλισθῇ}} \gr{ἐν ἐμοί.
		% }
		
		% \gr{
		% 7 Τούτων δὲ πορευομένων ἤρξατο ὁ Ἰησοῦς λέγειν τοῖς ὄχλοις περὶ Ἰωάννου· Τί ἐξήλθατε εἰς τὴν ἔρημον θεάσασθαι; κάλαμον ὑπὸ ἀνέμου σαλευόμενον;
		% 8 ἀλλὰ τί ἐξήλθατε ἰδεῖν; ἄνθρωπον ἐν μαλακοῖς ἠμφιεσμένον; ἰδοὺ οἱ τὰ μαλακὰ φοροῦντες ἐν τοῖς οἴκοις τῶν βασιλέων εἰσίν.
		% 9 ἀλλὰ τί ἐξήλθατε; προφήτην ἰδεῖν; ναί, λέγω ὑμῖν, καὶ περισσότερον προφήτου.
		% 10 οὗτός ἐστιν περὶ οὗ γέγραπται· Ἰδοὺ ἐγὼ ἀποστέλλω τὸν ἄγγελόν μου πρὸ προσώπου σου, ὃς κατασκευάσει τὴν ὁδόν σου ἔμπροσθέν σου.
		% 11 ἀμὴν λέγω ὑμῖν, οὐκ ἐγήγερται ἐν γεννητοῖς γυναικῶν μείζων Ἰωάννου τοῦ βαπτιστοῦ· ὁ δὲ μικρότερος ἐν τῇ βασιλείᾳ τῶν οὐρανῶν μείζων αὐτοῦ ἐστιν.
		% 12 ἀπὸ δὲ τῶν ἡμερῶν Ἰωάννου τοῦ βαπτιστοῦ ἕως ἄρτι ἡ βασιλεία τῶν οὐρανῶν βιάζεται, καὶ βιασταὶ ἁρπάζουσιν αὐτήν.
		% 13 πάντες γὰρ οἱ προφῆται καὶ ὁ νόμος ἕως Ἰωάννου ἐπροφήτευσαν·
		% 14 καὶ εἰ θέλετε δέξασθαι, αὐτός ἐστιν Ἠλίας ὁ μέλλων ἔρχεσθαι.
		% 15 ὁ ἔχων ὦτα ἀκουέτω.
		% }
		
		% \gr{
		% 16 Τίνι δὲ ὁμοιώσω τὴν γενεὰν ταύτην; ὁμοία ἐστὶν παιδίοις καθημένοις ἐν ταῖς ἀγοραῖς ἃ προσφωνοῦντα τοῖς ἑτέροις
		% 17 λέγουσιν· Ηὐλήσαμεν ὑμῖν καὶ οὐκ ὠρχήσασθε· ἐθρηνήσαμεν καὶ οὐκ ἐκόψασθε·}
			% \footnote{%
				% \label{LCHRIST-088-aulesamen}%
				% I read this verse as saying: the people said to John, we're treating you in the ways we normally treat others, but you aren't responding like others do; we aren't able to control you like we do other people. You do your own thing and don't listen whether we're playing the flute or anything else. Usually when we play the flute, the people dance, but you don't do that. I think that we ought to be more like John here---we make the choices about how we spend our time, how we act, and what we do. We don't listen when the world is piping the flute or singing a song to try to influence us.
				% }
		% \gr{18 ἦλθεν γὰρ Ἰωάννης μήτε ἐσθίων μήτε πίνων, καὶ λέγουσιν· Δαιμόνιον ἔχει·
		% 19 ἦλθεν ὁ υἱὸς τοῦ ἀνθρώπου ἐσθίων καὶ πίνων, καὶ λέγουσιν· Ἰδοὺ ἄνθρωπος φάγος καὶ οἰνοπότης, τελωνῶν φίλος καὶ ἁμαρτωλῶν. καὶ ἐδικαιώθη ἡ σοφία ἀπὸ τῶν ἔργων αὐτῆς.
		% }
		
	% % ----------------------
	% \subsection{Luke 7:18--23}
	% % ----------------------
		% \label{LCHRIST-088-L}
		
		% \gr{
		% 18 Καὶ ἀπήγγειλαν Ἰωάννῃ οἱ μαθηταὶ αὐτοῦ περὶ πάντων τούτων. καὶ προσκαλεσάμενος δύο τινὰς τῶν μαθητῶν αὐτοῦ ὁ Ἰωάννης
		% 19 ἔπεμψεν πρὸς τὸν κύριον λέγων· Σὺ εἶ ὁ ἐρχόμενος ἢ ἄλλον προσδοκῶμεν;
		% 20 παραγενόμενοι δὲ πρὸς αὐτὸν οἱ ἄνδρες εἶπαν· Ἰωάννης ὁ βαπτιστὴς ἀπέστειλεν ἡμᾶς πρὸς σὲ λέγων· Σὺ εἶ ὁ ἐρχόμενος ἢ ἄλλον προσδοκῶμεν;
		% 21 ἐν ἐκείνῃ τῇ ὥρᾳ ἐθεράπευσεν πολλοὺς ἀπὸ νόσων καὶ μαστίγων καὶ πνευμάτων πονηρῶν, καὶ τυφλοῖς πολλοῖς ἐχαρίσατο βλέπειν.
		% 22 καὶ ἀποκριθεὶς εἶπεν αὐτοῖς· Πορευθέντες ἀπαγγείλατε Ἰωάννῃ ἃ εἴδετε καὶ ἠκούσατε· τυφλοὶ ἀναβλέπουσιν, χωλοὶ περιπατοῦσιν, λεπροὶ καθαρίζονται, κωφοὶ ἀκούουσιν, νεκροὶ ἐγείρονται, πτωχοὶ εὐαγγελίζονται·
		% 23 καὶ μακάριός ἐστιν ὃς ἐὰν μὴ} \hyperref[SKANDALIZO]{\grl{σκανδαλισθῇ}} \gr{ἐν ἐμοί.
		% }
		
% % -----------------------------------------------------------
% \section{89 | Christ's Testimony of John}
% % -----------------------------------------------------------
	% \label{LCHRIST-089}
	
	% % ----------------------
	% \subsection{Luke 7:24--35}
	% % ----------------------
		% \label{LCHRIST-089-L}
		
		% \gr{
		% 24 Ἀπελθόντων δὲ τῶν ἀγγέλων Ἰωάννου ἤρξατο λέγειν πρὸς τοὺς ὄχλους περὶ Ἰωάννου· Τί ἐξήλθατε εἰς τὴν ἔρημον θεάσασθαι; κάλαμον ὑπὸ ἀνέμου σαλευόμενον;
		% 25 ἀλλὰ τί ἐξήλθατε ἰδεῖν; ἄνθρωπον ἐν μαλακοῖς ἱματίοις ἠμφιεσμένον; ἰδοὺ οἱ ἐν ἱματισμῷ ἐνδόξῳ καὶ τρυφῇ ὑπάρχοντες ἐν τοῖς βασιλείοις εἰσίν.
		% 26 ἀλλὰ τί ἐξήλθατε ἰδεῖν; προφήτην; ναί, λέγω ὑμῖν, καὶ περισσότερον προφήτου.
		% 27 οὗτός ἐστιν περὶ οὗ γέγραπται· Ἰδοὺ ἀποστέλλω τὸν ἄγγελόν μου πρὸ προσώπου σου, ὃς κατασκευάσει τὴν ὁδόν σου ἔμπροσθέν σου.
		% 28 λέγω ὑμῖν, μείζων ἐν γεννητοῖς γυναικῶν Ἰωάννου οὐδείς ἐστιν· ὁ δὲ μικρότερος ἐν τῇ βασιλείᾳ τοῦ θεοῦ μείζων αὐτοῦ ἐστιν.
		% 29 (καὶ πᾶς ὁ λαὸς ἀκούσας καὶ οἱ τελῶναι ἐδικαίωσαν τὸν θεόν, βαπτισθέντες τὸ βάπτισμα Ἰωάννου·
		% 30 οἱ δὲ Φαρισαῖοι καὶ οἱ νομικοὶ τὴν βουλὴν τοῦ θεοῦ ἠθέτησαν εἰς ἑαυτούς, μὴ βαπτισθέντες ὑπ᾽ αὐτοῦ.)
		% }
		
		% \gr{
		% 31 Τίνι οὖν ὁμοιώσω τοὺς ἀνθρώπους τῆς γενεᾶς ταύτης, καὶ τίνι εἰσὶν ὅμοιοι;
		% 32 ὅμοιοί εἰσιν παιδίοις τοῖς ἐν ἀγορᾷ καθημένοις καὶ προσφωνοῦσιν ἀλλήλοις, ἃ λέγει· Ηὐλήσαμεν ὑμῖν καὶ οὐκ ὠρχήσασθε· ἐθρηνήσαμεν καὶ οὐκ ἐκλαύσατε·
		% 33 ἐλήλυθεν γὰρ Ἰωάννης ὁ βαπτιστὴς μὴ ἐσθίων ἄρτον μήτε πίνων οἶνον, καὶ λέγετε· Δαιμόνιον ἔχει·
		% 34 ἐλήλυθεν ὁ υἱὸς τοῦ ἀνθρώπου ἐσθίων καὶ πίνων, καὶ λέγετε· Ἰδοὺ ἄνθρωπος φάγος καὶ οἰνοπότης, φίλος τελωνῶν καὶ ἁμαρτωλῶν.
		% 35 καὶ ἐδικαιώθη ἡ σοφία ἀπὸ πάντων τῶν τέκνων αὐτῆς.
		% }
		
	% % ----------------------
	% \subsection{John 3:22--30}
	% % ----------------------
		% \label{LCHRIST-089-J}
		
		% \gr{
		% 22 Μετὰ ταῦτα ἦλθεν ὁ Ἰησοῦς καὶ οἱ μαθηταὶ αὐτοῦ εἰς τὴν Ἰουδαίαν γῆν, καὶ ἐκεῖ διέτριβεν μετ᾽ αὐτῶν καὶ ἐβάπτιζεν.
		% 23 ἦν δὲ καὶ ὁ Ἰωάννης βαπτίζων ἐν Αἰνὼν ἐγγὺς τοῦ Σαλείμ, ὅτι ὕδατα πολλὰ ἦν ἐκεῖ, καὶ παρεγίνοντο καὶ ἐβαπτίζοντο·
		% 24 οὔπω γὰρ ἦν βεβλημένος εἰς τὴν φυλακὴν ὁ Ἰωάννης.
		% }
		
		% \gr{
		% 25 Ἐγένετο οὖν ζήτησις ἐκ τῶν μαθητῶν Ἰωάννου μετὰ Ἰουδαίου περὶ καθαρισμοῦ.
		% 26 καὶ ἦλθον πρὸς τὸν Ἰωάννην καὶ εἶπαν αὐτῷ· Ῥαββί, ὃς ἦν μετὰ σοῦ πέραν τοῦ Ἰορδάνου, ᾧ σὺ μεμαρτύρηκας, ἴδε οὗτος βαπτίζει καὶ πάντες ἔρχονται πρὸς αὐτόν.
		% 27 ἀπεκρίθη Ἰωάννης καὶ εἶπεν· Οὐ δύναται ἄνθρωπος λαμβάνειν οὐδὲ ἓν ἐὰν μὴ ᾖ δεδομένον αὐτῷ ἐκ τοῦ οὐρανοῦ.
		% 28 αὐτοὶ ὑμεῖς μοι μαρτυρεῖτε ὅτι εἶπον· Οὐκ εἰμὶ ἐγὼ ὁ χριστός, ἀλλ᾽ ὅτι Ἀπεσταλμένος εἰμὶ ἔμπροσθεν ἐκείνου.
		% 29 ὁ ἔχων τὴν νύμφην νυμφίος ἐστίν· ὁ δὲ φίλος τοῦ νυμφίου ὁ ἑστηκὼς καὶ ἀκούων αὐτοῦ, χαρᾷ χαίρει διὰ τὴν φωνὴν τοῦ νυμφίου. αὕτη οὖν ἡ χαρὰ ἡ ἐμὴ πεπλήρωται.
		% 30 ἐκεῖνον δεῖ αὐξάνειν, ἐμὲ δὲ ἐλαττοῦσθαι.
		% }
		
% % -----------------------------------------------------------
% \section{90 | Law and the Kingdom of God}
% % -----------------------------------------------------------
	% \label{LCHRIST-090}
	
	% % ----------------------
	% \subsection{Matthew 11:12--13}
	% % ----------------------
		% \label{LCHRIST-090-MT}
		
		% \gr{
		% 12 ἀπὸ δὲ τῶν ἡμερῶν Ἰωάννου τοῦ βαπτιστοῦ ἕως ἄρτι ἡ βασιλεία τῶν οὐρανῶν βιάζεται,} %
			% \footnote{%
				% \label{LCHRIST-090-biazo}%
				% \gr{βιάζω} appears twice in the NT---here and at Luke 16:16, below. See BDAG 175b, PDF 128. Apparently there is some controversy over the interpretation of this word; the beginning of BDAG's entry says, ``The principal semantic problem is whether \gr{βιάζω} is used negatively (`in malam partem') or positively (`in bonam partem'), a problem compounded by the question of the functio of these verses in their literary context. In Greek literature \gr{βιάζω} is most often used in the unfavorable sense of attack or forcible constraint.'' Then it gives the definitions:
					% \begin{compactitem}
						% \item[1] \textbf{to inflict violence on, \textit{dominate, constraint}}...with \gr{βιάζω} taken as passive, \textbf{Mt 11:12} is frequently understood in the unfavorable sense \textit{the reign/kingdom of heaven is unfairly treated, is oppressed}
						% \item[2] \textbf{to gain an objective by force, \textit{use force}}...\textit{Makes its way with triumphant force} is preferred for \textbf{Mt 11:12}
						% \item[3] \textbf{go after something with enthusiasm, \textit{seek fervently, try hard}}, the sense \textit{is sought with burning zeal} is preferred by [some authors] for \textbf{Mt 11:12}
						% \item[4] \textbf{\textit{constrain (warmly)}}
					% \end{compactitem} 
				% }
		% \gr{καὶ βιασταὶ ἁρπάζουσιν αὐτήν.
		% 13 πάντες γὰρ οἱ προφῆται καὶ ὁ νόμος ἕως Ἰωάννου ἐπροφήτευσαν·
		% }
		
	% % ----------------------
	% \subsection{Luke 16:14--18}
	% % ----------------------
		% \label{LCHRIST-090-L}
		
		% \gr{
		% 14 Ἤκουον δὲ ταῦτα πάντα οἱ Φαρισαῖοι φιλάργυροι ὑπάρχοντες, καὶ ἐξεμυκτήριζον αὐτόν.
		% 15 καὶ εἶπεν αὐτοῖς· Ὑμεῖς ἐστε οἱ δικαιοῦντες ἑαυτοὺς ἐνώπιον τῶν ἀνθρώπων, ὁ δὲ θεὸς γινώσκει τὰς καρδίας ὑμῶν· ὅτι τὸ ἐν ἀνθρώποις ὑψηλὸν βδέλυγμα ἐνώπιον τοῦ θεοῦ.
		% }
		
		% \gr{
		% 16 Ὁ νόμος καὶ οἱ προφῆται μέχρι Ἰωάννου· ἀπὸ τότε ἡ βασιλεία τοῦ θεοῦ εὐαγγελίζεται καὶ πᾶς εἰς αὐτὴν βιάζεται.
		% 17 Εὐκοπώτερον δέ ἐστιν τὸν οὐρανὸν καὶ τὴν γῆν παρελθεῖν ἢ τοῦ νόμου μίαν κεραίαν πεσεῖν.
		% }
		
		% \gr{
		% 18 Πᾶς ὁ ἀπολύων τὴν γυναῖκα αὐτοῦ καὶ γαμῶν ἑτέραν μοιχεύει, καὶ ὁ ἀπολελυμένην ἀπὸ ἀνδρὸς γαμῶν μοιχεύει.
		% }
		
% % -----------------------------------------------------------
% \section{91 | On Testimony}
% % -----------------------------------------------------------
	% \label{LCHRIST-091}
	
	% % ----------------------
	% \subsection{John 3:31--36}
	% % ----------------------
		% \label{LCHRIST-091-J}
		
		% \gr{
		% 31 Ὁ ἄνωθεν ἐρχόμενος ἐπάνω πάντων ἐστίν. ὁ ὢν ἐκ τῆς γῆς ἐκ τῆς γῆς ἐστιν καὶ ἐκ τῆς γῆς λαλεῖ· ὁ ἐκ τοῦ οὐρανοῦ ἐρχόμενος ἐπάνω πάντων ἐστίν·
		% 32 ὃ ἑώρακεν καὶ ἤκουσεν τοῦτο μαρτυρεῖ, καὶ τὴν μαρτυρίαν αὐτοῦ οὐδεὶς λαμβάνει.
		% 33 ὁ λαβὼν αὐτοῦ τὴν μαρτυρίαν ἐσφράγισεν ὅτι ὁ θεὸς ἀληθής ἐστιν.
		% 34 ὃν γὰρ ἀπέστειλεν ὁ θεὸς τὰ ῥήματα τοῦ θεοῦ λαλεῖ, οὐ γὰρ ἐκ μέτρου δίδωσιν τὸ πνεῦμα.}%
			% \footnote{%
				% \label{LCHRIST-091-34}%
				% Some alternate translations of this verse:
					% \begin{compactdesc}
						% \item[ASV] For he whom God hath sent speaketh the words of God: for he giveth not the Spirit by measure.
						% \item[NASB] For He whom God has sent speaks the words of God; for He gives the Spirit without measure.
						% \item[NCV] The One whom God sent speaks the words of God, because God gives him the Spirit fully.
						% \item[NIV] For the one whom God has sent speaks the words of God, for God gives the Spirit without limit.
						% \item[NRSV] He whom God has sent speaks the words of God, for he gives the Spirit without measure.
					% \end{compactdesc}
				% I'm not totally sure how to relate it to the restoration; in the Greek there is no dative pronoun for the receiver of the Spirit which God is giving in the second half of the verse. It could be implicit, that we're stilling talking about ``the one whom God sent,'' but it could also be a more general statement of God's willingness to give the Spirit without measure (not in a measured/limited way) to anyone. It makes me think of verses like \hyperref[ALMA-18-35]{Alma 18:35}, where the text is explicit about someone receiving a ``portion'' of the Spirit. Maybe John 3:34 here is referring to our growth, that eventually we will receive a fulness? That would fall in line with restoratio thought especially as it appears in \hyperref[DC93]{DC 93}.
				% }
		% \gr{35 ὁ πατὴρ ἀγαπᾷ τὸν υἱόν, καὶ πάντα δέδωκεν ἐν τῇ χειρὶ αὐτοῦ.
		% 36 ὁ πιστεύων εἰς τὸν υἱὸν ἔχει ζωὴν αἰώνιον· ὁ δὲ ἀπειθῶν τῷ υἱῷ οὐκ ὄψεται ζωήν, ἀλλ᾽ ἡ ὀργὴ τοῦ θεοῦ μένει ἐπ᾽ αὐτόν.
		% }
		
% % -----------------------------------------------------------
% \section{92 | Woes to Unrepentant Cities}
% % -----------------------------------------------------------
	% \label{LCHRIST-092}
	
	% % ----------------------
	% \subsection{Matthew 11:20--24}
	% % ----------------------
		% \label{LCHRIST-092-MT}
		
		% \gr{
		% 20 Τότε ἤρξατο ὀνειδίζειν τὰς πόλεις ἐν αἷς ἐγένοντο αἱ πλεῖσται δυνάμεις αὐτοῦ, ὅτι οὐ μετενόησαν·
		% 21 Οὐαί σοι, Χοραζίν· οὐαί σοι, Βηθσαϊδά· ὅτι εἰ ἐν Τύρῳ καὶ Σιδῶνι ἐγένοντο αἱ δυνάμεις αἱ γενόμεναι ἐν ὑμῖν, πάλαι ἂν ἐν σάκκῳ καὶ σποδῷ μετενόησαν.
		% 22 πλὴν λέγω ὑμῖν, Τύρῳ καὶ Σιδῶνι ἀνεκτότερον ἔσται ἐν ἡμέρᾳ κρίσεως ἢ ὑμῖν.
		% 23 καὶ σύ, Καφαρναούμ, μὴ ἕως οὐρανοῦ ὑψωθήσῃ; ἕως ᾅδου καταβήσῃ· ὅτι εἰ ἐν Σοδόμοις ἐγενήθησαν αἱ δυνάμεις αἱ γενόμεναι ἐν σοί, ἔμεινεν ἂν μέχρι τῆς σήμερον.
		% 24 πλὴν λέγω ὑμῖν ὅτι γῇ Σοδόμων ἀνεκτότερον ἔσται ἐν ἡμέρᾳ κρίσεως ἢ σοί.
		% }
		
	% % ----------------------
	% \subsection{Luke 10:13--15}
	% % ----------------------
		% \label{LCHRIST-092-L}
		
		% \gr{
		% 13 Οὐαί σοι, Χοραζίν· οὐαί σοι, Βηθσαϊδά· ὅτι εἰ ἐν Τύρῳ καὶ Σιδῶνι ἐγενήθησαν αἱ δυνάμεις αἱ γενόμεναι ἐν ὑμῖν, πάλαι ἂν ἐν σάκκῳ καὶ σποδῷ καθήμενοι μετενόησαν.
		% 14 πλὴν Τύρῳ καὶ Σιδῶνι ἀνεκτότερον ἔσται ἐν τῇ κρίσει ἢ ὑμῖν.
		% 15 καὶ σύ, Καφαρναούμ, μὴ ἕως οὐρανοῦ ὑψωθήσῃ; ἕως τοῦ ᾅδου καταβιβασθήσῃ.
		% }
		
% % -----------------------------------------------------------
% \section{93 | Christ's Thanksgiving Prayer}
% % -----------------------------------------------------------
	% \label{LCHRIST-093}
	
	% % ----------------------
	% \subsection{Matthew 11:25--27}
	% % ----------------------
		% \label{LCHRIST-093-MT}
		
		% \gr{
		% 25 Ἐν ἐκείνῳ τῷ καιρῷ ἀποκριθεὶς ὁ Ἰησοῦς εἶπεν· Ἐξομολογοῦμαί σοι, πάτερ κύριε τοῦ οὐρανοῦ καὶ τῆς γῆς, ὅτι ἔκρυψας ταῦτα ἀπὸ σοφῶν καὶ συνετῶν, καὶ ἀπεκάλυψας αὐτὰ νηπίοις·
		% 26 ναί, ὁ πατήρ, ὅτι οὕτως εὐδοκία ἐγένετο ἔμπροσθέν σου.
		% 27 Πάντα μοι παρεδόθη ὑπὸ τοῦ πατρός μου, καὶ οὐδεὶς ἐπιγινώσκει τὸν υἱὸν εἰ μὴ ὁ πατήρ, οὐδὲ τὸν πατέρα τις ἐπιγινώσκει}%
			% \footnote{%
				% \label{LCHRIST-093-epiginosko}%
				% \hyperref[EPIGINOSKO]{BDAG} lists this verse as meaning \#1b, where \gr{ἐπιγινώσκω} is basically equivalent to \gr{γινώσκω}. In CGL the uses in this verse would probably fall under meaning \#1, ``understand the identity of, \textbf{recognise}.'' It's not the recognizing that happens when you see someone you haven't seen for a long time and recognize who they are by memory---I think this has to be more like, the Son is able to recognize the Father because he's the one that knows the Father and so is the only one who would be in a position to recognize him. But that doesn't really explain the first use, where it says \gr{οὐδεὶς ἐπιγινώσκει τὸν υἱὸν εἰ μὴ ὁ πατήρ}. Maybe the \gr{ἐπι}- has more force there to create more emphasis, to say that no one recognizes the Son \textit{as} the Son except the Father? I'm not really sure.
				% }
		% \gr{εἰ μὴ ὁ υἱὸς καὶ ᾧ ἐὰν βούληται ὁ υἱὸς ἀποκαλύψαι.
		% }
		
% % -----------------------------------------------------------
% \section{94 | Healing at the Pool of Bethesda}
% % -----------------------------------------------------------
	% \label{LCHRIST-094}
	
	% % ----------------------
	% \subsection{John 5:2--18}
	% % ----------------------
		% \label{LCHRIST-094-J}
		
		% \gr{
		% 2 ἔστιν δὲ ἐν τοῖς Ἱεροσολύμοις ἐπὶ τῇ προβατικῇ κολυμβήθρα ἡ ἐπιλεγομένη Ἑβραϊστὶ Βηθεσδά, πέντε στοὰς ἔχουσα·
		% 3 ἐν ταύταις κατέκειτο πλῆθος τῶν ἀσθενούντων, τυφλῶν, χωλῶν, ξηρῶν.
		% 5 ἦν δέ τις ἄνθρωπος ἐκεῖ τριάκοντα ὀκτὼ ἔτη ἔχων ἐν τῇ ἀσθενείᾳ αὐτοῦ·
		% 6 τοῦτον ἰδὼν ὁ Ἰησοῦς κατακείμενον, καὶ γνοὺς ὅτι πολὺν ἤδη χρόνον ἔχει, λέγει αὐτῷ· Θέλεις ὑγιὴς γενέσθαι;
		% 7 ἀπεκρίθη αὐτῷ ὁ ἀσθενῶν· Κύριε, ἄνθρωπον οὐκ ἔχω ἵνα ὅταν ταραχθῇ τὸ ὕδωρ βάλῃ με εἰς τὴν κολυμβήθραν· ἐν ᾧ δὲ ἔρχομαι ἐγὼ ἄλλος πρὸ ἐμοῦ καταβαίνει.
		% 8 λέγει αὐτῷ ὁ Ἰησοῦς· Ἔγειρε ἆρον τὸν κράβαττόν σου καὶ περιπάτει.
		% 9 καὶ εὐθέως ἐγένετο ὑγιὴς ὁ ἄνθρωπος καὶ ἦρε τὸν κράβαττον αὐτοῦ καὶ περιεπάτει.
		% }
		
		% \gr{
		% Ἦν δὲ σάββατον ἐν ἐκείνῃ τῇ ἡμέρᾳ.
		% 10 ἔλεγον οὖν οἱ Ἰουδαῖοι τῷ τεθεραπευμένῳ· Σάββατόν ἐστιν, καὶ οὐκ ἔξεστίν σοι ἆραι τὸν κράβαττον.
		% 11 ὃς δὲ ἀπεκρίθη αὐτοῖς· Ὁ ποιήσας με ὑγιῆ ἐκεῖνός μοι εἶπεν Ἆρον τὸν κράβαττόν σου καὶ περιπάτει.
		% 12 ἠρώτησαν οὖν αὐτόν· Τίς ἐστιν ὁ ἄνθρωπος ὁ εἰπών σοι· Ἆρον καὶ περιπάτει;
		% 13 ὁ δὲ ἰαθεὶς οὐκ ᾔδει τίς ἐστιν, ὁ γὰρ Ἰησοῦς ἐξένευσεν ὄχλου ὄντος ἐν τῷ τόπῳ.
		% 14 μετὰ ταῦτα εὑρίσκει αὐτὸν ὁ Ἰησοῦς ἐν τῷ ἱερῷ καὶ εἶπεν αὐτῷ· Ἴδε ὑγιὴς γέγονας· μηκέτι ἁμάρτανε, ἵνα μὴ χεῖρόν σοί τι γένηται.
		% 15 ἀπῆλθεν ὁ ἄνθρωπος καὶ ἀνήγγειλεν τοῖς Ἰουδαίοις ὅτι Ἰησοῦς ἐστιν ὁ ποιήσας αὐτὸν ὑγιῆ.
		% 16 καὶ διὰ τοῦτο ἐδίωκον οἱ Ἰουδαῖοι τὸν Ἰησοῦν ὅτι ταῦτα ἐποίει ἐν σαββάτῳ.
		% 17 ὁ δὲ ἀπεκρίνατο αὐτοῖς· Ὁ πατήρ μου ἕως ἄρτι ἐργάζεται κἀγὼ ἐργάζομαι.
		% 18 διὰ τοῦτο οὖν μᾶλλον ἐζήτουν αὐτὸν οἱ Ἰουδαῖοι ἀποκτεῖναι ὅτι οὐ μόνον ἔλυε τὸ σάββατον, ἀλλὰ καὶ πατέρα ἴδιον ἔλεγε τὸν θεόν, ἴσον ἑαυτὸν ποιῶν τῷ θεῷ.
		% }
		
% % -----------------------------------------------------------
% \section{95 | Authority of the Son}
% % -----------------------------------------------------------
	% \label{LCHRIST-095}
	
	% % ----------------------
	% \subsection{John 5:19--29}
	% % ----------------------
		% \label{LCHRIST-095-J}
		
		% \gr{
		% 19 Ἀπεκρίνατο οὖν ὁ Ἰησοῦς καὶ ἔλεγεν αὐτοῖς· Ἀμὴν ἀμὴν λέγω ὑμῖν, οὐ δύναται ὁ υἱὸς ποιεῖν ἀφ᾽ ἑαυτοῦ οὐδὲν ἐὰν μή τι βλέπῃ τὸν πατέρα ποιοῦντα· ἃ γὰρ ἂν ἐκεῖνος ποιῇ, ταῦτα καὶ ὁ υἱὸς ὁμοίως ποιεῖ.
		% 20 ὁ γὰρ πατὴρ φιλεῖ τὸν υἱὸν καὶ πάντα δείκνυσιν αὐτῷ ἃ αὐτὸς ποιεῖ, καὶ μείζονα τούτων δείξει αὐτῷ ἔργα, ἵνα ὑμεῖς θαυμάζητε.
		% 21 ὥσπερ γὰρ ὁ πατὴρ ἐγείρει τοὺς νεκροὺς καὶ ζῳοποιεῖ, οὕτως καὶ ὁ υἱὸς οὓς θέλει ζῳοποιεῖ.}%
			% \footnote{%
				% \label{LCHRIST-095-zoopoiei}%
				% See also \hyperref[1CORINTHIANS-15-22]{1 Corinthians 15:22}, which uses a future passive of \gr{ζωοποιέω}, and also \hyperref[2NEPHI-25-25]{2 Nephi 25:25}.
				% }
		% \gr{22 οὐδὲ γὰρ ὁ πατὴρ κρίνει οὐδένα, ἀλλὰ τὴν κρίσιν πᾶσαν δέδωκεν τῷ υἱῷ,
		% 23 ἵνα πάντες τιμῶσι τὸν υἱὸν καθὼς τιμῶσι τὸν πατέρα. ὁ μὴ τιμῶν τὸν υἱὸν οὐ τιμᾷ τὸν πατέρα τὸν πέμψαντα αὐτόν.
		% 24 Ἀμὴν ἀμὴν λέγω ὑμῖν ὅτι ὁ τὸν λόγον μου ἀκούων καὶ πιστεύων τῷ πέμψαντί με ἔχει ζωὴν αἰώνιον, καὶ εἰς κρίσιν οὐκ ἔρχεται ἀλλὰ μεταβέβηκεν ἐκ τοῦ θανάτου εἰς τὴν ζωήν.
		% }
		
		% \gr{
		% 25 Ἀμὴν ἀμὴν λέγω ὑμῖν ὅτι ἔρχεται ὥρα καὶ νῦν ἐστιν ὅτε οἱ νεκροὶ ἀκούσουσιν τῆς φωνῆς τοῦ υἱοῦ τοῦ θεοῦ καὶ οἱ ἀκούσαντες ζήσουσιν.
		% 26 ὥσπερ γὰρ ὁ πατὴρ ἔχει ζωὴν ἐν ἑαυτῷ, οὕτως καὶ τῷ υἱῷ ἔδωκεν ζωὴν ἔχειν ἐν ἑαυτῷ·
		% 27 καὶ ἐξουσίαν ἔδωκεν αὐτῷ κρίσιν ποιεῖν, ὅτι υἱὸς ἀνθρώπου ἐστίν.
		% 28 μὴ θαυμάζετε τοῦτο, ὅτι ἔρχεται ὥρα ἐν ᾗ πάντες οἱ ἐν τοῖς μνημείοις ἀκούσουσιν τῆς φωνῆς αὐτοῦ
		% 29 καὶ ἐκπορεύσονται οἱ τὰ ἀγαθὰ ποιήσαντες εἰς ἀνάστασιν ζωῆς, οἱ δὲ τὰ φαῦλα πράξαντες εἰς ἀνάστασιν κρίσεως.
		% }
		
% % -----------------------------------------------------------
% \section{96 | Come to Me and Rest}
% % -----------------------------------------------------------
	% \label{LCHRIST-096}
	
	% % ----------------------
	% \subsection{Matthew 11:25--30}
	% % ----------------------
		% \label{LCHRIST-096-MT}
		
		% \gr{
		% 25 Ἐν ἐκείνῳ τῷ καιρῷ ἀποκριθεὶς ὁ Ἰησοῦς εἶπεν· Ἐξομολογοῦμαί σοι, πάτερ κύριε τοῦ οὐρανοῦ καὶ τῆς γῆς, ὅτι ἔκρυψας ταῦτα ἀπὸ σοφῶν καὶ συνετῶν, καὶ ἀπεκάλυψας αὐτὰ νηπίοις·
		% 26 ναί, ὁ πατήρ, ὅτι οὕτως εὐδοκία ἐγένετο ἔμπροσθέν σου.
		% 27 Πάντα μοι παρεδόθη ὑπὸ τοῦ πατρός μου, καὶ οὐδεὶς ἐπιγινώσκει τὸν υἱὸν εἰ μὴ ὁ πατήρ, οὐδὲ τὸν πατέρα τις ἐπιγινώσκει εἰ μὴ ὁ υἱὸς καὶ ᾧ ἐὰν βούληται ὁ υἱὸς ἀποκαλύψαι.
		% }
		
		% \gr{
		% 28 Δεῦτε πρός με πάντες οἱ κοπιῶντες καὶ πεφορτισμένοι, κἀγὼ ἀναπαύσω ὑμᾶς.
		% 29 ἄρατε τὸν ζυγόν μου ἐφ᾽ ὑμᾶς καὶ μάθετε ἀπ᾽ ἐμοῦ, ὅτι πραΰς εἰμι καὶ ταπεινὸς τῇ καρδίᾳ, καὶ εὑρήσετε ἀνάπαυσιν ταῖς ψυχαῖς ὑμῶν·
		% 30 ὁ γὰρ ζυγός μου χρηστὸς καὶ τὸ φορτίον μου ἐλαφρόν ἐστιν.
		% }
		
% % -----------------------------------------------------------
% \section{97 | Witnesses to Jesus}
% % -----------------------------------------------------------
	% \label{LCHRIST-097}
	
	% % ----------------------
	% \subsection{John 5:30--47}
	% % ----------------------
		% \label{LCHRIST-097-J}
		
		% \gr{
		% 30 Οὐ δύναμαι ἐγὼ ποιεῖν ἀπ᾽ ἐμαυτοῦ οὐδέν· καθὼς ἀκούω κρίνω, καὶ ἡ κρίσις ἡ ἐμὴ δικαία ἐστίν, ὅτι οὐ ζητῶ τὸ θέλημα τὸ ἐμὸν ἀλλὰ τὸ θέλημα τοῦ πέμψαντός με.
		% }
		
		% \gr{
		% 31 Ἐὰν ἐγὼ μαρτυρῶ περὶ ἐμαυτοῦ, ἡ μαρτυρία μου οὐκ ἔστιν ἀληθής·
		% 32 ἄλλος ἐστὶν ὁ μαρτυρῶν περὶ ἐμοῦ, καὶ οἶδα ὅτι ἀληθής ἐστιν ἡ μαρτυρία ἣν μαρτυρεῖ περὶ ἐμοῦ.
		% 33 ὑμεῖς ἀπεστάλκατε πρὸς Ἰωάννην, καὶ μεμαρτύρηκε τῇ ἀληθείᾳ·
		% 34 ἐγὼ δὲ οὐ παρὰ ἀνθρώπου τὴν μαρτυρίαν λαμβάνω, ἀλλὰ ταῦτα λέγω ἵνα ὑμεῖς σωθῆτε.
		% 35 ἐκεῖνος ἦν ὁ λύχνος ὁ καιόμενος καὶ φαίνων, ὑμεῖς δὲ ἠθελήσατε ἀγαλλιαθῆναι πρὸς ὥραν ἐν τῷ φωτὶ αὐτοῦ·
		% 36 ἐγὼ δὲ ἔχω τὴν μαρτυρίαν μείζω τοῦ Ἰωάννου, τὰ γὰρ ἔργα ἃ δέδωκέν μοι ὁ πατὴρ ἵνα τελειώσω αὐτά, αὐτὰ τὰ ἔργα ἃ ποιῶ, μαρτυρεῖ περὶ ἐμοῦ ὅτι ὁ πατήρ με ἀπέσταλκεν,
		% 37 καὶ ὁ πέμψας με πατὴρ ἐκεῖνος μεμαρτύρηκεν περὶ ἐμοῦ. οὔτε φωνὴν αὐτοῦ πώποτε ἀκηκόατε οὔτε εἶδος αὐτοῦ ἑωράκατε,
		% 38 καὶ τὸν λόγον αὐτοῦ οὐκ ἔχετε ἐν ὑμῖν μένοντα, ὅτι ὃν ἀπέστειλεν ἐκεῖνος τούτῳ ὑμεῖς οὐ πιστεύετε.
		% }
		
		% \gr{
		% 39 Ἐραυνᾶτε τὰς γραφάς, ὅτι ὑμεῖς δοκεῖτε ἐν αὐταῖς ζωὴν αἰώνιον ἔχειν· καὶ ἐκεῖναί εἰσιν αἱ μαρτυροῦσαι περὶ ἐμοῦ·
		% 40 καὶ οὐ θέλετε ἐλθεῖν πρός με ἵνα ζωὴν ἔχητε.
		% 41 δόξαν παρὰ ἀνθρώπων οὐ λαμβάνω,
		% 42 ἀλλὰ ἔγνωκα ὑμᾶς ὅτι τὴν ἀγάπην τοῦ θεοῦ οὐκ ἔχετε ἐν ἑαυτοῖς.
		% 43 ἐγὼ ἐλήλυθα ἐν τῷ ὀνόματι τοῦ πατρός μου καὶ οὐ λαμβάνετέ με· ἐὰν ἄλλος ἔλθῃ ἐν τῷ ὀνόματι τῷ ἰδίῳ, ἐκεῖνον λήμψεσθε.
		% 44 πῶς δύνασθε ὑμεῖς πιστεῦσαι, δόξαν παρ᾽ ἀλλήλων λαμβάνοντες, καὶ τὴν δόξαν τὴν παρὰ τοῦ μόνου θεοῦ οὐ ζητεῖτε;
		% 45 μὴ δοκεῖτε ὅτι ἐγὼ κατηγορήσω ὑμῶν πρὸς τὸν πατέρα· ἔστιν ὁ κατηγορῶν ὑμῶν Μωϋσῆς, εἰς ὃν ὑμεῖς ἠλπίκατε.
		% 46 εἰ γὰρ ἐπιστεύετε Μωϋσεῖ, ἐπιστεύετε ἂν ἐμοί, περὶ γὰρ ἐμοῦ ἐκεῖνος ἔγραψεν.
		% 47 εἰ δὲ τοῖς ἐκείνου γράμμασιν οὐ πιστεύετε, πῶς τοῖς ἐμοῖς ῥήμασιν πιστεύσετε;
		% }
		
% % -----------------------------------------------------------
% \section{98 | Corn on the Sabbath}
% % -----------------------------------------------------------
	% \label{LCHRIST-098}
	
	% % ----------------------
	% \subsection{Matthew 12:1--8}
	% % ----------------------
		% \label{LCHRIST-098-MT}
		
		% \gr{
		% 1 Ἐν ἐκείνῳ τῷ καιρῷ ἐπορεύθη ὁ Ἰησοῦς τοῖς σάββασιν διὰ τῶν σπορίμων· οἱ δὲ μαθηταὶ αὐτοῦ ἐπείνασαν καὶ ἤρξαντο τίλλειν στάχυας καὶ ἐσθίειν.
		% 2 οἱ δὲ Φαρισαῖοι ἰδόντες εἶπαν αὐτῷ· Ἰδοὺ οἱ μαθηταί σου ποιοῦσιν ὃ οὐκ ἔξεστιν ποιεῖν ἐν σαββάτῳ.
		% 3 ὁ δὲ εἶπεν αὐτοῖς· Οὐκ ἀνέγνωτε τί ἐποίησεν Δαυὶδ ὅτε ἐπείνασεν καὶ οἱ μετ᾽ αὐτοῦ;
		% 4 πῶς εἰσῆλθεν εἰς τὸν οἶκον τοῦ θεοῦ καὶ τοὺς ἄρτους τῆς προθέσεως ἔφαγον, ὃ οὐκ ἐξὸν ἦν αὐτῷ φαγεῖν οὐδὲ τοῖς μετ᾽ αὐτοῦ, εἰ μὴ τοῖς ἱερεῦσιν μόνοις;
		% 5 ἢ οὐκ ἀνέγνωτε ἐν τῷ νόμῳ ὅτι τοῖς σάββασιν οἱ ἱερεῖς ἐν τῷ ἱερῷ τὸ σάββατον βεβηλοῦσιν καὶ ἀναίτιοί εἰσιν;
		% 6 λέγω δὲ ὑμῖν ὅτι τοῦ ἱεροῦ μεῖζόν ἐστιν ὧδε.
		% 7 εἰ δὲ ἐγνώκειτε τί ἐστιν· Ἔλεος θέλω καὶ οὐ θυσίαν, οὐκ ἂν κατεδικάσατε τοὺς ἀναιτίους.
		% 8 κύριος γάρ ἐστιν τοῦ σαββάτου ὁ υἱὸς τοῦ ἀνθρώπου.}%
			% \footnote{%
				% \label{LCHRIST-098-kurios}%
				% All three accounts of this event end with Christ saying a variation of \gr{κύριος γάρ ἐστιν τοῦ σαββάτου ὁ υἱὸς τοῦ ἀνθρώπου} (the same parts are in all three accounts). I am not sure what it means exactly, other than to say that a \gr{κύριος} is someone who has authority over something or who governs it, and so Christ has some sort of authority over the sabbath. Even if that's right, I'm not sure what the significance of it is beyond the principle that Christ teaches here, which I think appears most clearly in \hyperref[LCHRIST-098-MK]{Mark 2:27}.
				% }
		
	% % ----------------------
	% \subsection{Mark 2:23--28}
	% % ----------------------
		% \label{LCHRIST-098-MK}
		
		% \gr{
		% 23 Καὶ ἐγένετο αὐτὸν ἐν τοῖς σάββασιν παραπορεύεσθαι διὰ τῶν σπορίμων, καὶ οἱ μαθηταὶ αὐτοῦ ἤρξαντο ὁδὸν ποιεῖν τίλλοντες τοὺς στάχυας.
		% 24 καὶ οἱ Φαρισαῖοι ἔλεγον αὐτῷ· Ἴδε τί ποιοῦσιν τοῖς σάββασιν ὃ οὐκ ἔξεστιν;
		% 25 καὶ λέγει αὐτοῖς· Οὐδέποτε ἀνέγνωτε τί ἐποίησεν Δαυὶδ ὅτε χρείαν ἔσχεν καὶ ἐπείνασεν αὐτὸς καὶ οἱ μετ᾽ αὐτοῦ;
		% 26 πῶς εἰσῆλθεν εἰς τὸν οἶκον τοῦ θεοῦ ἐπὶ Ἀβιαθὰρ ἀρχιερέως καὶ τοὺς ἄρτους τῆς προθέσεως ἔφαγεν, οὓς οὐκ ἔξεστιν φαγεῖν εἰ μὴ τοὺς ἱερεῖς, καὶ ἔδωκεν καὶ τοῖς σὺν αὐτῷ οὖσιν;
		% 27 καὶ ἔλεγεν αὐτοῖς· Τὸ σάββατον διὰ τὸν ἄνθρωπον ἐγένετο καὶ οὐχ ὁ ἄνθρωπος διὰ τὸ σάββατον·}%
			% \footnote{%
				% \label{LCHRIST-098-anthopon}%
				% For me, this is the most important principle that comes out of this event. \gr{Τὸ σάββατον διὰ τὸν ἄνθρωπον ἐγένετο}, and not the reverse---in other words, people and their needs (like eating) are more fundamental than the sabbath and the rules that govern its observance. The sabbath was given or commanded to people for specific purposes, to give them some kind of benefit, but it does not supersede them and their other concerns. There is still room here for sabbath observance, but it's an important principle. I wonder also how widely a similar sort of principle applies to other commandments. What's also interesting to me is that thinking about this principle makes me want to think more deliberately about what I do on the sabbath; it doesn't make me want to just say, oh good, I can just fulfill all my ``needs'' on the sabbath.
				
				% Note also that the principles taught by Christ during \hyperref[LCHRIST-099]{this healing} are similar to the one taught here.
				% }
		% \gr{28 ὥστε κύριός ἐστιν ὁ υἱὸς τοῦ ἀνθρώπου καὶ τοῦ σαββάτου.
		% }
		
	% % ----------------------
	% \subsection{Luke 6:1--5}
	% % ----------------------
		% \label{LCHRIST-098-L}
		
		% \gr{
		% 1 Ἐγένετο δὲ ἐν σαββάτῳ διαπορεύεσθαι αὐτὸν διὰ σπορίμων, καὶ ἔτιλλον οἱ μαθηταὶ αὐτοῦ καὶ ἤσθιον τοὺς στάχυας ψώχοντες ταῖς χερσίν.
		% 2 τινὲς δὲ τῶν Φαρισαίων εἶπαν· Τί ποιεῖτε ὃ οὐκ ἔξεστιν τοῖς σάββασιν;
		% 3 καὶ ἀποκριθεὶς πρὸς αὐτοὺς εἶπεν ὁ Ἰησοῦς· Οὐδὲ τοῦτο ἀνέγνωτε ὃ ἐποίησεν Δαυὶδ ὁπότε ἐπείνασεν αὐτὸς καὶ οἱ μετ᾽ αὐτοῦ ὄντες;
		% 4 ὡς εἰσῆλθεν εἰς τὸν οἶκον τοῦ θεοῦ καὶ τοὺς ἄρτους τῆς προθέσεως λαβὼν ἔφαγεν καὶ ἔδωκεν τοῖς μετ᾽ αὐτοῦ, οὓς οὐκ ἔξεστιν φαγεῖν εἰ μὴ μόνους τοὺς ἱερεῖς;
		% 5 καὶ ἔλεγεν αὐτοῖς· Κύριός ἐστιν τοῦ σαββάτου ὁ υἱὸς τοῦ ἀνθρώπου.
		% }
		
% % -----------------------------------------------------------
% \section{99 | Healing a Withered Hand}
% % -----------------------------------------------------------
	% \label{LCHRIST-099}
	
	% % ----------------------
	% \subsection{Matthew 12:9--14}
	% % ----------------------
		% \label{LCHRIST-099-MT}
		
		% \footnote{%
			% \label{LCHRIST-099-same}%
			% See also the story of the \hyperref[LCHRIST-098]{corn} for more teachings related to this event. The \hyperref[LCHRIST-098-MK]{Mark} account is especially important, which is where we read that \gr{Τὸ σάββατον διὰ τὸν ἄνθρωπον ἐγένετο}.
			% }%
		% \gr{%
		% 9 Καὶ μεταβὰς ἐκεῖθεν ἦλθεν εἰς τὴν συναγωγὴν αὐτῶν·
		% 10 καὶ ἰδοὺ ἄνθρωπος χεῖρα ἔχων ξηράν. καὶ ἐπηρώτησαν αὐτὸν λέγοντες· Εἰ ἔξεστι τοῖς σάββασιν θεραπεύειν; ἵνα κατηγορήσωσιν αὐτοῦ.
		% 11 ὁ δὲ εἶπεν αὐτοῖς· Τίς ἔσται ἐξ ὑμῶν ἄνθρωπος ὃς ἕξει πρόβατον ἕν, καὶ ἐὰν ἐμπέσῃ τοῦτο τοῖς σάββασιν εἰς βόθυνον, οὐχὶ κρατήσει αὐτὸ καὶ ἐγερεῖ;
		% 12 πόσῳ οὖν διαφέρει ἄνθρωπος προβάτου. ὥστε ἔξεστιν τοῖς σάββασιν καλῶς ποιεῖν.
		% 13 τότε λέγει τῷ ἀνθρώπῳ· Ἔκτεινόν σου τὴν χεῖρα· καὶ ἐξέτεινεν, καὶ ἀπεκατεστάθη ὑγιὴς ὡς ἡ ἄλλη.
		% 14 ἐξελθόντες δὲ οἱ Φαρισαῖοι συμβούλιον ἔλαβον κατ᾽ αὐτοῦ ὅπως αὐτὸν ἀπολέσωσιν.
		% }
		
	% % ----------------------
	% \subsection{Mark 3:1--6}
	% % ----------------------
		% \label{LCHRIST-099-MK}
		
		% \gr{
		% 1 Καὶ εἰσῆλθεν πάλιν εἰς συναγωγήν, καὶ ἦν ἐκεῖ ἄνθρωπος ἐξηραμμένην ἔχων τὴν χεῖρα.
		% 2 καὶ παρετήρουν αὐτὸν εἰ τοῖς σάββασιν θεραπεύσει αὐτόν, ἵνα κατηγορήσωσιν αὐτοῦ.
		% 3 καὶ λέγει τῷ ἀνθρώπῳ τῷ τὴν χεῖρα ἔχοντι ξηράν· Ἔγειρε εἰς τὸ μέσον.
		% 4 καὶ λέγει αὐτοῖς· Ἔξεστιν τοῖς σάββασιν ἀγαθοποιῆσαι ἢ κακοποιῆσαι, ψυχὴν σῶσαι ἢ ἀποκτεῖναι; οἱ δὲ ἐσιώπων.
		% 5 καὶ περιβλεψάμενος αὐτοὺς μετ᾽ ὀργῆς, συλλυπούμενος ἐπὶ τῇ πωρώσει τῆς καρδίας αὐτῶν, λέγει τῷ ἀνθρώπῳ· Ἔκτεινον τὴν χεῖρα· καὶ ἐξέτεινεν, καὶ ἀπεκατεστάθη ἡ χεὶρ αὐτοῦ.
		% 6 καὶ ἐξελθόντες οἱ Φαρισαῖοι εὐθὺς μετὰ τῶν Ἡρῳδιανῶν συμβούλιον ἐδίδουν κατ᾽ αὐτοῦ ὅπως αὐτὸν ἀπολέσωσιν.
		% }
		
	% % ----------------------
	% \subsection{Luke 6:6--11}
	% % ----------------------
		% \label{LCHRIST-099-L}
		
		% \gr{
		% 6 Ἐγένετο δὲ ἐν ἑτέρῳ σαββάτῳ εἰσελθεῖν αὐτὸν εἰς τὴν συναγωγὴν καὶ διδάσκειν· καὶ ἦν ἄνθρωπος ἐκεῖ καὶ ἡ χεὶρ αὐτοῦ ἡ δεξιὰ ἦν ξηρά·
		% 7 παρετηροῦντο δὲ οἱ γραμματεῖς καὶ οἱ Φαρισαῖοι εἰ ἐν τῷ σαββάτῳ θεραπεύει, ἵνα εὕρωσιν κατηγορεῖν αὐτοῦ.
		% 8 αὐτὸς δὲ ᾔδει τοὺς διαλογισμοὺς αὐτῶν, εἶπεν δὲ τῷ ἀνδρὶ τῷ ξηρὰν ἔχοντι τὴν χεῖρα· Ἔγειρε καὶ στῆθι εἰς τὸ μέσον· καὶ ἀναστὰς ἔστη.
		% 9 εἶπεν δὲ ὁ Ἰησοῦς πρὸς αὐτούς· Ἐπερωτῶ ὑμᾶς, εἰ ἔξεστιν τῷ σαββάτῳ ἀγαθοποιῆσαι ἢ κακοποιῆσαι, ψυχὴν σῶσαι ἢ ἀπολέσαι;
		% 10 καὶ περιβλεψάμενος πάντας αὐτοὺς εἶπεν αὐτῷ· Ἔκτεινον τὴν χεῖρά σου· ὁ δὲ ἐποίησεν, καὶ ἀπεκατεστάθη ἡ χεὶρ αὐτοῦ.
		% 11 αὐτοὶ δὲ ἐπλήσθησαν ἀνοίας, καὶ διελάλουν πρὸς ἀλλήλους τί ἂν ποιήσαιεν τῷ Ἰησοῦ.
		% }
		
% % -----------------------------------------------------------
% \section{100 | The Chosen Healer-Servant}
% % -----------------------------------------------------------
	% \label{LCHRIST-100}
	
	% % ----------------------
	% \subsection{Matthew 12:15--21}
	% % ----------------------
		% \label{LCHRIST-100-MT}
		
		% \gr{
		% 15 Ὁ δὲ Ἰησοῦς γνοὺς ἀνεχώρησεν ἐκεῖθεν. καὶ ἠκολούθησαν αὐτῷ πολλοί, καὶ ἐθεράπευσεν αὐτοὺς πάντας,
		% 16 καὶ ἐπετίμησεν αὐτοῖς ἵνα μὴ φανερὸν αὐτὸν ποιήσωσιν,
		% 17 ἵνα πληρωθῇ τὸ ῥηθὲν διὰ Ἠσαΐου τοῦ προφήτου λέγοντος·
		% 18 Ἰδοὺ ὁ παῖς μου ὃν ᾑρέτισα, ὁ ἀγαπητός μου εἰς ὃν εὐδόκησεν ἡ ψυχή μου· θήσω τὸ πνεῦμά μου ἐπ᾽ αὐτόν, καὶ κρίσιν τοῖς ἔθνεσιν ἀπαγγελεῖ.
		% 19 οὐκ ἐρίσει οὐδὲ κραυγάσει, οὐδὲ ἀκούσει τις ἐν ταῖς πλατείαις τὴν φωνὴν αὐτοῦ.
		% 20 κάλαμον συντετριμμένον οὐ κατεάξει καὶ λίνον τυφόμενον οὐ σβέσει, ἕως ἂν ἐκβάλῃ εἰς νῖκος τὴν κρίσιν.
		% 21 καὶ τῷ ὀνόματι αὐτοῦ ἔθνη ἐλπιοῦσιν.
		% }
		
% -----------------------------------------------------------
\pagebreak{}
% -----------------------------------------------------------